\documentclass[11pt]{article}
\usepackage[a4paper,margin=18mm]{geometry}
\usepackage[no-math]{fontspec}
\setmainfont{Latin Modern Roman}
\usepackage{amsmath,amssymb,amsthm,xfrac,xspace,hyperref,graphicx,comment}
\excludecomment{DefTralics}
\providecommand{\dpl}{\displaystyle}
\providecommand{\backmatter}{}
\providecommand{\loccit}{\emph{loc. cit.}}

\providecommand{\bibliofont}{\small}
\newcommand{\ie}{i.e.,\xspace}
\newcommand{\eg}{e.g.,\xspace}
\newcommand{\etc}{etc.\xspace}
\newcommand{\cf}{cf.\xspace}
\newcommand{\resp}{resp.\xspace}
\newcommand{\smaller}{\small}
\newcommand{\Subsection}{\subsection}
\newcommand{\Subsubsection}{\subsubsection}
\newcommand{\skpt}{\smallskip}
\emergencystretch=3em

% Begin literal retained input author-preamble.tex
\usepackage[scr=rsfs,cal=euler]{mathalfa}

% Literal retained local macro yhmath


\hyphenation{After alge-braic Assu-me assump-tion assump-tions being deter-mines domains Dynam-i-cal-ly exists exist expand-ing instance invari-ant rena-me under}
\newenvironment{enumeratei}
{\bgroup\def\theenumi{\roman{enumi}}\def\theenumii{\arabic{enumii}}\begin{enumerate}}
{\end{enumerate}\egroup}
\newenvironment{enumeratea}
{\bgroup\def\theenumi{\alph{enumi}}\begin{enumerate}}
{\end{enumerate}\egroup}

\newcommand\mto{\mathchoice{\longmapsto}{\mapsto}{\mapsto}{\mapsto}}

\newcommand{\Psfrac}[2]{(\sfrac{#1}{#2})}
\newcommand{\psfrac}[2]{\sfrac{(#1)}{#2}}
\newcommand{\spfrac}[2]{\sfrac{#1}{(#2)}}
\newcommand{\pspfrac}[2]{\sfrac{(#1)}{(#2)}}

\newcommand{\llbracket}{[\![}
\newcommand{\rrbracket}{]\!]}

\let\epsilon\varepsilon
\let\sep :
\let\ts\textstyle\let\varstar\ast

\DeclareFontFamily{U}{mathc}{}
\DeclareFontShape{U}{mathc}{m}{it}%
{<->s*[1.03] mathc10}{}

\DeclareMathAlphabet{\mathscr}{U}{mathc}{m}{it}

\let\e\varepsilon
\let\F F

\newcommand{\set}[1]{\left\{#1\right\}}
\newcommand{\aset}[1]{\{#1\}}
\newcommand{\bset}[1]{\bigl\{#1\bigr\}}
\newcommand{\Bset}[1]{\Bigl\{#1\Bigr\}}
\newcommand{\norm}[1]{{\left\Vert#1\right\Vert}}
\newcommand{\anorm}[1]{{\Vert#1\Vert}}
\newcommand{\bnorm}[1]{{\bigl\Vert#1\bigr\Vert}}
\newcommand{\Bnorm}[1]{{\Bigl\Vert#1\Bigr\Vert}}
\newcommand{\abs}[1]{\left\vert#1\right\vert}
\newcommand{\aabs}[1]{\vert#1\vert}
\newcommand{\babs}[1]{\bigl\vert#1\bigr\vert}
\newcommand{\Babs}[1]{\Bigl\vert#1\Bigr\vert}
\newcommand{\bbabs}[1]{\biggl\vert#1\biggr\vert}
\newcommand{\pu}{{\mathbb{P}^1}}
\newcommand{\rest}[1]{ \arrowvert_{#1}}
\newcommand{\unsur}[1]{\frac{1}{#1}}
\newcommand{\lrpar}[1]{\left(#1\right)}
\newcommand{\alrpar}[1]{(#1)}
\newcommand{\blrpar}[1]{\bigl(#1\bigr)}
\newcommand{\Blrpar}[1]{\Bigl(#1\Bigr)}
\newcommand{\bblrpar}[1]{\biggl(#1\biggr)}
\newcommand{\bra}[1]{\left\langle #1\right\rangle}
\newcommand{\inv}{^{-1}}

\DeclareMathOperator{\aut}{Aut}
\DeclareMathOperator{\supp}{Supp}
\DeclareMathOperator{\mult}{mult}
\DeclareMathOperator{\Crit}{Crit}
\DeclareMathOperator{\id}{id}
\DeclareMathOperator{\dist}{dist}
\DeclareMathOperator{\Hal}{Hal}
\DeclareMathOperator{\Stab}{Stab}

\newcommand{\disk}{\mathbb{D}}
\newcommand{\ee}{\mathbb{E}}
\newcommand{\pp}{\mathbb{P}}\let\P\pp

\def\C{\mathbf{C}}
\def\R{\mathbf{R}}
\def\Q{\mathbf{Q}}
\def\Z{\mathbf{Z}}
\def\N{\mathbf{N}}
\def\bfk{{\mathbf{k}}}

\newcommand{\bfx}{\boldsymbol{x}}
\newcommand{\bfy}{\boldsymbol{y}}
\newcommand{\bfz}{\boldsymbol{z}}

\newcommand{\cX}{\mathcal{X}}
\newcommand{\cW}{\mathcal{W}}

\newcommand{\dil}{{\mathrm{Dil}}}
\newcommand{\Lat}{{\mathrm{Lat}}}
\newcommand\Lip{\mathrm{Lip}}
\newcommand{\NS}{{\mathrm{NS}}}
\newcommand{\Sing}{{\mathrm{Sing}}}
\newcommand{\Tang}{{\mathrm{Tang}}}
\newcommand{\STang}{{\mathrm{STang}}}
\newcommand{\NT}{{\mathrm{NT}}}

\newcommand\Aut{\mathsf{Aut}}
\newcommand\Diff{{\mathsf{Diff}}}
\newcommand\ii{{\mathsf{i}}}
\newcommand\jj{{\mathsf{j}}}
\newcommand{\Ima}{{\mathsf{Im}}}
\newcommand{\Rea}{{\mathsf{Re}}}
\newcommand\PGL{{\mathsf{PGL}}}
\newcommand\GL{{\mathsf{GL}}}
\newcommand\SL{{\mathsf{SL}}}
\newcommand{\vol}{{\mathsf{vol}}}

\newcommand{\m}{\mathscr{m}}

\theoremstyle{plain}
\newtheorem{thm}{Theorem}[section]
\newtheorem{cor}[thm]{Corollary}
\newtheorem{pro}[thm]{Proposition}
\newtheorem{lem}[thm]{Lemma}

\theoremstyle{definition}
\newtheorem{defi}[thm]{Definition}
\newtheorem{example}[thm]{Example}
\newtheorem{rem}[thm]{Remark}

% End literal retained input author-preamble.tex

\newfontfamily\jepScriptMfont{DejaVu Math TeX Gyre}
\renewcommand{\m}{\ensuremath{\mathord{\text{\upshape\jepScriptMfont 𝓂}}}}
\numberwithin{equation}{section}
\begin{document}
\begin{center}{\Large Stable item 2320}\end{center}
\paragraph{Source and attribution.}
Serge Cantat, Romain Dujardin, \emph{Hyperbolicity for large automorphism groups of projective surfaces}, Journal de l’École polytechnique — Mathématiques 12 (2025), publisher version, DOI 10.5802/jep.294. \href{https://jep.centre-mersenne.org/articles/10.5802/jep.294/}{Publisher article}. Authors' text: \href{https://creativecommons.org/licenses/by/4.0/}{CC BY 4.0}.
\paragraph{Editorial problem boundary.}
Prove precisely the original Theorem4.3 below: both printed non-escape conclusions and local uniformity of the eventual time for a uniformly expanding finite orbit under C2 random diffeomorphisms satisfying (M2,+). Retain the full logarithmic-distance drift construction under the stated moment condition, all finite-orbit and expansion definitions, and the exact prior Margulis theorem used. Later algebraic-surface classification and Wehler applications are unselected context.
\paragraph{Difficulty (editorial): H3.}
The prior Margulis occupation theorem does not provide the needed drift from tangent expansion alone under a logarithmic moment of order just above one. The authors build the logarithmic-distance function and combine a bounded-map second-order Taylor comparison with a separate large-norm moment-tail estimate, choosing the neighborhood and cutoff so both errors preserve the expansion drift. This nonlinear drift construction is the remaining nonroutine step beyond the precise prior occupation theorem; no stronger exponential-moment assumption is introduced.
\paragraph{Editorial transcription note.}
Original author language, mathematics and ordering are preserved. Publisher front matter is replaced by this independent article wrapper. Bibliography is recovered from matching cached publisher HTML, including its TeX formulas. No new proof or mathematical repair is supplied. Surrounding original results are retained for conservative context and are not extra selected corpus claims.
\clearpage

% Begin literal retained input author-body.tex
\section{Introduction}\label{sec:intro}

Let $X$ be a compact complex
surface and denote by $\Aut(X)$ its group of
automorphisms, \ie of holomorphic diffeomorphisms. Let $\Gamma$ be a subgroup of $\Aut(X)$. We~say that $\Gamma$ is \emph{non-elementary} if the subgroup
$\Gamma^\varstar\leq \GL(H^\varstar (X, \C))$ induced by the action of~$\Gamma$ on the De Rham cohomology of $X$ contains a non-abelian free group; the existence of
a non-elementary subgroup of $\Aut(X)$
implies that $X$ is projective (see \cite{Simons}). In~a series of articles \cite{stiffness, finite_orbits, invariant} we have
explored the dynamics of such a non-elementary group $\Gamma$ on $X$, notably by
means of random walk techniques. In~this paper, we~study the hyperbolicity properties of such random actions and their consequences.

\subsection{Wehler examples} \label{subs:wehler}
To understand the motivation behind our general results,
it is interesting to start with the Wehler family $ \cW$ of
surfaces of degree $(2,2,2)$ in $\P^1\times \P^1\times \P^1$, which has been a
recurring example in our work (see \eg \cite[\S 3]{stiffness}).
This family $ \cW$ depends on $26$ parameters and is naturally parameterized by
$\P^{26}(\C)$; we~shall denote by $ \cW_0\subset \cW$ the
Zariski open subset of smooth Wehler surfaces which do not contain any fiber of the three coordinate projections $\P^1\times \P^1\times \P^1\to \P^1$. Note that $\Aut(\P^1)^3$ acts on $\P^1\times\P^1\times \P^1$ as well as on $\cW$ and $\cW_0$.
For $X\in \cW_0$, the three natural projections $X\to \P^1\times \P^1$ are ramified covers of degree $2$; their deck transformations yield three holomorphic
involutions $\sigma_1$, $\sigma_2$, and $\sigma_3$;
the group $\Gamma$ generated by these involutions is non-elementary and isomorphic to
$\Z/2\Z\ast\Z/2\Z\ast\Z/2\Z$.

Since every $X\in \cW_0$ is a K3 surface,
there is a canonical $\Aut(X)$-invariant volume form $\vol_X$ on $X(\C)$; furthermore,
when $X$ is defined over $\R$ there is a canonical area form $\vol_{X(\R)}$
on $X(\R)$ which is invariant under the action of $\Aut(X_\R)$ (see Example~\ref{eg:K3_intro} below).
Slightly abusing notation, we~respectively denote by $\vol_X$ and $\vol_{X(\R)}$
the associated measures on $X$ and $X(\R)$, normalized to
have mass~$1$.

Our first main result is a complete description of orbit closures for
most parameters $X\in \cW_0$.
Recall that a 2-dimensional real submanifold
$Y\subset X$ is \emph{totally real} if for every $x\in Y$, $T_xY$ spans $T_xX$ as a
complex vector space.

\begin{thm}\label{thm:orbits_wehler}
There exists a dense and Zariski open subset
$ \cW_{\mathrm{exp}}\subset \cW_0(\C)$
such that for every $X\in \cW_{\mathrm{exp}}$, the action of $\Gamma=\langle \sigma_1, \sigma_2, \sigma_3\rangle$ on $X$ satisfies the following properties. There exists
a $\Gamma$-invariant finite set $F\subset X$ and
a $\Gamma$-invariant totally real analytic surface $Y\subset X$ (with possibly finitely many singular points)
such that for {\rm{every}} $x\in X$,
\begin{enumeratea}
\item\label{thm:orbits_wehlera} either $x\in F$ (and its orbit is finite);
\item or $\overline{\Gamma(x)}$ is a union of connected components of $Y$;
\item\label{thm:orbits_wehlerb} or $\overline{\Gamma(x)} = X$.
\end{enumeratea}
\end{thm}

As a general convention in the paper, by~``dense'' we mean dense for the Euclidean topology; when working with the Zariski topology, we~specify ``Zariski dense''. Note also that a Zariski open subset of a variety
$W$ is dense if and only if it intersects every component of $W$. In~this statement both $F$ and $Y$ may be empty, depending on~$X$.
For instance, \cite[Th.\,A]{finite_orbits} says that $F$ is empty for a very general Wehler surface $X\in \cW_0(\C)$,
\ie for $X$ in the complement of countably many proper Zariski closed subsets. A typical situation for
case~(b) is that $X$ is defined over $\R$ and $Y=X(\R)$.

\begin{rem}
To check that an orbit $\Gamma(x)$ is dense, it suffices to find a point $x'=(x_1',x_2',x_3')$ in $\Gamma(x)$ such that (1) the fiber $X_{x_3'}$ of the third projection containing $x'$ is a smooth curve and (2) $\sigma_1\circ\sigma_2$ acts on this genus $1$ curve $X_{x_3'}$ as a translation with dense orbits. Then, the closure of $\Gamma(x)$ is infinite and is not contained in a totally real surface, so case~(c) occurs.
Now, for $x_3$ outside a countable union of real analytic curves, $\sigma_1\circ\sigma_2$ has dense orbits along $X_{x_3}$. Thus, it is easy to produce examples of dense orbits.
On the other hand, given a specific Wehler surface $X$, it is a priori hard to decide whether there is a $\Gamma$ invariant real analytic surface $Y\subset X$.
\end{rem}

If we restrict to real parameters in $ \cW$, we
also have a fairly complete understanding
of the asymptotic distribution of random orbits. By~this we mean the following.
Let~$\nu$ be the probability measure on $\Gamma$ defined by
$\nu = \frac{1}{3}(\delta_{\sigma_1}+ \delta_{\sigma_2}+\delta_{\sigma_3})$. For any $x$ in $X(\R)$, and for any sequence $(g_i)$ of
automorphisms $g_i\in \Gamma$ chosen independently with distribution $\nu$, consider the trajectory $(g_n\cdots g_0(x))_{n\geq 0}$.
Let $X'(\R)$ be a union of connected components of $X(\R)$. We~say that these random trajectories, starting at $x$, are \emph{equidistributed} in $X'(\R)$ if for $\nu^\N$-almost every $(g_i)$, the empirical measures
$\unsur{n} \sum_{k=0}^{n-1} \delta_{g_k\cdots g_0(x)}$ converge
to the normalized volume form induced by $\vol_{X(\R)}$ on $X'(\R)$ as $n\to\infty$. The appearance of $X'(\R)$ is due to the fact that $\Gamma$ may not act transitively on the
components of $X(\R)$.

\begin{thm}\label{thm:equidistribution_wehler}
There exists a dense
and Zariski open subset $ \cW_{\mathrm{exp}}(\R)\subset \cW_0(\R)$
such that for {every} $X\in \cW_{\mathrm{exp}}(\R)$, there exists
a $\Gamma$-invariant finite set $F\subset X(\R)$ such that for {every} $x\in X(\R)$:
\begin{enumeratea}
\item\label{thm:equidistribution_wehlera} either $x\in F$;
\item\label{thm:equidistribution_wehlerb} or the random trajectories starting at $x$ are equidistributed in a union of connected components of $X(\R)$.
\end{enumeratea}
\end{thm}

An interesting point in Theorems~\ref{thm:orbits_wehler} and~\ref{thm:equidistribution_wehler} is that their
conclusions hold for \emph{every} $x\in X$. Let us explain how these theorems fall within
the progression of \cite{stiffness, finite_orbits, invariant} and what the last missing ingredient was until the present paper.

First, the existence of the maximal finite invariant set $F$ follows from \cite[Th.\,C]{finite_orbits}. One key point here is that $\Gamma$ contains \emph{parabolic elements}, that is automorphisms whose
action on $H^*(X;\C)$ is virtually unipotent and
of infinite order (see Section~\ref{sec:preliminaries}).

Now, the scheme of proof of Theorem~\ref{thm:equidistribution_wehler} is as follows.
The random walk on $\Gamma$ induced by $\nu$ gives rise to a random dynamical system on $X$. We~refer to~\cite{Kifer, br} for general references on this topic,
and to Sections~4 and~7 of \cite{stiffness} for our holomorphic context. In~particular, we~shall use the notions of stationary and invariant measures~$\mu$, of fibered entropy $h_\mu(X,\nu)$, etc.
Fix $x\in X\setminus F$. By~Breiman's ergodic theorem, for almost every sequence $(g_n)_{n\geq 0}$ with respect to the measure $\nu^\N$,
every cluster value of the sequence of
empirical measures
$\unsur{n} \sum_{k=0}^{n-1} \delta_{g_k\cdots g_0(x)}$ is a $\nu$-stationary measure (see~\cite[\S 2.2]{benoist-quint_book}). We~proved in~\cite{stiffness} that every $\nu$-stationary measure is $\Gamma$-invariant,\footnote{Here, we~use the fact that on
a typical Wehler surface there is no $\Gamma$-invariant curve, see~\cite[Lem.\,2.3]{finite_orbits}.}
and in~\cite{invariant} we showed that any invariant measure is either supported on $F$, or of the
form $\vol_{X'(\R)}$, for some union of components $X'(\R)$ of $X(\R)$. Therefore, any cluster value of
$\unsur{n} \sum_{k=0}^{n-1} \delta_{g_k\cdots g_0(x)}$ is a convex combination of point masses on $F$ and
$\vol_{X'(\R)}$. Thus the last step is to show that if $\Gamma(x)$ is infinite, the limiting empirical
measures give no mass to $F$.

For Theorem~\ref{thm:orbits_wehler} the situation is similar:
most of the work was done in \cite[\S 8]{invariant}, except that there we could not exclude
that the accumulation locus of an infinite orbit could be contained in a finite invariant set.
Note that since we are talking about orbit closures and not asymptotic distribution,
the full classification of stationary measures,
which is much harder and not yet complete in the
complex surface $X(\C)$, is not required here.

These difficulties were already addressed for homogeneous random dynamical systems in~\cite{benoist-quint3, Eskin-Lindenstrauss:RandomWalk} and in the context of non-linear actions on real
surfaces in~\cite{liu_thesis, chung}.
The key is to show that if $X$ belongs to the
dense Zariski open set $ \cW_{\mathrm{exp}}$ of Theorem~\ref{thm:orbits_wehler}
(\resp $ \cW_{\mathrm{exp}}(\R)$ of Theorem~\ref{thm:equidistribution_wehler}), the maximal finite invariant set
$F$ is \emph{repelling} for the random dynamics. Since we do not know the set $F$, nor its cardinality (examples of Wehler surfaces with large finite invariant sets were recently constructed
in~\cite{fuchs-litman-silverman-tran}), we~make a large detour and prove a uniform
hyperbolicity property for
the dynamics on \emph{the whole of} $X$, which is interesting in its own right: this is the \emph{uniform expansion} property that we present in detail in Section \ref{subs:UE_intro}.
Establishing this property relies on ergodic-theoretic arguments, the first of which is an
automatic hyperbolicity property that we describe in the next paragraph.

\subsection{Hyperbolicity of invariant measures} \label{subs:hyperbolicity}

It is a fundamental (and widely open) problem in conservative dynamics to show the typicality of non-zero Lyapunov exponents
on a set of positive Lebesgue measure. In~deterministic dynamics, a recent breakthrough is the work of
Berger and Turaev \cite{berger-turaev}. Adding some randomness makes such a hyperbolicity result easier to obtain:
see \cite{blumenthal-xue-young} for random perturbation of the standard map, and
\cite{barrientos-malicet, obata-poletti}
for random conservative diffeomorphisms on closed real surfaces. The results of Barrientos and Malicet \cite{barrientos-malicet} and of
Obata and Poletti \cite{obata-poletti} are perturbative in nature, so they do not give explicit examples. In~our context, the rigidity properties of holomorphic diffeomorphisms will
enable us to exhibit explicit criteria ensuring such a non-uniform hyperbolicity.

In~\cite{invariant} we have classified
invariant measures for non-elementary groups containing parabolic elements. We~say
that a measure $\mu$ on $X$ is
{\emph{Zariski diffuse}}
if it gives zero mass to proper Zariski closed subsets. If $\mu$ is $\Gamma$-invariant and ergodic for some
$\Gamma\subset \Aut(X)$, this is equivalent to its support $\supp(\mu)$
being Zariski dense. Roughly speaking, our classification of invariant measures says that every Zariski diffuse, ergodic, invariant probability measure is
given by an analytic $4$-form on $X$ or by an analytic $2$-form on some invariant, real analytic subset $Y\subset X$ of dimension $2$.
Here we proceed to a finer study of the dynamical properties of these invariant measures. For this, we
fix a probability measure $\nu$ on $\Aut(X)$ satisfying the moment condition
\begin{equation}\label{eq:moment}
\tag{M}
\int \lrpar{\log\norm{f}_{C^1(X)}+ \log\anorm{f\inv}_{C^1(X)}} \, d\nu (f) < + \infty;
\end{equation}
(see Sections~\ref{subs:moment} and~\ref{par:finite-orbits-C2-actions} and Remark~\ref{rem:moment-condition-holo} for discussions of stronger moment conditions),
and we view any invariant measure $\mu$
as a $\nu$-stationary measure, that is,
$\int f_\varstar\mu \, d\nu(f) = \mu$. Then by~\eqref{eq:moment},
the Lyapunov exponents of $\mu$ are well defined: for $\nu^\N$-almost every sequence $(g_i)$,
and $\mu$-almost every $x\in X$, $\frac{1}{n}\log \norm{D_x(g_{n-1}\circ \cdots \circ g_0)}$ converges towards the upper Lyapunov exponent $\lambda^+(x)\in \R$; by ergodicity of $\mu$, $\lambda^+(x)$ is almost surely equal to some constant $\lambda^+(\mu)$. Similarly, one defines the lower Lyapunov exponent $\lambda^-(\mu)$, and $\mu$ is said to be hyperbolic if $\lambda^+(\mu)>0>\lambda^-(\mu)$.

We denote by $\Gamma_\nu \subset \Aut(X)$
the closed subgroup generated by $\supp(\nu)$.\footnote{\label{page:discrete}Note that $\Aut(X)$ is discrete unless
$X$ is a torus, see \cite[\S 3]{stiffness}
so in most cases $\Gamma_\nu = \langle \supp(\nu) \rangle$.}

\begin{thm}\label{thm:hyperbolic}
Let $X$ be a compact complex surface and
$\Gamma$ be a non-elementary subgroup of $\Aut(X)$ containing parabolic elements.
Let $\mu$ be a Zariski diffuse ergodic $\Gamma$-invariant probability measure on $X$.
Let $\nu$ be any probability measure on $\Aut(X)$
satisfying $\Gamma_\nu=\Gamma$ and the moment condition~\eqref{eq:moment}.

Then, viewed as a $\nu$-stationary measure, $\mu$ is hyperbolic and its fiber entropy $h_\mu(X,\nu)$ is positive.
\end{thm}

A variant of this result will also be obtained when $\Gamma_\nu$ contains a Kummer example
instead of a parabolic element (see Theorem~\ref{thm:hyperbolic_kummer}).

\begin{example}\label{eg:K3_intro} When $X$ is a torus or a K3 surface, the canonical bundle $K_X$ is trivial and, up to multiplication by a complex number of modulus $1$, there is
a unique section $\Omega_X$ of $K_X$ that satisfies $\int_X\Omega_X\wedge \overline{\Omega_X}=1$. The volume form $\vol_X:=\Omega_X\wedge \overline{\Omega_X}$
is $\Aut(X)$-invariant. Likewise, every Enriques surface $S$ inherits such an invariant volume form $\vol_S$ from its universal cover $X$ (a $2$-to-$1$ cover by a K3 surface).
Under the assumptions of Theorem~\ref{thm:hyperbolic}, $\vol_X$ is $\Gamma$-ergodic,
thus we conclude that {\textit{it is hyperbolic}}. Other examples are provided by some rational surfaces (see the discussion on Coble surfaces in~\cite{Simons}).

In these situations the 2-form $\Omega_X$ also induces a natural measure $\vol_Y$ on any totally real surface $Y\subset X$ (see~\cite[Rem.\,2.3]{invariant}). For instance, if $X$ is projective and defined over $\R$, $\Gamma$ is contained in $\Aut(X_\R)$, and $Y$ is a $\Gamma$-invariant connected component of $X(\R)$, Theorem~\ref{thm:hyperbolic} asserts that $\vol_Y$ is hyperbolic.
\end{example}

\subsection{Uniform expansion}\label{subs:UE_intro}

Fix a Riemannian
metric on $X$. We~say that the measure~$\nu$ on $\Aut(X)$
is \emph{uniformly expanding} if there exists $c>0$ and an integer $n_0$ such that for \emph{every} $x\in X$ and \emph{every} $v\in T_xX\setminus\{ 0\}$,\vspace*{-3pt}\enlargethispage{.5\baselineskip}
\begin{equation}\label{eq:UE}
\int_{\Aut(X)} \log\Psfrac{\norm{D_xf (v)}}{\norm{v}} \,d\nu^{(n_0)}(f) \geq c;
\end{equation}
here $\nu^{(n)}$ denotes the $n^\mathrm{th}$ convolution power of $\nu$.
This notion is taken from
\cite{chung, dolgopyat-krikorian, liu_thesis, zzhang} (see also \cite{Eskin-Lindenstrauss:RandomWalk, Prohaska-Sert:TAMS} for the linear context) and has a number
of strong ergodic and topological consequences on the action of $\Gamma_\nu$.
So far, uniform expansion has been verified only in the context of homogeneous dynamics, or for certain perturbative situations, or with the help of numerical methods.
The geometric analysis of stationary measures developed in~\cite{stiffness} together with Theorem~\ref{thm:hyperbolic}
will be used to obtain the following result.

\begin{thm}\label{thm:criterion_uniform_expansion_parabolic}
Let $X$ be a compact complex surface which is not rational.
Let $\nu$ be a probability measure on $\Aut(X)$. Assume that:
{\em{(i)}} $\nu$ satisfies the moment condition~\eqref{eq:moment} and {\em{(ii)}} the group
$\Gamma = \Gamma_\nu$ is non-elementary and contains parabolic elements.

Then $\nu$ is uniformly expanding if and only if the following two conditions hold:
\begin{enumerate}
\item\label{thm:criterion_uniform_expansion_parabolic1} every finite $\Gamma$-orbit is uniformly expanding;
\item\label{thm:criterion_uniform_expansion_parabolic2} there is no $\Gamma$-invariant algebraic curve.
\end{enumerate}
\end{thm}

Here, by~definition,
a finite orbit $F$ of $\Gamma$ is said to be \emph{uniformly expanding} if Condition~\eqref{eq:UE}
holds for every $x\in F$. This is the repulsion property alluded to at the end of Section \ref{subs:wehler}.

Checking Condition~(2)
of Theorem~\ref{thm:criterion_uniform_expansion_parabolic} is not hard in practice and boils down to cohomological computations
(see Section \ref{subs:invariant_curves}). Therefore, most of the complexity in applying this theorem to practical situations
comes from the analysis of finite orbits. The simplest instance is when there are no finite orbits at all:

\begin{cor}\label{cor:no_algebraic_invariant_subset}
Under the assumptions of Theorem~\ref{thm:criterion_uniform_expansion_parabolic}, if there is no proper algebraic $\Gamma_\nu$-invariant subset, then
$\nu$ is uniformly expanding.
\end{cor}

By~\cite[Th.\,A]{finite_orbits}
the automorphism group of a very general Wehler surface has no proper Zariski closed invariant set. Since uniform expansion is an open property in the $C^1$ topology,
it holds on an open and dense set -- in the Euclidean topology -- of Wehler examples. In~the next few paragraphs we explain why it actually holds on a dense Zariski open set,
which is the main point in Theorems~\ref{thm:orbits_wehler} and~\ref{thm:equidistribution_wehler}.

First, for a given finite $\Gamma$-orbit $F$,
if $\nu$ is symmetric and satisfies a slightly stronger moment condition~\eqref{eq:moment+}, Theorem~\ref{thm:finite_UE} provides a
checkable necessary and sufficient condition for $F$ to be uniformly expanding:
it is equivalent to the tangent action of $\Gamma$ being proximal and strongly irreducible. It
follows that when $\nu$ is symmetric and $X$ is not rational
{\emph{the uniform expansion property
depends only on $\Gamma$, and not on $\nu$}} (Corollary~\ref{cor:independence_nu}).
Anticipating on these results, for a non-elementary subgroup $\Gamma\subset\Aut(X)$ on a non-rational surface
$X$, we~can say that \emph{the action is uniformly expanding} if this property holds for some (hence any)
symmetric probability measure~$\nu$ satisfying~\eqref{eq:moment+} and
generating~$\Gamma$.

In Section \ref{subs:finitary} we show that uniform expansion can be checked algorithmically.
The starting point is the fact that if $X$ is not a torus and Condition~(2) of
Theorem~\ref{thm:criterion_uniform_expansion_parabolic} holds,
then by~\cite[Th.\,C]{finite_orbits}, there are only finitely many finite orbits.
The difficulty is that there is no a priori bound on their number so far, even in the Wehler family.
Fortunately, we~prove that the number of non-expanding finite orbits can be controlled:

\begin{pro}
Let $X$ be a smooth projective surface and $\Gamma $
be a non-elementary subgroup of $\Aut(X)$
containing parabolic elements and without invariant algebraic curve. Then there is a computable number $N(X, \Gamma)$ such that any finite orbit of length greater than $N(X, \Gamma)$ is uniformly expanding.
Moreover, in the Wehler family, the number $N(X, \bra{\sigma_1, \sigma_2, \sigma_3})$ is uniformly bounded.
\end{pro}

See Theorem~\ref{thm:effective} for details on what we mean by computable. To conclude
from Condition~(2) that uniform
expansion holds on a Zariski open subset of $\mathcal W$, we~use the fact from~\cite{finite_orbits}
that on a Zariski open subset $\mathcal W_N\subset \mathcal W$, all finite orbits have length
greater than $N$ (see Theorem~\ref{thm:wehler_zariski} below). We~do not know the value of $N$ for the Wehler family but we do not expect it to be large.\footnote{We show that $N$ depends on the cardinality of a set of special ``non-twisting'' fibers of elliptic fibrations with automorphisms, which seems to be quite scarce, see~\cite[Rem.\,7.7.14]{duistermaat}.}
In particular the equations defining $ \cW_{\mathrm{exp}}$ and $ \cW_{\mathrm{exp}}(\R)$
could in principle be written down explicitly.

\subsection{Ergodicity} \label{subs:intro_ergodicity}
Given an action of a general non-elementary group $\Gamma$ on a compact complex surface $X$,
one may ask the following two basic questions:
does there exist a dense orbit? Is the action ergodic with respect to Lebesgue measure?
(The latter makes sense even when there is no invariant volume form.)
If $\Gamma$ contains a parabolic element, by~\cite{invariant} the answer to both questions is `yes',
but without parabolic elements, the answer is unknown.
A natural obstruction to the existence of
a dense orbit could be the presence of a non-trivial Fatou
component for $\Gamma$. No example of such a Fatou component is known so far;
note that examples do exist for algebraic actions on \emph{affine} surfaces
(see~\cite[\S 4.1]{Cantat:BHPS} or~\cite[Th.\,E]{rebelo-roeder}).

As a matter of fact, the failure of ergodicity is
associated to a lack of expansion: indeed a
theorem of Dolgopyat and Krikorian~\cite[\S 10]{dolgopyat-krikorian}
asserts that a conservative uniformly expanding action on a (real) surface must be ergodic. It is not difficult
to extend their argument to the complex setting (see Theorem~\ref{thm:dolgopyat-krikorian}). In~Theorem~\ref{thm:criterion_uniform_expansion} we state
a general criterion (\ie without parabolic elements) for uniform expansion
which shows that under the conditions \eqref{thm:criterion_uniform_expansion_parabolic1} and \eqref{thm:criterion_uniform_expansion_parabolic2} of
Theorem~\ref{thm:criterion_uniform_expansion_parabolic},
the failure of uniform expansion is due to the existence of a
$\Gamma$-invariant measure with exceptional properties (see Theorems~\ref{thm:criterion_uniform_expansion} and~\ref{thm:entropy_every_f}). We~expect it to be an extremely rare phenomenon.
Incidentally,
this shows that the question of ergodicity for general non-elementary groups (\ie without parabolic elements)
ultimately boils down to the
classification of $\Gamma$-invariant measures.

Another consequence of our results, together with~\cite{dolgopyat-krikorian}, is that
a generic real Wehler example is \emph{stably ergodic} among $C^2$ volume preserving actions, that is,
if $X$ belongs to the open set $ \cW_{\mathrm{exp}}(\R)$
of Theorem~\ref{thm:equidistribution_wehler} and
$\sigma'_1$, $\sigma'_2$, $\sigma'_3$ are $C^2$
volume preserving diffeomorphisms sufficiently close to
of $\sigma_1$, $\sigma_2$, $\sigma_3$ in the $C^1$ topology, then
$\Gamma' := \langle \sigma'_1,\sigma'_2,\sigma'_3\rangle $ is ergodic for $\vol_{X(\R)}$.

In an opposite direction, the examples from \cite[\S 9]{invariant} of $\Aut(X_\R)$-invariant
domains with boundary in $X(\R)$ (which admit an invariant curve)
provide explicit counterexamples to uniform expansion.
For another example, start with the round sphere $x^2+y^2+z^2=3$ in the affine space, and view it as a singular
Wehler surface. The three involutions act by changing the signs of the coordinates and the points $(\epsilon_1,\epsilon_2,\epsilon_3)$
with $\epsilon_i=\pm 1$ form an orbit of size $8$. Now, choose a smooth Wehler surface $X$ containing these $8$ points and tangent to the sphere at each of them; this imposes $16$ linear conditions on the coefficients of the equation defining $X$, thus such examples exist.
For such a surface, the $8$ points form a finite, non-expanding orbit of $\langle\sigma_1, \sigma_2, \sigma_3\rangle$ (the action of the stabilizer of $(1,1,1)$ on the tangent space is identical to that of the round sphere so it factorizes through a finite group).

\subsection{Organization of the paper}
The first part of this paper (Sections~\ref{sec:uniform_generalities} to~\ref{sec:criterion}) is
devoted to a general study of the notion of uniform expansion on compact (real) manifolds.
Much of this material is inspired from other sources; the novelty here is that we strive for
optimal moment conditions. We~see several reasons for this. First, it is an important trend
in random dynamics to look for optimal conditions in the measure rigidity results (an explicit motivation of~\cite{Eskin-Lindenstrauss:RandomWalk} is to extend the results of Benoist and Quint to measures with finite first moment). Next,
when a random dynamical system is generated by a probability measure $\nu$ with finite support $S=\supp(\nu)$ and one considers a finite index subgroup $\Gamma_0$ of $\Gamma:=\langle S\rangle$, then the support of the measure induced by $\nu$ on $\Gamma_0$ is infinite (albeit with exponential moments). Also, when looking for random dynamical systems with atypical features, the first examples are usually given by probability measures with only weak moment conditions.
Finally, we~expect fine moment estimates to be important in the process of trying to
improve Theorem~\ref{thm:margulis_totalement_reel} to include the case of some
singular subvarieties $Y\subset X$ (possibly in the spirit of~\hbox{\cite[Th.\,B']{benard-desaxce}}).

In Section~\ref{sec:uniform_generalities} we give several equivalent definitions of uniform expansion:
this is inspired by Liu~\cite{liu_thesis} and Chung~\cite{chung}. In~Section~\ref{sec:inducing} we show that uniform expansion is preserved
when restricting to a finite index
subgroup or taking a finite extension (Proposition~\ref{pro:induced_UE2}); this
is useful when dealing with invariant sets made of finitely many connected components. Section~\ref{sec:margulis} deals with the construction of
\emph{Margulis functions}. In~a nutshell, a Margulis function
near a finite uniformly expanding invariant set $F$ is a function
$u:M\setminus F\to \R_+$ that tends to infinity at $F$ and decreases on average along orbits.
The existence of such a function
guarantees that empirical measures of random orbits do not accumulate at $F$. These functions have played an
important role in random dynamics since the work of Eskin and Margulis~\cite{eskin-margulis}.
Here, thanks to the work by Bénard and De Saxcé~\cite{benard-desaxce}, we~construct such Margulis functions under optimal moment conditions (Theorem~\ref{thm:margulis});
note that the usual average
decay property $\int u (f(x)) d\nu(f) \leq a u(x)+b$, $a<1$, is then replaced by $\int u (f(x)) d\nu(f) \leq u(x) -\nobreak \gamma$, $\gamma>0$. This repulsion property does not hold if $F$ is an invariant submanifold (see Example~\ref{eg:margulis_submanifold}). However in the holomorphic context, Margulis functions
can be constructed for invariant
totally real manifolds of maximal dimension (Theorem~\ref{thm:margulis_totalement_reel}): a~typical situation
is that of $X(\R)\subset X$ for real projective manifolds. In~Section~\ref{sec:criterion}, we elaborate on an ergodic-theoretic criterion for uniform expansion borrowed
from~\cite{chung}.

In the second part of the paper (Sections~\ref{sec:preliminaries} to~\ref{sec:applications}), we
consider groups of automorphisms of projective surfaces. Theorem~\ref{thm:hyperbolic}
is established in Section~\ref{sec:hyperbolicity}. In~Section~\ref{sec:characterization}, we~prove a general version of
Theorem~\ref{thm:criterion_uniform_expansion_parabolic} and
study uniform expansion along periodic orbits; this makes essential use of the results of the first part.
The focus in Section \ref{subs:finitary} is on finding
algorithmically checkable conditions for uniform expansion along finite orbits (Theorem~\ref{thm:effective});
this leads to a precise description of the locus of uniform expansion in the Wehler family (Theorem~\ref{thm:wehler_zariski}). In~Section \ref{subs:thin}, we~construct uniformly expanding actions by perturbing Kummer examples in the Wehler family;
in particular this work for ``thin'' subgroups of $\Aut(X)$ containing no parabolic element. In~Section~\ref{sec:applications} we study orbit closures and equidistribution by proving
general versions of Theorems~\ref{thm:orbits_wehler} and~\ref{thm:equidistribution_wehler};
we also explain the adaptation to the complex setting of the ergodicity
theorem of Dolgopyat and Krikorian~\cite{dolgopyat-krikorian}.

The paper ends with an appendix on the rigidity of zero entropy measures.\enlargethispage{.5\baselineskip}

\subsection{Notes and comments} Theorem~\ref{thm:hyperbolic} was included in the first preprint version
of~\cite{stiffness}. We~were informed of ongoing projects by Aaron Brown, Alex Eskin,
Simion Filip and Federico Rodriguez Hertz, as well as Megan Roda,
on the classification of stationary measures for uniformly expanding actions.
This should fit nicely with our work; indeed, parts of this article are written so as to to be easily combined with such a classification (see \eg Theorem~\ref{thm:equidistribution_complex}),

\subsubsection*{Acknowledgements}
We are grateful to Jean-François Quint for useful comments on Margulis functions and to
the anonymous referees for their detailed reports and constructive suggestions.

\part{Uniform expansion for discrete group actions on manifolds}

\section{Generalities}\label{sec:uniform_generalities}

In this section, $M$ denotes a compact manifold. The group $\Diff^1(M)$
of $C^1$ diffeomorphisms of $M$ endowed with the $C^1$ topology is a polish space (see~\cite{rosendal:book}) and we consider a Borel
probability measure $\nu$ on it. We~fix a Riemannian metric on $M$. We~denote by $\norm{\cdot}$
the norm induced by the metric on the tangent bundle $TM$, and by $T^1M$ the unit tangent bundle.

\subsection{Moment conditions}\label{subs:moment}
If $f$ is a $C^1$-diffeomorphism of $M$, we~denote by $f_\varstar $ its action on $TM$. Note that
if $v\in TM$ is a tangent vector based at $x$ (that is, $v\in T_xM$), then
$f_\varstar v = D_xf(v)$ is based at $f(x)$. By~definition, $\norm{f}_{C^1(X)}$ is the supremum of
$v\mto \norm{f_\varstar v}$ on $T^1M$.
For $f\in \Diff^1(M)$ we put
\begin{equation}\label{eq:subadditive_function_L}
\mathrm L(f) = \log\norm{f}_{C^1(X)}+ \log\anorm{f\inv}_{C^1(X)};
\end{equation}
this quantity is subadditive: $\mathrm L(f\circ g)\leq \mathrm L(f)+\mathrm L(g)$.
For $p\geq 1$ we consider the moment conditions
\begin{align}\label{eq:momentp}
\tag{M$_p$}
\int \mathrm L(f)^p \, d\nu (f) < + \infty, \quad \quad \quad \quad\\
\label{eq:moment+}
\tag{M$_+$}
\exists p>1, \ \mathrm{(M}_p\mathrm{)} \text{ holds}, \quad \quad \quad\quad \\
\label{eq:momentexp}
\tag{M$_{\exp}$}
\exists t>0, \ \int \norm{f}^t_{C^1(X)}+ \anorm{f\inv}^t_{C^1(X)} \, d\nu (f) < + \infty.
\end{align}
When $p=1$, \eqref{eq:momentp} coincides with the moment condition \eqref{eq:moment} from the introduction. For $p>1$, \eqref{eq:momentp}
implies \eqref{eq:moment+} which implies \eqref{eq:moment}.
The subadditivity of $L$ and the convexity inequality
$(r\inv \sum_{i=1}^r L_i)^p\leq r\inv \sum_{i=1}^r L_i^p$ imply
\begin{equation}\label{eq:Mp_iteree}
\int \mathrm L(f)^p \, d\nu^{(r)} (f) \leq r^{p} \int \mathrm L(f)^p \, d\nu (f)
\end{equation}
for $p\in [1,+\infty[$ and $r\in \N^*$, where $\nu^{(r)}$ denotes the $r^\mathrm{th}$ convolution power of $\nu$.

\subsection{Notation for random compositions} \label{subs:notation}

Set $\Omega=\Diff^1(M)^\N$; its elements
are sequences $\omega = (f_n)_{n\geq 0}$ of diffeomorphisms. We~use the
probabilistic notation $\ee(\cdot)$ and $\pp(\cdot)$ for the expectation and probability with respect to $\nu^\N$
on the probability space $\Omega$. We~let
$(\mathcal F_n)_{n\geq 1}$ be the increasing sequence of $\sigma$-algebras in $\Omega$
generated by cylinders of length $n$, so that an event is $\mathcal F_n$-measurable if it depends only on
the first~$n$ terms $f_0, \ldots , f_{n-1}$ of $\omega=(f_n)_{n\geq 0}$.
For $\omega = (f_n)_{n\geq 0} \in \Omega$ we put $f^0_\omega=\id$ and
\begin{equation*}
f_\omega^n = f_{n-1}\circ \cdots \circ f_0
\end{equation*}
for $n\geq 1$; in particular $f_\omega^1 = f_0$. For $x$ in $M$ and $ v \in T_xM\setminus\{ 0\}$ we set
\begin{equation*}
x_{\omega, n} = f_\omega^n(x) \quad {\text{and}} \quad
v_{\omega, n} = \frac{ (f_\omega^n)_\varstar (v)} {\norm{ (f_\omega^n)_\varstar (v)}}\in T^1_{x_{\omega, n}}M.
\end{equation*}
For any sequence of integers
$0= k_0< k_1< \cdots < k_p = n$ the chain rule gives
\begin{equation}\label{eq:chain_rule}
\log\norm{(f_\omega^n)_\varstar v} = \sum_{j=0}^{p-1} \log \bnorm{ \blrpar{f^{k_{j+1}- k_j}_{\sigma^{k_j}\omega}}_\varstar { v_{\omega, k_j}}} .
\end{equation}
If $x$ is a point of $M$, we~denote by $\delta_x$ the Dirac mass at $x$. If $\mu$ is a probability measure on $M$, $\nu\star\mu$
is the measure defined by $(\nu\star\mu)(B)=\int (f_*\mu)(B) \; d\nu(f)$ for every Borel subset $B$ of $M$.

\subsection{Equivalent conditions for uniform expansion}

Recall that the probability measure $\nu$ on $\Diff^1(M)$ is \emph{uniformly expanding}
if there exists a real number $c>0$ and an integer $n_0\geq 1$ such that
\begin{equation}\label{eq:UEbisrepetita}
\text{ for every } v\in T^1X, \quad
\int \log {\norm{ f_\varstar ( v)}} \; d\nu^{(n_0)} (f) \geq c.
\end{equation}
Then, the cocycle relation for
$ \log \Psfrac{\norm{ f_\varstar ( v)}}{\norm{ v}}$ implies that
\begin{equation}\label{eq:UE_iteree}
\int \log {\norm{ f_\varstar ( v)}} \; d\nu^{(kn_0)}(f) \geq kc
\end{equation}
for every $k\geq 1$. Thus, $\nu$ is uniformly expanding if and only if
$\nu^{(n)}$ is uniformly expanding for some (and hence for all) $n$. It follows that the uniform expansion property
does not depend on our choice of a Riemannian metric on $M$.

\begin{rem} \label{rem:restriction_UE}
If $\nu$ is uniformly expanding and the submanifold $N\subset M$ is invariant under every diffeomorphism in the support of $\nu$,
then $\nu$ induces a uniformly expanding measure on $\Diff^1(N)$.
\end{rem}

\begin{lem}\label{lem:weak_UE}
Let $\nu$ be a probability measure on $\Gamma$ satisfying \eqref{eq:moment}.
It is uniformly expanding if and only if
\begin{equation}\label{eq:weak_UE}
\forall v\in T^1M, \ \exists n = n(v) \ \text{ such that } \
\int \log \norm{f_\varstar v } \; d\nu^{(n)}(f) >0.
\end{equation}
\end{lem}

This is Lemma 4.3.1 of \cite{liu_thesis}, but Liu assumes that the support of $\nu$ is compact; thus
we briefly reproduce his proof, assuming only \eqref{eq:moment}.

\begin{proof} We have to show that~\eqref{eq:weak_UE} implies~\eqref{eq:UEbisrepetita}.
Since $\abs{\log \norm{f_{\varstar}v}} \leq \mathrm L(f)$ for every $v\in T^1X$, the
dominated convergence theorem implies that, for every $n$,
\begin{equation*}
v\mto \int \log \norm {f_{\varstar}( v)} \; d\nu^{(n)}(f)
\end{equation*}
is continuous.
Thus by compactness, there exists a finite open cover $V_1, \ldots, V_p$ of $T^1M$, positive
real numbers $c_i$, and integers
$n_i$ such that
\begin{equation*}
\int \log \norm{f_\varstar ( v) } \; d\nu^{(n_i)}(f) \geq c_i
\end{equation*}
for every
$ v\!\in\! V_i$.
Set $c_0 \!=\! \min(c_i)$ and $n_0 \!=\! \max(n_i)$. For $ v\in T^1X$ and
$\omega\in \Omega$, define the stopping time $\tau_1( v, \omega)$ to
be the first integer $n\!\geq\! 1$ such~that~\hbox{$\int \log \norm{f_\varstar v } \, d\nu^{(n)}(f) \!\geq\! c_0$}, and then define inductively
\begin{equation*}
\tau_{k+1}( v, \omega) = \tau_k( v, \omega) + \tau_1( v_{\omega, k}, \sigma^k (\omega)).
\end{equation*}
By construction, $\tau_1$ depends on $v$ (hence on $x$) but not on $\omega$,
while $\tau_k$ depends on both $v$ and $\omega$ when $k\geq 2$; in addition
$\tau_k(v,\omega)\leq k n_0$ for all $k\geq 1$. For $n\geq 1$, define
$K_n( v,\omega)$, or $K(n)$ for short, by~$K_n( v,\omega)=\max\set{k \sep \tau_k \leq n}$. Then
$K(n) \geq n/n_0$ and $n- K(n)\leq n_0-1$.
With the convention $\tau_0 = 0$, the chain rule~\eqref{eq:chain_rule} gives
\begin{align*}\label{eq:stopping_bdd}
\notag\ee\lrpar{\log\norm{(f^n_\omega)_\varstar v}} &= \ee\bblrpar{\sum_{j = 0}^{K(n) -1} \log\bnorm{\blrpar{f_{\sigma^{\tau_j}\omega}^{\tau_{j+1} - \tau_j} }_\varstar v_{\omega, \tau_j}}} +
\ee\Blrpar{ \blrpar{f_{\sigma^{\tau_{K(n)}}\omega}^{n- \tau_{K(n)}} }_\varstar v_{\omega, \tau_{K(n)}}} \\
&\geq \frac{n}{n_0}c_0 - \max_{1\leq q\leq n_0} \ee\lrpar{ L (f_\omega^q) } \\
&\geq \frac{n}{n_0}c_0 - n_0 \int L(f)d\nu(f).
\end{align*}
Thus, for $n\geq \sfrac{n_0}{2} + \Psfrac{n_0^2}{c_0}\int L(f)d\nu(f)$, we~have $\ee\lrpar{\log\norm{(f^n_\omega)_\varstar v}} \geq \sfrac{c_0}{2}>0$ independently of $ v$, as was to be shown.
\end{proof}

\begin{lem}\label{lem:weak_UE2}
Under the moment condition \eqref{eq:moment+}, $\nu$ is uniformly expanding if and only
if
\begin{equation}\label{eq:weak_UE2}
\forall v \in T^1X, \exists c>0 \ \text{ such that} \quad
\pp\Blrpar{ \unsur{n}\log \norm{ (f^n_\omega)_\varstar v} \geq c} \underset{n\to\infty}{\to} 1.
\end{equation}
Under the moment condition~\eqref{eq:moment}, Property~\eqref{eq:weak_UE2} implies uniform expansion.
\end{lem}

\begin{proof}
Let us first show that~\eqref{eq:weak_UE2} implies~\eqref{eq:weak_UE} under the assumption~\eqref{eq:moment}.
Fix $ v\in T^1X$,
set $\Omega_n = \set{ \omega\in \Omega \sep  \unsur{n}\log \norm{ (f^n_\omega)_\varstar v} \geq c}$, and split
$\ee\lrpar{\unsur{n}\log \norm{ (f^n_\omega)_\varstar v}}$ into the sum of an integral over $\Omega_n$ and an integral over $\Omega_n^\complement$. The first one is larger than $c \pp(\Omega_n)$, and $\pp(\Omega_n)$ tends to $1$ as $n$ goes to $+\infty$. The second one satisfies
\[
\Babs{\ee\Blrpar{\unsur{n}\log \norm{ (f^n_\omega)_\varstar v} \mathbf 1_{\Omega_n^\complement} } }\leq
\ee \Blrpar{ \unsur{n} \mathrm L (f^n_\omega) \mathbf 1_{\Omega_n^\complement} } .
\]
The moment condition and Kingman's subadditive ergodic theorem show that $ \unsur{n} \mathrm L(f^n_\omega)$ is uniformly integrable and converges almost surely to some finite constant; since $ \mathbf \pp({\Omega_n^\complement}) $ converges to $0$, we~conclude that $ \ee\lrpar{\unsur{n}\log \norm{ (f^n_\omega)_\varstar v}}\geq c/2$ for large $n$.

For the converse implication we use a martingale convergence argument, as in~\cite[Lem.\,4.3.5]{liu_thesis} and~\cite[Prop.\,2.2]{chung}.\footnote{Chung only assumes the moment condition \eqref{eq:moment} however it seems to us that a stronger assumption is needed for the
control of the martingale differences.} Choose $p>1$ such that
(M$_{p}$) holds.
For convenience, let us first replace $\nu$ by $\nu^{(n_0)}$, where $n_0$ is given by the expansion property~\eqref{eq:UEbisrepetita}.
Define (for some fixed unit vector $v$)
\begin{equation}\label{eq:def_Xk}
X_k = \log\norm{ (f^1_{\sigma^k\omega})_\varstar v_{\omega, k}} -
\int \log\norm{ f_\varstar ( v_{\omega, k})} d\nu(f).
\end{equation}
These increments $ X_k $ are uniformly bounded in $L^p$ because
\begin{equation*}
\ee\lrpar{\abs{\log\norm{ (f^1_{\sigma^k\omega})_\varstar v_{\omega, k}}}^p }^{1/p} \leq
\ee\lrpar{ L(f^1_{\sigma^k\omega})^p}^{1/p} = \lrpar{\int L(f)^p d\nu(f)}^{1/p}
\end{equation*}
and the second term in~\eqref{eq:def_Xk} is pointwise bounded by
\begin{equation*}
\abs{\int \log\norm{ f_\varstar v_{\omega, k}} d\nu(f)} \leq \int L(f) d\nu(f) \leq \bblrpar{\int L(f)^p d\nu(f)}^{1/p}.
\end{equation*}
Thus, the sums $S_n = \sum_{k=0}^{n-1} X_k$ are all in $L^p$. Since $\ee({X_{n}\mid \mathcal F_{n}})= 0$ and $S_n$ is $\mathcal{F}_{n-1}$-measurable, $(S_n)$ is a martingale relative to the filtration $(\mathcal F_{n-1})$.
It follows from Theorem 2.22 in \cite[\S 2.7]{hall-heyde} that
$\frac{1}{n}S_n$ converges to $0$ in probability and in $L^p$.
Now, the chain rule gives
\begin{equation*}
\frac{1}{n} S_n(\omega)=
\unsur{n}\log \norm { (f^n_\omega)_\varstar v} - \unsur{n} \int \log \norm { f_\varstar v} d\nu^{(n)}(f),
\end{equation*}
and \eqref{eq:UE_iteree} asserts that
$\int \log {\norm{f_\varstar v}} d\nu^{(n)}(f) \geq c n$, so we conclude that for any $c'<c$
\begin{equation}\label{eq:proba_expansion_goes_to_1}
\pp\Blrpar{ \unsur{n}\log \norm{ (f^n_\omega)_\varstar v} \geq c'} \underset{n\to\infty}{\to} 1,
\end{equation}
as desired. Recall however that we are working with $\nu^{(n_0)}$:
coming back to $\nu$ this means that \eqref{eq:proba_expansion_goes_to_1} holds
along the subsequence $(nn_0)$. We~then write $n = kn_0+ r$, with $0\leq r \leq n_0-1$, so that
\begin{equation*}\label{eq:differential}
{(f^{n }_\omega)_\varstar v} = (f^{r}_{\sigma^{kn_0} \omega})_\varstar
(f^{kn_0}_\omega)_\varstar v
\end{equation*}
and what we have to show is that applying $f^{r}_{\sigma^{kn_0} \omega}$
does not affect
the linear growth of $ \log\norm{ (f^{kn_0}_\omega)_\varstar v }$.
But the inequality~\eqref{eq:Mp_iteree}, applied with $p=1$, gives
\begin{equation*}
\P \Blrpar{\exists 0\leq r\leq n_0-1, \ \abs{\log\norm{ ( f^{r}_{\sigma^{kn_0}})_\varstar }}\geq \e k}
\leq \sum_{r=0}^{n_0-1}
\nu^{(r)}\lrpar{ L(f) \geq \e k}
\leq \frac{Cn_0^2}{\e k} ,
\end{equation*}
and we are done.
\end{proof}

\begin{rem} \label{rem:given_v} In the first part of the proof,
the implication~$\eqref{eq:weak_UE2}\Rightarrow\eqref{eq:weak_UE}$ is true for a given $v$, while the converse implication requires uniform expansion on the whole of~$X$.
\end{rem}

\begin{rem}
This proof shows that if $\nu$ satisfies (M$_2$), then the convergence in probability in~\eqref{eq:weak_UE2} can be replaced by an almost sure convergence. (Indeed by \cite[Th.\,3, p.\,243]{Feller:vol2},
$\frac{1}{n}S_n$ converges almost surely to $0$ when the $X_k$ are uniformly~$L^2$.)
\end{rem}

\begin{rem}
So far, we~have not really used that we are dealing with diffeomorphisms: the results from this section
hold for a semigroup action, by~replacing $\mathrm L(f)$ by $\norm{f}_{C^1(X)}$.
\end{rem}

\section{Inducing on a finite index subgroup} \label{sec:inducing}

\subsection{Hitting times and hitting measures (see \cite[Chap.\,5]{benoist-quint_book}) }\label{par:hitting_measure}

Let $\nu$ be a Borel
probability measure on $\Diff^1(M)$ and let $G$ be the
closed sub-semigroup of $\Diff^1(M)$ generated by $\nu$.
Let $H\subset G$ be a closed finite index
sub-semigroup; this means that
there is a continuous and transitive action $G\times F \to F$ on some finite set $F$ such that $H$ is the stabilizer of some element $x_0\in F$; the index of
$H$ is $[G:H]=\abs{F}$ and $F$ is the quotient space.
For instance, $H$ can be the stabilizer of a point $x$ in a finite $G$-orbit.

The hitting time $T_{H}$ of $H$ for the random walk
induced by $\nu$ (starting from the neutral element) is
\begin{equation*}
T_H(\omega) = \min\bset {n\geq 1\sep f^n_\omega \in H}.
\end{equation*}
Lemmas 5.4 and 5.5 in~\cite{benoist-quint_book} show that $T_H$ is almost surely finite,
admits an exponential moment, and satisfies $\ee(T_{H}) = [G:H]$. By~definition the \emph{hitting measure} (or induced measure)
$\nu_H$ is the probability measure on $H$ describing the distribution of $f^{T_{H}(\omega)}_\omega$.

Define the \emph{$k$-th hitting time} $T_{H,k}$ of $H$ by $T_{H,1}=T_H$ and the induction\vspace*{-3pt}
\begin{equation*}
T_{H,k+1}(\omega) = \min\aset{n\geq T_{H,k}+1\sep  f_\omega^{n} \in H}.
\end{equation*}
The convolution $\nu_H^{(k)}$ describes the distribution
of $f^{T_{H,k}(\omega)}_\omega$.
If $H$ is a finite index semigroup and $g\in H$, $hg$ belongs to $H$ if and only if $h$ belongs to~$H$. Thus,
$T_{H, k+1}(\omega) - T_{H, k}(\omega) = T_{H, 1}(\sigma^{T_{H, k}(\omega)}(\omega))$ and the Markov property implies that the random variables $(T_{H, k+1} - T_{H, k})$ are independent and identically distributed: each of them is distributed as $T_H$. Since their expectation equals $[G:H]$,
the law of large numbers gives\vspace*{-3pt}\enlargethispage{.6\baselineskip}
\begin{equation*}\label{eq:limit_kth_hitting_time}
\lim_{k\to+\infty} \frac{1}{k}T_{H, k}(\omega)=[G:H]
\end{equation*} $\nu^\N$-almost surely.

\begin{thm}\label{thm:M$_2$} The hitting measure on a finite index subgroup satisfies the following properties
\begin{enumerate}
\item\label{thm:M$_2$1} if $\nu_H$ satisfies~\eqref{eq:momentp} for some $p\geq 1$, then so does $\nu$;
\item\label{thm:M$_2$2} if $\nu$ satisfies~\eqref{eq:momentp}, then $\nu_H$ satisfies (M$_{p'}$) for any $1 < p' < p$;
\item\label{thm:M$_2$3} $\nu$ satisfies~\eqref{eq:moment}, or~\eqref{eq:moment+}, or~\eqref{eq:momentexp} if and only if $\nu_H$ does.
\end{enumerate}
Moreover, $\nu_H$ generates $H$ as a semigroup, which means that $H$ is the smallest closed sub-semigroup of $G$ containing the support of $\nu_H$.
\end{thm}

This result still holds if we substitute any subbaditive function to $ \log\norm{f}_{C^1(X)}$ in the definition of $L$ (see Equation~\ref{eq:subadditive_function_L}), with exactly the same proof.
\begin{proof}
Consider the finite quotient $F$ of $G$ by $H$ and denote the action of~$G$ on~$F$ by left translations by ($u\mto au$, $a\in G$); by definition $H$ is the stabilizer of some $x_0\in F$. Set $K=\vert F\vert = [G: H]$.

For each $u\in F$, choose a sequence of measurable subsets $A_1(u)$, $A_2(u)$, $\ldots$, $A_k(u)$ in $G$, with $k=k(u)\leq K$
such that $\nu(A_i(u))>0$ for each $i$ and, for all sequences $a_i\in A_i(u)$, $(a_k \cdots a_1) u= x_0$
while $(a_j \cdots a_1) u\neq x_0$ if $j<k$. Since $F$ is finite, there is a real
number $\e >0$ such that $\nu(A_1(u))\cdots \nu(A_{k(u)}(u))$ is larger than $\e$ for all $u$. Shrinking the $A_i(u)$ if necessary, we~may assume that $\mathrm L(g)\leq C$ for some $C>0$ and all $g$ in $\bigcup_{u,i}A_i(u)$.

We split the integral of $ \mathrm L(f)^p$ as a finite sum\vspace*{-3pt}
\[
\int \mathrm L(f)^pd\nu(f) = \sum_{u\in F}\int_{\set{fx_0=u }} \mathrm L(f)^pd\nu(f).
\]
If\vspace*{-5pt}
\begin{equation*}
(a_{k(u)}, \ldots, a_1)\in A_{k(u)}(u)\times \cdots \times A_1(u),
\end{equation*}
then $\mathrm L(f)\leq \mathrm L(a_{k(u)} \cdots a_1 f)+K C$ because $\mathrm L$ is subadditive; thus,\vspace*{-3pt}
\begin{align*}
\int_{\set{fx_0=u }} \mathrm L(f)^pd\nu(f) & \leq \int_{\set{fx_0=u }} (\mathrm L(a_{k(u)} \cdots a_1 f)+KC)^pd\nu(f).
\end{align*}
By construction, the product $a_{k(u)} \cdots a_1 f$ is a first return in $H$. Thus,
integrating over the $A_i(u)$, the distribution of $a_{k(u)} \cdots a_1 f$ contributes positively
to $\nu_H$, and we get
\begin{align*}
\e \int_{\set{fx_0=u }} \mathrm L(f)^pd\nu(f) & \leq \int_H (\mathrm L(g)+KC)^pd\nu_H(g).
\end{align*}
Assertion \eqref{thm:M$_2$1} follows from this estimate.

For Assertion \eqref{thm:M$_2$2}, we~must bound
$\int \mathrm L(f)^{p'} d\nu_H(f)= \ee ( \mathrm L(f_\omega^{T_H(\omega)})^{p'} )$. By~subadditivity of $\mathrm L$ and convexity of $s\mto s^{p}$, we~have
\begin{align*} \label{eq:int_T}
\mathrm L\blrpar{f_\omega^{T_H(\omega)}}^{p}
& \leq T_H(\omega)^{p-1} \sum_{i=0}^{T_H(\omega)-1} {\mathrm L(f_i)}^{p}.
\end{align*}
Raising this inequality to power $p'/p$ gives
\begin{align*}
\mathrm L\blrpar{f_\omega^{T_H(\omega)}}^{p'} & \leq T_H(\omega)^{p'(1-1/p)}
\bblrpar{\sum_{i=0}^{T_H(\omega)-1} {\mathrm L(f_i)}^{p}}^{p'/p}.
\end{align*}
On the other hand, Lemma 5.4 of \cite{benoist-quint_book} says that $\ee(\sum_{i=1}^{T_H(\omega)}\varphi \circ\sigma^i ) = \ee(T_H)\ee(\varphi)$, for any integrable function~$\varphi$. This shows that $\sum_{i=0}^{T_H(\omega)-1} {\mathrm L(f_i)}^{p}$ is integrable, and in particular its $p'/p$ power
is in $\mathrm L^{p/p'}(\Omega; \nu^\N)$. Since the hitting time $T_H$ admits moments of all orders, we~can apply the Hölder
inequality with parameters $r=p/p'>1$ and $q$ such that $1/q+1/r=1$: it shows that $\mathrm L\blrpar{f_\omega^{T_H(\omega)}}^{p'} $ is integrable, as desired.

Assertion \eqref{thm:M$_2$3} follows from Assertions \eqref{thm:M$_2$1} and \eqref{thm:M$_2$2} and \cite[Cor.\,5.6]{benoist-quint_book}.
For the last assertion, fix an element $h$ of $H$ and an open neighborhood $U$ of $h$ in $H$.
Since $H$ is of finite index in $G$, it is open and closed, so $U$ is also a neighborhood of $h$ in $G$. By~assumption the support of $\nu$ generates a dense sub-semigroup of
$G$, so the random walk induced by $\nu$ starting at the neutral element visits $U$,
thus $\nu_H$ generates $H$ as a semigroup. \end{proof}

\Subsection{Uniform expansion of the induced measure}

\begin{pro}\label{pro:induced_UE}
Let $\nu$ be a probability measure on $\Diff^1(M)$ satisfying \eqref{eq:moment}. Assume that $\nu$ is uniformly expanding and
let $n_0$ be as in~\eqref{eq:UEbisrepetita}. Then, the measure induced by
$\nu^{(n_0)}$ on $H$ is uniformly expanding.
\end{pro}

In fact,
Proposition~\ref{pro:induced_UE2} below shows that, under condition~\eqref{eq:moment+}, $\nu$ is uniformly expanding if
and only if $\nu_H$ is. The proof of Proposition~\ref{pro:induced_UE} is based on a simple martingale argument, while
Proposition~\ref{pro:induced_UE2} relies on the criterion of Lemma~\ref{lem:weak_UE2}.

\begin{proof} We use ideas from \cite[\S 4.3]{liu_thesis} and \cite[Prop.\,2.2]{chung}.
To ease notation we rename~$\nu^{(n_0)}$ into $\nu$ so that~\eqref{eq:UEbisrepetita} holds with $n_0= 1$ and some $c>0$; as above, we~denote by $\nu_H$ the measure induced by $\nu$ (\ie by $\nu^{(n_0)}$) on $H$. Fix
$ v\in T^1X$, and
define a sequence of random variables $(Y_k)_{k\geq 0}$ by
\begin{equation*}
Y_k(\omega) = \log\norm{ (f^1_{\sigma^k\omega})_\varstar v_{\omega, k}}- c.
\end{equation*}
Then for all $k\geq 1$,
$\ee({Y_{k}\mid \mathcal F_{k}})\geq 0$, so that the sequence $(S_n)_{n\geq 1}$
defined by $ S_n = \sum_{k=0}^{n-1} Y_k$ is a submartingale
relative to the filtration $(\mathcal F_n)$: $\ee( S_{n+1}\mid \mathcal F_n)\geq S_{n}$.
The moment condition~\eqref{eq:moment} implies that
$\ee(\abs{S_{n+1}-S_n}\mid \mathcal F_n) = \ee(\abs{Y_n}\mid \mathcal F_n)$ is uniformly bounded.
Since the hitting time $T_H$ is integrable, we~can apply the optional stopping theorem \cite[Th.\,4.7.5]{durrett_book}, which implies that
$\ee(S_{T_H})\geq \ee(S_1)\geq 0$. Unwinding the definitions and applying the chain rule, we~see that
\begin{equation*}
\ee(S_{T_H}) = \int \log \norm{ f_\varstar v}\, d\nu_H(f) - c [G:H],
\end{equation*}
where we use $\ee (T_H) = [G:H]$. Therefore
$\int \log \norm{f_\varstar v}\, d\nu_H(f) \geq c [G:H]>0$, and
$\nu_H$ is uniformly expanding.
\end{proof}

\begin{pro}\label{pro:induced_UE2}
Let $\nu$ be a probability measure on $\Diff^1(M)$ satisfying \eqref{eq:moment+}. Let~$\nu_H$ be the measure induced on a closed
finite index sub-semigroup. Then $\nu$ is uniformly expanding if and only if $\nu_H$ is uniformly expanding.
\end{pro}

\begin{proof}
Let us show that if $\nu_H$ is uniformly expanding then $\nu$ is uniformly expanding.
The converse implication is similar and is left to the reader (in
this direction, Proposition~\ref{pro:induced_UE} will actually be sufficient for our purposes).
Fix $ v\in T^1M$. In~view of Lemma~\ref{lem:weak_UE2}, we~have to show that for some $c>0$,
\begin{equation*}\label{eq:expansion_n}
\pp\Blrpar{ \unsur{n}\log \norm{ (f^n_\omega)_\varstar v} \geq c}
\underset{n\to\infty}{\to} 1.
\end{equation*}
Consider the sequence of hitting times $(T_{H,k})$ defined in Section \ref{par:hitting_measure} and
denote it by $(T_k)$
for simplicity (hence $T_1=T_H$). By~Theorem~\ref{thm:M$_2$},
$\nu_H$ satisfies \eqref{eq:moment+}, so we can apply
Lemma~\ref{lem:weak_UE2} to get a real number $c>0$ such that
\begin{equation}\label{eq:expansion_Tk}
\pp\Blrpar{ \unsur{k} \log\anorm{\alrpar{f^{T_{k}(\omega)}}_\varstar v} \geq c} \underset{k\to\infty}{\to} 1.
\end{equation}
Let $\gamma = [G:H]$ and fix a positive real number
$\e < \sfrac{c}{\gamma}$. Let also $\e_1\ll\e$ which will be specified later.
For $K\geq 1$, set
\begin{equation*}
\Omega_1(K)=\Bset{ \omega \sep  \forall k\geq K, \; \Babs{\frac{T_k(\omega)}{k} -\gamma} < \e_1 }.
\end{equation*}
If $K>\gamma/(2\e_1)$ and $\omega\in \Omega_1(K)$, then for $n\geq K\gamma+1$ we get
\begin{equation}\label{eq:estimate_on_T}
T_{\lfloor n/(\gamma+\e_1)\rfloor}(\omega)- \frac{4\e_1}{\gamma}n \leq n \leq T_{\lfloor n/(\gamma+\e_1)\rfloor}(\omega)+ \frac{4\e_1}{\gamma}n.
\end{equation}

Now, define $\Omega_2(n )$ to be the set of sequences $\omega$ such that
the inequality involved in \eqref{eq:expansion_Tk} is satisfied at time $T_{\lfloor n/(\gamma+\e_1)\rfloor}(\omega)$; in other words, $\omega\in \Omega_2(n)$ if and only if
\begin{equation}\label{eq:estimate_on_T2}
\log\bnorm{\blrpar{f^{T_{\lfloor n/(\gamma+\e_1)\rfloor}(\omega)}}_\varstar v} \geq c \lfloor n/(\gamma+\e_1)\rfloor.
\end{equation}
Then $\pp(\Omega_1(\lfloor\sqrt{n}\rfloor)\cap \Omega_2(n))$ converges to $1$ as $n$ goes to $+\infty$
and as soon as $n\geq \gamma\sqrt{n}+1$, any $\omega\in \Omega_1(\lfloor\sqrt{n}\rfloor)\cap \Omega_2(n)$
satisfies~\eqref{eq:estimate_on_T} and~\eqref{eq:estimate_on_T2}.

Now consider the set $\Omega_3(n)\subset \Omega_1(\lfloor\sqrt{n}\rfloor)\cap \Omega_2(n)$ made of those
$\omega$ such that $\mathrm L\blrpar{ f^{n-T_{k}}_{\sigma^{T_{k} }\omega}} < \e n$ for $k= \lfloor n/(\gamma+\e_1)\rfloor $. Then
\begin{equation*}
\pp(\Omega_3(n)^\complement)\leq \pp\blrpar{\mathrm L\blrpar{ f^{n-T_{k}}_{\sigma^{T_{k} }\omega}} \geq \e n}\leq
\pp\Blrpar{\max_{0\leq q\leq \sfrac{4\e_1n}{\gamma}} \ts\sum_{i=0}^{q-1} L_i \geq \e n},
\end{equation*}
where $(L_i)_{i\geq 0}$ is a sequence of independent random variables, each of which being
distributed as $\mathrm L(g)$ for $d\nu(g)$. Since the $L_i$ are non-negative, $\pp(\Omega_3(n)^\complement)\leq\pp\blrpar{ \sum_{i=0}^{4\e_1 n/\gamma} L_i \geq\nobreak \e n}$. If $\e_1$ is chosen such that $\Psfrac{4\e_1}{\gamma} \ee(L_1) < \e$, the law of large numbers entails that
\begin{equation*}
\pp \Blrpar{ \sum_{i=0}^{4\e_1 n/\gamma} L_i \geq \e n} \underset{n\to\infty}\to 0,
\end{equation*}
thus $ \pp \lrpar{ \Omega_3(n)}$ tends to $1$ as $n\to \infty$.

Then, for $\omega \in \Omega_3(n)$, the estimates~\eqref{eq:estimate_on_T},~\eqref{eq:estimate_on_T2}
and $\mathrm L\alrpar{ f^{n-T_{k}}_{\sigma^{T_{k} }\omega}} < \e n$ imply that
$\unsur{n}\log \norm{(f^n_\omega)_\varstar v}\geq \sfrac{c}{\gamma}-\e$
and the conclusion follows.
\end{proof}

\section{Margulis functions}\label{sec:margulis}

In this section we develop some tools for the proof of the equidistribution
Theorem~\ref{thm:equidistribution_wehler}.
Under appropriate assumptions, we~show that the measures
$\nu^n\ast \delta_x$ and $\unsur{n}\sum_{k=1}^n \delta_{f^k_\omega(x)}$
do not cluster at a $\Gamma$-periodic orbit, except when $\Gamma(x)$ is itself finite.
The basic tool is the construction of a proper function, defined on the complement of such a periodic orbit,
which ``essentially decreases''
along random trajectories. After~\cite{eskin-margulis}
it is often
referred to as a \emph{Margulis function}, even if this strategy has a long history in the Markov chain literature (see
\cite{meyn-tweedie}). Our presentation is greatly influenced by \cite{benoist-quint3} and
\cite{benard-desaxce}.

\subsection{A general recurrence criterion}
For concreteness, instead of general Mar\-kov chains, we~consider the setting of group actions.

\begin{thm}[Bénard-De Saxcé \cite{benard-desaxce}]\label{thm:margulis} Let $U$ be a locally compact topological space. Let $\Gamma$ be a group of homeomorphisms of $U$, and $\nu$ be a probability measure on $\Gamma$.
Assume that there exists a function $u:U\to \R_+$ satisfying the assumptions:
\begin{eqnarray}
\label{eq:weak_margulis}
\exists A>0, \ \exists \gamma>0, \ \forall x\in U, \ u(x)\geq A \implies
\int u(f(x))d\nu(f) \leq u(x) - \gamma, \\
\label{eq:margulis_moment}
\exists B>0, \exists \eta>0, \forall x\in U, \quad \int \abs{u(f(x)) - u(x) } ^{1+\eta} d\nu(f) \leq B.
\end{eqnarray}
Then for every $\e>0$ there exists $R>0$ such that for all $x$ in $U$,
\begin{enumerate}
\item\label{thm:margulis1} there exists $n_x\geq 0$, such that $(\nu^{n}\ast \delta_x )(\set{u\geq R}) \leq \e$ for all $n\geq n_x$;
\item\label{thm:margulis2} for $\nu^\N$-almost every $\omega$,
\[
\limsup_{n\to\infty} \unsur{n} \# \set{k\in \set{1, \ldots, n} \sep  u \lrpar{f^k_\omega(x) }\geq R} \leq \e.
\]
\end{enumerate}
Furthermore the integer $n_x$ in \eqref{thm:margulis1} depends only on $u(x)$.
\end{thm}

In the most interesting cases, $u$ will be a proper function on $U$. Then, Equation~\eqref{eq:weak_margulis}
expresses that, on average, the random dynamics does not send points too far off at infinity,
and Equation~\eqref{eq:margulis_moment} can be understood as a moment condition in the $u$\nobreakdash-vari\-able; then, the conclusions \eqref{thm:margulis1} and \eqref{thm:margulis2} correspond to ``non-escape of mass'' and ``quantitative recurrence'' properties.

For the proof, see \cite[Prop.\,1.2]{benard-desaxce}, and the comments following it. More
precisely, we~refer to~\cite[Prop.\,2.5]{benard-desaxce} for the conclusion~\eqref{thm:margulis1},
including the uniformity statement on~$n_x$, and to~\cite[Prop.\,2.7]{benard-desaxce} for~\eqref{thm:margulis2}.

The original decay property for the Margulis function $u$ in
\cite{eskin-margulis, benoist-quint3}
is
\begin{equation}
\label{eq:margulis}
\exists 0<a<1, \ \exists b>0, \ \forall x\in U, \quad \int u(f(x))d\nu(f) \leq au(x) +b
\end{equation}
instead of~\eqref{eq:weak_margulis}. One easily checks that if $u$ satisfies~\eqref{eq:weak_margulis} and the following strong integrability property
\begin{equation}
\label{eq:margulis_moment_exp}
\exists B>0,\; \forall x\in U, \quad \int \exp\lrpar{\abs{u(f(x)) - u(x) } }d\nu(f) \leq B,
\end{equation}
then $e^{\delta u}$
satisfies~\eqref{eq:margulis} for small $\delta>0$. Under this assumption, Theorem~\ref{thm:margulis} was established in
\cite{benoist-quint3}.

\subsection{Finite orbits of $C^2$ actions}\label{par:finite-orbits-C2-actions}
Let $\nu$ be a probability measure on the group of $C^2$ diffeomorphisms of a compact Riemannian
manifold $M$ of dimension $d$.
As in Section \ref{subs:moment} we consider the moment conditions
\begin{align}
\label{eq:moment2p}
\tag{M$_{2, p}$}
\int \lrpar{ \log\norm{f}_{C^2 }+ \log\norm{f\inv}_{C^2 }} ^p \, d\nu (f) < + \infty, \\
\label{eq:moment2+}
\tag{M$_{2, +}$}
\exists p>1, \ \mathrm{(M}_{2,p}\mathrm{)} \text{ holds}, \quad \quad \quad\quad \quad \quad \quad\quad
\end{align}
\begin{rem}\label{rem:moment-condition-holo}
For a holomorphic action on a compact
complex manifold $X$, these conditions are equivalent to their respective
$C^1$ analogues~\eqref{eq:momentp} and~\eqref{eq:moment+}, because a uniform control on the first derivatives provides a uniform control on the higher derivatives as well.
Here is an outline of the argument.
Cover $X$ by finitely many charts $\Omega_i$.
Then there exists $c>0$ such that for every $x\in X$ and every $f\in \Aut(X)$, the balls
$B(x,r)$ and $f(B(x,r))$
are contained in a single chart, as soon as $r \leq c \norm{f}_{C^1}\inv$. Then the
Cauchy estimates imply that $\norm{f\rest{B(x, r/2)}}_{C^2}\leq (C/r) \norm{f\rest{B(x, r)}}_{C^1} \leq C \norm{f}_{C^1}^2$ and the result follows.
\end{rem}

Before stating our next result, recall that the notion of uniform expansion along a finite orbit was defined in Section~\ref{subs:UE_intro}.

\begin{thm}\label{thm:margulis_UE}
Let $\Gamma$ be a group of $C^2$ diffeomorphisms of a compact Riemannian
manifold $M$, and $\nu$ be a measure on
$\Gamma$ satisfying the moment condition (M$_{2,+}$).
Let~$F$ be a finite orbit of $\Gamma$ such that $\nu$ is uniformly expanding on $F$.
Then for every $x\in M\setminus F$, for every $\e>0$
there exists a compact set $K\Subset M\setminus F$ such that:
\begin{enumerate}
\item\label{thm:margulis_UE1} $(\nu^{n}\ast \delta_x) (K) \geq 1- \e$ for $n\geq n_x$, and
\item\label{thm:margulis_UE2} for $\nu^\N$-almost every $\omega$,
\[
\limsup_{n\to\infty} \unsur{n} \# \set{k\in \set{1, \ldots, n}\sep f^k_\omega(x) \in K} \geq 1-\e.
\]
\end{enumerate}
Furthermore the integer $n_x$ in \eqref{thm:margulis_UE1} is locally uniform in $M\setminus F$.
\end{thm}

This result seems to be new: it appears under stronger (exponential) moment assumptions
in \eg \cite{liu_thesis, chung}. Note that such a result is not expected to hold under the
first moment condition (M$_{2, 1}$), as explained in
Examples 1 and 2 of Section~2 in~\cite{benard-desaxce}.

\begin{proof}
First, the proof of Proposition 3.3 in~\cite{benard-desaxce} shows that if
the conclusions \eqref{thm:margulis_UE1} and \eqref{thm:margulis_UE2} hold for $\nu^{(n_0)}$, then they hold for $\nu$. So we can replace $\nu$ by $\nu^{(n_0)}$ and hence assume that the uniform expansion property~\eqref{eq:UE} holds (on $F$) for $n_0=1$.

Let $d(\cdot,\cdot)$ be the Riemannian distance on $M$. According to Theorem~\ref{thm:margulis}, we~only need to show that
$u\colon x\mto - \log d(x, F)$ is a
proper function $M\setminus F\to \R_+$ satisfying Properties~\eqref{eq:weak_margulis}
and~\eqref{eq:margulis_moment}.

\subsubsection*{Preliminaries}
We set
$N(f)=\norm{f}_{C^2}+\norm{f\inv}_{C^2}$ and note that $N(f)\geq \Lip(f) + \Lip(f\inv)$ for every $f\in \Gamma$. In~particular, for every $x\in X$
\begin{equation}
\label{eq:lipschitz}
\frac{1}{N(f)} \leq \frac{d (f(x), F)}{d(x,F)} \leq N(f).
\end{equation}
For $R>0$, set $\Gamma(R)=\{f\in \Gamma \sep  N(f) \leq R\}$. We~choose $\eta >0$ such that the moment condition~\eqref{eq:moment2p} is satisfied with $p=1+\eta$. Then,
\begin{equation*}
I_\eta:=\int_{\Gamma} \left(\log(N(f))\right)^{1+\eta} d\nu(f)
\end{equation*}
is a finite positive number. In~what follows, we~choose $R>1$ such that
\begin{equation*}
\frac{2I_\eta}{(\log(R))^\eta}< \frac{c}{4},
\end{equation*}
where $c$ is the expansion factor in Equation~\eqref{eq:UE} (along the finite orbit $F$).

Take $s>0$ such that
\begin{itemize}
\item $s$ is smaller than the injectivity radius of $M$ at $y$, for every $y\in F$;
\item the balls $B(y;s)$, for $y$ in $F$, are pairwise disjoint;
\item $C_0R^2 s < c/4$, where $c$ is the expansion factor as above, and $C_0$ is the constant appearing
below in the Taylor expansion (Equation~\eqref{eq:Taylor}).
\end{itemize}
Then, define $V$ and $V'$ by
\begin{equation*}
V=\bigcup_{y\in F} B(y;s), \quad V' = \bigcup_{y\in F} B(y;s/R).
\end{equation*}
By~\eqref{eq:lipschitz} we have $f(V')\subset V$ for every $f\in \Gamma(R)$.

If $x$ belongs to $V$, we~denote by $\pi(x)$ the unique point of $F$ at distance $\leq s$ from~$x$, and we denote by $w_x$ the unique vector in $T_{\pi(x)}M$ such that $\exp_{\pi(x)}(w_x)=x$ and $\norm{w_x}=d(x,\pi(x))$.

\subsubsection*{First estimate}
For $f$ in $\Gamma(R)$ and $x\in V'$, Taylor's second order formula yields
\begin{equation}\label{eq:Taylor}
\babs{d(f(x),f(\pi(x))) - \norm{f_\varstar (w_x)}}\leq C_0 N(f)d(x, \pi(x))^2,
\end{equation}
for some uniform constant $C_0$, that does not depend on $f$. This gives
\begin{equation*}
\Babs{\frac{d(f(x),F)}{d(x, F)} - \frac{\norm{f_\varstar (w_x)}}{\norm{w_x}}} \leq C_0 N(f) d(x,F).
\end{equation*}
Using the Lipschitz estimate~\eqref{eq:lipschitz} and the fact that $\abs{\log(a) -\log(b)}\leq N\abs{a-b}$ when $a, b\in [N\inv,\infty[$, we~obtain
\begin{equation}\label{eq:delta_1}
\Babs{\log\Blrpar{\frac{d(f(x),F)}{d(x, F)}} - \log\Blrpar{\frac{\norm{f_\varstar (w_x)}}{\norm{w_x}}}} \leq C_0 N(f)^2 d(x,F).
\end{equation}
By the definition of $\Gamma(R)$ and the requirements on $s$, we~get
\begin{align*}
\int_{f\in \Gamma(R)} \Babs{\log\Blrpar{\frac{d(f(x),F)}{d(x, F)}}- \log\Blrpar{\frac{\norm{f_\varstar (w_x)}}{\norm{w_x}}}} d\nu(f) & \leq C_0 R^2 d(x,F) \leq \frac{c}{4},
\end{align*}
because $d(x,F)\leq s$.

\subsubsection*{Second estimate}
Now, for any $f$ in $\Gamma$ we also have
\begin{equation}\label{eq:lipschitz_log_ratio_estimate}
\Babs{\log\Blrpar{\frac{d(f(x),F)}{d(x, F)}} - \log\Blrpar{\frac{\norm{f_\varstar (w_x)}}{\norm{w_x}}}} \leq 2\log(N(f))
\end{equation}
hence Markov's inequality and our choice of $R$ give
\begin{align*}
\int_{f\in \Gamma(R)^\complement} \Babs{\log\Blrpar{\frac{d(f(x),F)}{d(x, F)}} - \log\Blrpar{\frac{\norm{f_\varstar (w_x)}}{\norm{w_x}}}} d\nu(f)& \leq \frac{2}{\log(R)^\eta}I_\eta \leq \frac{c}{4}.
\end{align*}

\subsubsection*{Conclusion}
Summing the integrals over $f$ in $\Gamma(R)$ and $\Gamma(R)^\complement$, we~obtain
\begin{equation}\label{eq:distanceF}
\int_{f\in \Gamma} \Babs{\log\Blrpar{\frac{d(f(x),F)}{d(x, F)}} - \log\Blrpar{\frac{\norm{f_\varstar (w_x)}}{\norm{w_x}}}} d\nu(f) \leq \frac{c}{2}.
\end{equation}
Since $w_x$ is a vector tangent to $M$ at $\pi(x)\in F$,
the uniform expansion along $F$ yields
\begin{equation*}\int \log\Blrpar{\frac{\norm{f_\varstar (w_x)}}{\norm{w_x}}} d\nu(f) \geq c
\end{equation*}
and
then~\eqref{eq:distanceF} implies that
\begin{equation*}
\int - \log {d(f(x),F)} d\nu(f) \leq - \log(d(x, F)) - c/2.
\end{equation*}
In other words, $u\colon x\mto -\log(d(x,F))$ satisfies Property~\eqref{eq:weak_margulis} (with $A=-\log(s)$). Property~\eqref{eq:margulis_moment} is obtained from~\eqref{eq:lipschitz} and the moment condition.
Thus, as announced above, $u$ satisfies the assumptions of Theorem~\ref{thm:margulis}, and we are done. \end{proof}

The local uniformity of $n_x$ in Theorem~\ref{thm:margulis_UE} has the following interesting consequence.

\begin{pro}\label{pro:finite_mass}
Under the assumptions of Theorem~\ref{thm:margulis_UE}, any stationary Radon measure on $M\setminus F$ has finite mass.
\end{pro}

\begin{proof}
Let $\mu$ be such a stationary measure. Fix $\e>0$, say $\e =1/2$ and
let $K$ be as in Theorem~\ref{thm:margulis_UE}. The stationarity of $\mu$ implies that for every $n\geq 0$,
\begin{equation*}
(\nu^{(n)}\times\mu)\lrpar{ \set{(g, x)\sep gx \in K}} = \mu(K),
\end{equation*}
hence for every Borel set $B\subset M\setminus F$,
\begin{equation}\label{eq:margulis_uniformity}
\int_B \nu^{(n)}\lrpar{\set{g\sep gx\in K}} d\mu(x) \leq \int_X \nu^{(n)}\lrpar{\set{g\sep gx\in K}} d\mu(x) = \mu(K).
\end{equation}
Now if $B$ is an arbitrary compact subset of $M\setminus F$, the uniformity statement in
Theorem~\ref{thm:margulis_UE} implies that there exists $n = n_B$ such that
for every $x\in B$,
\begin{equation*}
\nu^{(n_B)}\lrpar{\set{g\sep gx\in K}}\geq \frac12.
\end{equation*}
Plugging this into~\eqref{eq:margulis_uniformity}, we~obtain $\frac12 \mu(B) \leq \mu(K)$.
Since $B$ is arbitrary, this implies that $\mu(M\setminus F)\leq 2\mu(K)$ and we are done.
\end{proof}

\subsection{Totally real invariant manifolds}
We now consider a situation which is specific to the complex setting.

\begin{thm}\label{thm:margulis_totalement_reel}
Let $X$ be a compact complex manifold of dimension $d$. Let $\Gamma$ be a group of holomorphic diffeomorphisms of $X$,
endowed with a probability measure $\nu$ satisfying~\eqref{eq:moment+}. Let $Y\subset X$
be a $\Gamma$-invariant,
real analytic, totally real submanifold of maximal (real) dimension $d$, such that
$\nu$ is uniformly expanding on $Y$.
Then for any
$x\in X\setminus Y$ and any $\e>0$, there exists a compact subset $K\Subset X\setminus Y$ such that the
conclusions \eqref{thm:margulis_UE1} and \eqref{thm:margulis_UE2} of Theorem~\ref{thm:margulis_UE} hold.
\end{thm}

By ``uniformly expanding along $Y$'' we mean that the restriction of $\Gamma$ to $Y$ is
uniformly expanding viewed as an action on $Y$, or equivalently
that the uniform expansion condition~\eqref{eq:UE} holds in $X$ for every $x\in Y$; the equivalence between the two conditions comes from the fact that for every $x\in Y$, the
complex span of $T_xY$ is $T_xX$. When $Y$ is singular, we~require that~\eqref{eq:UE} holds in $X$ along
$\Sing(Y)$.

Note also that this statement is specific to totally real submanifolds and holomorphic actions. In~other words,
there is no analogue of Theorem~\ref{thm:margulis} when $F$ is replaced by an arbitrary submanifold:
see Example~\ref{eg:margulis_submanifold} below.

\begin{proof}
We suppose $Y$ smooth and show that there exists $n\geq 1$ such that
$x\mto - \log d(x, Y)$ defines a
Margulis function (\ie satisfies~\eqref{eq:weak_margulis} and~\eqref{eq:margulis_moment})
for $\nu^{(n)}$. Then, as explained before, \cite{benard-desaxce} shows that \eqref{eq:weak_margulis} and~\eqref{eq:margulis_moment} are automatically satisfied with $n=1$. As in Theorem~\ref{thm:margulis_UE}, Property~\eqref{eq:margulis_moment} follows from the invariance of $Y$ and
the bilipschitz property; so we focus on~\eqref{eq:weak_margulis}.

For every $x\in Y$ there exists a local chart in which
the equation of $Y$ becomes $\Ima (z) = 0$, where $\Ima(z) = \Ima (z_1, \ldots , z_d) =
(\Ima (z_1), \ldots , \Ima (z_d))$ (see \cite[Prop.\,1.3.8 \& 1.3.11]{BER}). We~fix
a finite family $\phi_i: U_i\to \C^d$ of such charts,
covering a neighborhood of $Y$. The charts being bilipschitz,
there exists an absolute constant $D$ such that if $x\in U_i$,
\[
\abs{\log d(\phi_i(x), \phi_i(Y)) - \log d(x,Y)}\leq D.
\] Then from~\eqref{eq:UE_iteree}, replacing
$\nu$ by $\nu^{(n)}$ we may assume that the uniform expansion holds for $n=1$ and the expansion
constant $c$ is bigger than $10D$. We~will work in local charts to show that $-\log d(\cdot, Y)$ is a Margulis function.

Let $d_{U_i}$ denote the Euclidean distance in the $i$-th chart (pulled back by $\phi_i$). In~$U_i$, write
$\phi_i(x) = z = (z_1, \ldots, z_d)$ and $\phi_i(Y) = \{\Ima (z) = 0\}$.
Let $\pi(\phi_i(x)) = (\Rea(z_1), \ldots , \Rea(z_d))$ be the projection of $\phi_i(x)$ on $Y$,
so that
\begin{equation} \label{eq:dui}
d_{U_i}(x,Y) = \norm{ \phi_i(x) - \pi(\phi_i(x))} = \norm{(\Ima (z_1), \ldots , \Ima (z_d))} = \norm{\Ima(\phi_i(x))}.
\end{equation}
As before let $\Gamma(R) = \set{f\in \Gamma \sep N(f)\leq R}$, where
$N(f) = \norm{f}_{C^2} + \norm{f\inv}_{C^2}$,
and fix $f\in \Gamma(R)$. If $x$ is sufficiently close to $Y$, then so does $f(x)$, hence
$f(x)$ belongs to some chart $U_j$ and working in this chart we get
$d_{U_j}(f(x),Y) = \norm{\Ima(\phi_j(f(x)))}$.
Applying Taylor's formula to the coordinate expression $\tilde f$ of $f$, we~obtain
\begin{align*}
\notag \phi_j(f(x)) &= \tilde f(\phi_i(x)) \\& \notag = \tilde f (\pi( \phi_i(x))) + d \tilde f_{\pi( \phi_i(x))} (\phi_i(x) - \pi( \phi_i(x)))
+ O\blrpar{ \norm{ \phi_i(x) - \pi(\phi_i(x))}^2}.
\end{align*}
Now, observe that the vector $d \tilde f_{\pi( \phi_i(x))} (\phi_i(x) - \pi( \phi_i(x)))$ is purely imaginary
because $\phi_i(x) - \pi( \phi_i(x))$ is purely imaginary and $d \tilde f_{\pi( \phi_i(x))}$ is real,
since it preserves $Y$.
Thus, taking imaginary parts and using~\eqref{eq:dui} yields
\begin{equation*}
\Babs{ \frac{d_{U_j}(f(x), Y)}{d_{U_i}(x, Y) }- \anorm{df_{\pi( x)}(v_x) }}\leq C R d_{U_i}(x, Y),
\end{equation*}
where
\[
v_x = \phi_i^\varstar\Blrpar{\frac{\phi_i(x) - \pi( \phi_i(x))}{\norm{\phi_i(x) - \pi( \phi_i(x))}}},\quad
\pi(x) = \phi_i\inv \pi( \phi_i(x)),
\]
and the constant $C$ depends only on the charts.
Arguing as in~\eqref{eq:delta_1}, plugging in the bilipschitz estimate for the distance to $Y$, and increasing $C$ if necessary we get
\begin{equation*}
\Babs{\log \frac{d (f(x), Y)}{d (x, Y) }- \log \anorm{df_{\pi( x)}(v_x) }}\leq C R^2 d(x, Y) + 2D.
\end{equation*}
Finally, using the moment condition to deal
with the contribution of $\Gamma\setminus \Gamma(R)$ as in Theorem~\ref{thm:margulis}, we~obtain
\begin{equation*}
\int_\Gamma
\Babs{\log \frac{d(f(x), Y)}{d(x, Y) }- \log \anorm{df_{\pi(x)}(v_x) }} d\nu(f) \leq C R^2 d(x, Y)
+ 2D + \frac{C}{(\log R)^\eta},
\end{equation*}
and we conclude that $\log d(\cdot , Y)$ is a Margulis function by first fixing a large $R$ and then choosing
$x$ sufficiently close to $Y$, as in Theorem~\ref{thm:margulis}.
\end{proof}

\begin{example}\label{eg:margulis_submanifold}
\textit{There exists a group $\Gamma = \bra{f, g}$ of diffeomorphisms of the $3$-torus $\R^3/\Z^3$ and a finitely supported measure $\nu$ on $\Gamma$ with $\langle \supp(\nu)\rangle = \Gamma$ such that:
\begin{itemize}
\item $\Gamma$ preserves $Y:=\R^2/\Z^2\times \set{0}$;
\item there exists a neighborhood $U$ of $Y$ on which the dynamics of $(\Gamma, \nu)$ is uniformly expanding;
\item for every $x\in U$ and almost every trajectory $\omega$, $f_\omega^n (x)$ converges to $Y$.
\end{itemize}
}
\end{example}

\begin{proof}

Let $0<\gamma<1$ and $\psi$ be a diffeomorphism of the circle $\R/\Z$ such that
\begin{enumeratei}
\item\label{hyp:i} $\psi$ fixes $0$, $\psi(z)=\gamma z$ on $[-1/8,1/8]$, and $\psi([-1/4,1/4]) = [-1/4,1/4]$;
\item\label{hyp:ii} in the interval $[-1/4, 1/4]$, the only fixed points of $\psi$ are $-1/4$, $0$ and $1/4$.
\end{enumeratei}
Then there is a diffeomorphism $\varphi\colon ]-1/4,1/4[\to \R$ that conjugates $\psi\rest{ ]-1/4,1/4[}$ to $t\mto \gamma t$.

Pick $A$ and $ B$ in $ \SL(2, \Z)$ and $(c_1,c_2)$ in $\Z^2$ such that
\begin{enumeratei}\setcounter{enumi}{2}
\item\label{hyp:iii}
$A$ and $B$ generate a non-elementary subgroup;
\item\label{hyp:iv}
$\gamma$ is not an eigenvalue of $B$
and $(c_1,c_2)\neq (0,0)$.
\end{enumeratei}
Define two diffeomorphisms $g$ and $h$ of $\R^3/\Z^3$ by
\begin{equation*}
g(x,y,z) = (A(x,y)+(c_1z, c_2z), \psi(z))\; \text{ and } \; h(x,y,z) = (B(x,y), \psi(z)).
\end{equation*}
Let $\nu$ be a probability measure supported on $\set{g, h, g\inv, h\inv}$ such that
$0<\nu(g\inv)<\nu(g)$ and $0<\nu(h\inv)<\nu(h)$.

First, let us prove that there exists $\Omega_0\subset \Omega$ of full
$\nu^\N$-measure such that $f_\omega^n(p)$ converges towards $Y$ for every $p = (x,y,z)\in \R^2/\Z^2 \times{}]-1/4, 1/4[$,
and $\omega\in \Omega_0$. Indeed, writing $f_\omega^n(p) = (x_n, y_n, z_n)$, we~have:
\begin{enumerate}
\item $z_n\in{} ]-1/4, 1/4[$ because $\psi$ preserves $]-1/4, 1/4[$;
\item $\varphi(z_n ) = \gamma^{\sum_{i=1}^{n} \e_i} \varphi(z)$, where
$(\e_n)$ is a sequence of independent random variables with
$\pp(\e =1 ) = \nu(g)+\nu(h)$ and $\pp(\e =-1 ) = \nu(g\inv)+\nu(h\inv)$. Since, $\nu(g\inv)+\nu(h\inv)< \nu(g)+\nu(h)$, $\varphi(z_n)$ converges almost surely to $0$.
\end{enumerate}

Now, we~show that the dynamics of $(\Gamma, \nu)$ is uniformly expanding in the product
$\R^2/\Z^2 \times{}]-1/4, 1/4[$. Indeed, if $p\in \R^2/\Z^2 \times{}]-1/4, 1/4[$ and $\omega\in \Omega_0$, there
is $n(\omega)$ such that $f^n_\omega(p)\in \R^2/\Z^2 \times{}]-1/8, 1/8[$ for $n\geq n(\omega)$.
Now, in $\R^2/\Z^2 \times{}]-1/8, 1/8[$ the dynamics is linear, and the tangent action is generated by
\begin{equation*}\arraycolsep3.5pt
\tilde g = \begin{pmatrix}A & \begin{pmatrix} c_1 \\ c_2 \end{pmatrix} \\ 0 & \gamma
\end{pmatrix}
\quad\text{and}\quad
\tilde h = \begin{pmatrix}B& 0 \\ 0 & \gamma
\end{pmatrix}.
\end{equation*}
We claim that the linear action of $(\tilde \Gamma, \tilde \nu)$ on $\R^3$ is uniformly expanding,
where $\tilde \Gamma = \langle{\tilde g, \tilde h}\rangle$ and $\tilde \nu$ is the measure naturally
corresponding to $\nu$. Indeed, the action is uniformly expanding on $\R^2\times \set{0}$ and
if it were not uniformly expanding on $\R^3$, by~Furstenberg-Kifer~\cite{furstenberg-kifer},
there would exist a $\tilde \Gamma$-invariant line transverse to $\R^2\times \set{0}$ along which the Lyapunov exponent would be non-positive. But the hypothesis~\eqref{hyp:iv} guarantees that such a line does not exist.
From this, we~deduce that there exists $c>0$ (any constant smaller than the Lyapunov exponent
of the random product generated by $A$ and $B$ will do)
such that
for every $p \in \R^2/\Z^2 \times{}]-1/4, 1/4[$, every unit tangent vector $v$ at $p$
and almost every $\omega$, $\unsur{n} \log \norm{(f_\omega^n)_\varstar v}\geq c$ if $n$ is large enough.
Applying Lemma~\ref{lem:weak_UE2} finishes the proof.
\end{proof}

\section{An ergodic-theoretic criterion for expansion}\label{sec:criterion}

\subsection{Construction of stationary measures}
\label{subs:stationary}

Let $M$ be a compact manifold endowed with a Riemannian metric; let $T^1M$ denote its unit tangent bundle
and $\pi\colon T^1M\to M$ be the canonical projection. As in Section~\ref{sec:uniform_generalities},
if $f$ is a diffeomorphism of $M$, we~denote by $f_\varstar $ its action on $TM$.
Let $\nu$ be a probability measure on $\Diff^1(M)$
satisfying the moment condition~\eqref{eq:moment}. We~apply a classical strategy to get the
following theorem (see \eg \cite[Prop.\,3.17]{chung}, and~\cite[Lem.\,3.3]{hurtado:burnside}).

\begin{thm}\label{thm:limit_measure_top_lyapunov}
Assume that there exists an increasing sequence $(n_k)\in \N^\N$ and a
sequence of unit tangent vectors $( u_k)\in (T^1M)^\N$
such that
\begin{equation*}\label{eq:exponential_dilatation}
\lim_{k\to\infty} \frac{1}{n_k} \int \log \norm{f_\varstar u_k} d\nu^{(n_k)}(f) = \chi_0.
\end{equation*}
Then, there exists a real number $\chi\geq \chi_0$, an ergodic $\nu$-stationary probability measure~$\hat{\mu}$
on $T^1M$, and a $\nu$-almost surely invariant sub-bundle $V\subset TM$ such that the top Lyapunov exponent of the projected measure $\mu:=\pi_\varstar \hat{\mu}$ in restriction to $V$ is equal to~$\chi$.

Likewise, there exists
a real number $\chi' \leq \chi_0$ that satisfies the same property for
some pair $(\hat{\mu}', V')$ and the projection $\mu':= \pi_*\hat{\mu}'$.
\end{thm}

Note that if $\hat\mu$ is a probability measure on $T^1M$ that is
$\nu$-stationary for the tangent action, then its projection $\mu $ on $M$ is $\nu$-stationary as well; and if $\hat\mu$ is ergodic, so is $\mu$.
When $\chi >0$, one typically obtains~$V=TM$.

\begin{proof}[Proof (see~\cite{chung, hurtado:burnside})]
Consider the
sequence of measures $\hat\mu_k$ on $T^1M$ defined by
\begin{equation*}
\hat\mu_k=\frac{1}{n_k} \sum_{j=0}^{n_k-1} \nu^{(j)}\star \delta_{u_k} =\frac{1}{n_k} \sum_{j=0}^{n_k-1}\int \frac{f_\varstar u_k}{\norm{f_\varstar u_k}} d\nu^{(j)}(f),
\end{equation*}
where $\nu^{(j)}\star \delta_{u_k}$ denotes the convolution for the action of $\Diff(M)$ on the unit tangent bundle. Since $T^1M$ is compact and the
$\hat\mu_k$ are probability measures, we~can extract a subsequence (still denoted by $\hat\mu_k$ for simplicity) that converges weakly towards a
probability measure $\hat\mu_\infty$ on $T^1M$. By~construction, this measure is $\nu$-stationary.

The function $\dil(f, u):=\log\norm{f_\varstar u}$ is continuous on
$\Diff^1(M)\times T^1M$. And by our moment assumption, so is the function
$u\in T^1M \mto \int \dil(f, u) d\nu(f)$.
For $ u\in T^1M$ the chain rule gives
\begin{align*}\notag
\frac{1}{n} \int \log\norm{(f^n_\omega)_\varstar u} d\nu^{\N}(\omega) & = \frac{1}{n} \sum_{j=0}^{n-1} \int \dil\Blrpar{f_j, \frac{(f^j_\omega)_\varstar u}{\anorm{(f^j_\omega)_\varstar u}}} d\nu^{\N}(\omega) \\
& = \int_{g\in \Diff(M)} \bblrpar{\frac{1}{n} \sum_{j=0}^{n-1} \int \dil\Blrpar{g, \frac{h_\varstar u}{\norm{h_\varstar u}}} d\nu^{(j)}(h) } d\nu(g).
\end{align*}
If we apply this equation to $n=n_k$ and $u=u_k$ the term between parentheses in
the last integral is equal to
$
\int \dil\lrpar{g,u} d\hat\mu_k(u),
$
so, letting $k$ go to $+\infty$, we~conclude that
\begin{equation*}
\lim_{k\to\infty} \frac{1}{n_k} \int \log \norm{f_\varstar u_k} d\nu^{(n_k)}(f) = \chi_0=
\int_{\Diff^1(M)}\int_{T^1M} \dil(g,u) d\hat\mu_\infty(u) d\nu(g).
\end{equation*}
Thus, there exists $\chi\geq \chi_0$ (\resp $\chi\leq \chi_0$) and
an ergodic component $\hat\mu$ of $\hat\mu_\infty$
such that
\begin{equation*}
\int_{\Diff^1(M)}\int_{T^1M} \dil(g,u) d\hat\mu (u) d\nu(g) = \chi.
\end{equation*}

As observed above, $\mu = \pi_\varstar \hat \mu$ is an
ergodic $\nu$-stationary probability measure. Denote by $\hat\mu_x$ the conditional measures obtained
by disintegration of $\hat\mu$ with respect to the fibers of $\pi$, that is,
$\hat\mu = \int \hat\mu_x d\mu(x)$. For $\mu$-almost every $x$, let $V(x)$ be the linear span of
$\supp(\hat \mu(x))$. Since $\supp(\hat\mu)$ is $\nu$-almost invariant and
$f_\varstar $ acts linearly along the fibers of $TM$, we~infer that $V$ is a $\nu$-almost invariant
measurable sub-bundle. The Furstenberg formula asserts that
the top Lyapunov exponent of $\mu$ in restriction
to $V$ is equal to $\chi$. For completeness let us recall the argument:
the ergodic theorem shows that for $(\nu^\N\times \mu)$-almost every $(\omega, x)$
and $\hat\mu_x$-almost every $u\in T^1_xM$,\vspace*{-3pt}
\begin{equation}\label{eq:birkhoff}
\lim_{n\to + \infty} \frac{1}{n} \sum_{j=0}^{n-1} \dil\left( f_j, (f^j_\omega)_\varstar u \right) = \int_{\Diff^1(M)^\N} \int_{T^1M} \dil(f^1_\omega, u) \, d\hat\mu( u) d\nu^{\N}(\omega) = \chi,
\end{equation}
where as usual
$\omega=(f_0, f_1, \ldots)$, $f^1_\omega=f_0$, and $f^j_\omega= f_{j-1}\circ \cdots \circ f_0$.
On the other hand the Oseledets theorem asserts that for
$(\nu^\N\times \mu)$-almost every $(\omega, x)$, there exists a proper subspace
$W(\omega, x)\subset V(x)$ such that for $u\notin W(\omega, x)$,
$\unsur{n} \log\norm{(f^n_\omega)_\varstar u}$ converges to the top Lyapunov exponent $\chi^+(\mu, V)$
of $\mu$ in restriction to $V$. Thus by~\eqref{eq:birkhoff}, $\chi^+(\mu, V)=\chi$, and the proof is complete.
\end{proof}

\subsection{Application: Chung's criterion}

The following theorem, taken from \cite[Prop.\,3.17]{chung},
plays an important role in this paper; a variant of this result appears in~\cite{brown-fisher-hurtado}.
It is stated in \cite{chung} for $C^2$ actions on surfaces but
it holds in greater generality. The proof follows directly from the second assertion of Theorem~\ref{thm:limit_measure_top_lyapunov}.\enlargethispage{\baselineskip}%

\begin{thm}[Chung] \label{thm:criterion_chung}
Let $M$ be a compact manifold. Let $\nu$ be a probability measure on $\Diff^1(M)$
that satisfies \eqref{eq:moment}. If $\nu$ is not uniformly expanding there exists an ergodic $\nu$-stationary measure $\mu$ on $M$ and a $\mu$-measurable
subbundle $W\subset TM$ such that
\begin{enumeratea}
\item\label{thm:criterion_chunga} $0<\dim(W)\leq \dim(M)$;
\item\label{thm:criterion_chungb} $W$ is $\nu$-almost surely invariant;
\item\label{thm:criterion_chungc} in restriction to $W$, the top Lyapunov exponent of $\mu$ is non-positive.
\end{enumeratea}
Conversely, if such a pair $(\mu, W)$ exists, then $\nu$ is not uniformly expanding.
\end{thm}

When $M$ is a surface and $\nu$ is supported by the group of diffeomorphisms preserving some fixed area form
the Lyapunov exponents of any ergodic stationary measure $\mu$ satisfy $\lambda^+(\mu)+\lambda^{-}(\mu)=0$. Thus,
in Chung's theorem, either $\lambda^{-}(\mu)=\lambda^{+}(\mu)=0$ and we can take $W=TM$ or $\lambda^{-}(\mu)<0<\lambda^+(\mu)$ and $W$ coincides with the stable line field
provided by the
Oseledets theorem; thus, $\mu$ is not hyperbolic or it is hyperbolic and its stable line field is non-random.\vspace*{-.5\baselineskip}

\part{Non-elementary actions on complex surfaces}

From now on we denote by $X$ a compact complex surface, endowed with a
group~$\Gamma$ of holomorphic
diffeomorphisms. Recall from~\cite{stiffness, Simons} that if $\Gamma$ is non-elementary, then~$X$ is necessarily projective and~$\Gamma\subset \aut(X)$.

\section{Preliminaries}\label{sec:preliminaries}

In this section we briefly recall some results from \cite{invariant} (see also \cite{cantat_groupes, Cantat:Milnor}).

\subsection{Parabolic automorphisms and their dynamics (see~\cite[\S 3]{invariant})} \label{subs:halphen}

Let $h$ be a parabolic automorphism of a compact projective surface $X$ (most of this discussion is valid for a
compact Kähler surface). Then, $h$ preserves a genus $1$ fibration
$\pi_h\colon X\to B$, and every $h$-invariant holomorphic (singular) foliation -- in particular any invariant fibration -- coincides with $\pi$. Let $h_B$ denote the automorphism of $B$ such that
\begin{equation*}
\pi\circ h=h_B\circ \pi.
\end{equation*}
If $X$ is not a torus there is a positive integer $m$ such that $h^m$ preserves every fiber of $\pi$, \ie $h_B^m=\id_B$. When $h_B = \id_B$ we say that $h$ is a \emph{Halphen twist}. The set of Halphen twists in a given subgroup $\Gamma\subset \Aut(X)$ is denoted by $\Hal(\Gamma)$.

\begin{rem}\label{rem:many_parabolic}
If $\Gamma$ is non-elementary and contains a Halphen twist (\resp a parabolic automorphism) $h$, then the conjugacy class of $h$ in $\Gamma$ contains Halphen twists (\resp parabolic automorphisms) associated with infinitely many distinct invariant fibrations (see~\cite[\S 3.1]{finite_orbits}).
\end{rem}

Suppose now that $h$ is a Halphen twist.
Then, $h$ acts by translation on every smooth fiber of $\pi$ (see~\cite[\S 3]{invariant}).
To be more precise, denote by $\Crit(\pi)\subset B$ the finite set of critical values of $\pi$ and set $B^\circ= B\setminus \Crit(\pi)$.
Fix some simply connected open subset $U\subset B^\circ$, endowed with a section $\sigma$ of
$\pi$ and a continuous choice of basis for $H_1(X_w, \Z)$.
Each fiber $X_w:=\pi^{-1}(w)$, $w\in U$, is an elliptic curve with zero $\sigma(w)$, and one can find a holomorphic function $\tau$ on $U$, with values in the upper half plane, such that~$X_w$ is isomorphic to $\C/\Lat(w)$ for $\Lat(w)=\Z\oplus \Z\tau(w)$. On $X_w$, $h$ is a translation $h_w(z)=z+t(w)$, for some holomorphic function $w\in U\mto t(w)\in \C/\Lat(w)$. Moreover, Lemma~\ref{lem:halphen_flat}(4) says that $h$ behaves like a ``complex Dehn twist'', with a shearing property in the direction which is transversal to the fibers; thus shearing (or~twisting) occurs along $X_w$ whenever $t$ and $\tau$ are ``transverse'' at $w$
(see Section~\ref{subs:finitary} for more details on the non-twisting locus).

The points $w$ for which $h_w$ is periodic are characterized by the relation $t(w)\in \Q\oplus \Q\tau(w)$. If
\begin{equation*}
t(w)-(\alpha+\beta\tau(w))\in \R \cdot (p+q\tau(w))
\end{equation*}
for some $(\alpha,\beta)\in \Q^2$ and $(p,q)\in \Z^2$, the closure of $\Z t(w)$ in $\C/\Lat(w)$ is an abelian Lie group of dimension $1$,
isomorphic to $\Z / k\Z \times \R /\Z$ for some $k>0$; then, the closure of each orbit of $h_w$ is a union of $k$ circles. This occurs along a countable union of analytic curves $\mathrm{R}^{\alpha, \beta}_{p, q}\subset U$. Otherwise, the orbits of $h_w$ are dense in $X_w$, and the unique $h_w$ invariant probability
measure is the Haar measure on $X_w$.

The following lemma summarizes this discussion.

\begin{lem}\label{lem:halphen_flat} Let $h$ be a Halphen twist with invariant fibration
$\pi:X\to B$. Then,
\begin{enumerate}
\item\label{lem:halphen_flat1} $h$ acts by translation on each fiber $X_w=\pi^{-1}(w)$, $w \in B^\circ$;
\item\label{lem:halphen_flat2} for $w$ in a dense countable subset of $B^\circ$, the orbits of $h_w$ are finite;
\item\label{lem:halphen_flat3} there is a dense, countable union of analytic curves $\mathrm R_j$ in $B^\circ$,
such that:
\begin{enumerate}
\item\label{lem:halphen_flat3a} for $w\notin\bigcup_j \mathrm R_j$, the action of $h$ in the fiber $X_w$ is a totally irrational
translation (it is uniquely ergodic, and its orbits are dense in $X_w$);
\item\label{lem:halphen_flat3b} for $w\in \bigcup_j \mathrm R_j$ the orbits of $h_w$ are either finite or dense in a finite union of circles;
\end{enumerate}
\item\label{lem:halphen_flat4} there is a finite subset $\mathrm{NT}_h$ such that for $x\notin \pi\inv\lrpar{
\mathrm{NT}_h}$
\[
\lim_{n\to \pm \infty} \norm{D_xh^n}\to +\infty
\]
locally uniformly in $x$; more precisely for every
$v\in T_xX \setminus T_x X_{\pi(x)}$, $\norm{D_x h^{n}(v)}$ grows linearly while
$\frac{1}{n}\pi_\varstar (D_xh^{n}(v))$ converges to $0$.
\end{enumerate}
If moreover $h$ preserves a totally real $2$-dimensional real analytic subset $Y\subset X$, then:
\begin{enumerate}\setcounter{enumi}{4}
\item\label{lem:halphen_flat5} the generic fibers of $\pi\rest{Y}$ are union of circles,
there exists an integer $m$ such that $h^m$ preserves each of these circles, and $h^m$ is uniquely ergodic along each of these circles,
except for countably many fibers.
\end{enumerate}
\end{lem}

Property~\eqref{lem:halphen_flat4} is the above mentioned twisting property of $h$. Property~\eqref{lem:halphen_flat5} occurs, for instance, when $X$ and $h$ are defined over $\R$
and $Y=X(\R)$ is the real part of $X$. There are also examples of
subgroups $\Gamma\subset \Aut(X)$ preserving a totally real surface $Y\subset X$ which is not the real part of $X$ for any real structure, see \cite[\S 9]{invariant}.

\subsection{Classification of invariant measures}

Recall from Example~\ref{eg:K3_intro} that if $X$ is a torus, a K3 surface, or an Enriques surface it admits a canonical $\Aut(X)$-invariant volume form $\vol_X$. The associated probability measure will also be denoted by $\vol_X$. Such an area form exists also on any totally real surface, by~virtue of the following lemma.

\begin{lem}[{see \cite[Rem.\,2.3]{invariant}}]\label{lem:volume_Y}
Let $X$ be an Abelian surface, or a K3 surface, or an Enriques surface with universal cover $\tilde{X}$.
Let $Y\subset X$ be a totally real surface of class~$C^1$, and $\Aut(X;Y)$ be the
subgroup of $\Aut(X)$ preserving $Y$.
If $Y$ is totally real, the canonical holomorphic 2-form
$\Omega_X$ (\resp $\Omega_{\tilde{X}}$) induces a smooth $\Aut(X;Y)$-invariant probability
measure $\vol_Y$ on $Y$.
\end{lem}

\begin{thm}[{see~\cite[Th.\,A]{invariant}}]\label{thm:classification_invariant}
Let $X$ be a projective surface. Let $\Gamma$ be a non-elementary subgroup of $\Aut(X)$ containing a parabolic element. Let $\mu$ be a $\Gamma$-invariant ergodic probability measure on $X$.
Then, $\mu$ satisfies exactly one of the following properties.

\begin{enumeratea}
\item\label{thm:classification_invarianta} $\mu$ is the average on a finite orbit of $\Gamma$;
\item\label{thm:classification_invariantb} $\mu$ is non-atomic and supported on a $\Gamma$-invariant algebraic curve $D\subset X$;
\item\label{thm:classification_invariantc} there is a $\Gamma$-invariant proper algebraic subset $Z$ of $X$, and a $\Gamma$-invariant,
totally real analytic surface $Y$ of
$X\setminus Z$ such that {\em {(1)}} $\mu(\overline{Y}) = 1$ and $\mu(Z)=0$; {\em {(2)}} $Y$ has finitely many irreducible components;
{\em {(3)}} the singular locus of $Y$ is locally finite in $X\setminus Z$;
{\em {(4)}} $\mu$
is absolutely continuous with respect to the Lebesgue measure on $Y$; and~{\em {(5)}} its density (with respect to any real analytic area form on the regular part of~$Y$)
is real analytic;
\item\label{thm:classification_invariantd} there is a $\Gamma$-invariant proper algebraic subset $Z$ of $X$ such that {\em {(1)}} $\mu(Z)=0$, {\em {(2)}} the support of $\mu$ is equal to $X$; {\em {(3)}} $\mu$ is absolutely continuous with respect to the Lebesgue measure on $X$; and {\em {(4)}} the density of $\mu$ with respect to any real analytic
volume form on $X$ is real analytic on $X\setminus Z$.
\end{enumeratea}
If $X$ is not a rational surface, then in case~\eqref{thm:classification_invariantc} (\resp \eqref{thm:classification_invariantd}) we can further conclude that the invariant measure is proportional to $\vol_Y$ (\resp $\vol_X$).
\end{thm}

\subsection{Invariant curves}\label{subs:invariant_curves}
By \cite[Lem.\,2.12]{stiffness}, any action of a non-elementary group~$\Gamma$ on a
projective surface $X$ admits a maximal invariant curve $\mathrm D_\Gamma$,
which can be easily detected from the action of $\Gamma$ on $H^2(X, \Z)$
since it corresponds to an invariant class. Bounds on the degrees of such invariant curves in terms of the action are given in~\hbox{\cite[\S 3]{finite_orbits}}. If in addition $\Gamma$
contains a parabolic element, $\mathrm D_\Gamma$ is the set of common components of the singular fibers of all elliptic fibrations associated to
parabolic elements in $\Gamma$ (see~\cite[\S 4.1]{invariant}).

\section{Hyperbolicity of invariant measures}\label{sec:hyperbolicity}

Here, $X$ is a compact Kähler surface. We~fix a Kähler form $\kappa_0$ on $X$; norms of tangent vectors and differentials will be computed with respect to it.

\Subsection{Ledrappier's invariance principle and invariant measures on $\P T\cX$}\label{subs:ledrappier_invariance_principle}
In this paragraph we collect some preliminary results for the proof of Theorems \ref{thm:hyperbolic} and \ref{thm:hyperbolic_kummer}.
The reader should also consult~\cite{barrientos-malicet} and~\cite{obata-poletti} for comparison;
\cite{obata-poletti} relies on the ``pinching and twisting'' formalism of Avila and Viana
(see \cite{viana} for an introduction.\footnote{Beware that the word ``twisting'' has a different meaning there.}
Most of this discussion is valid for a random holomorphic dynamical system on an arbitrary complex surface (not necessarily compact)
satisfying \eqref{eq:moment}.

We denote by $ \P TX$ the projectivized tangent bundle of $X$; if $f$ is an automorphism of $X$, we~denote by $\P(Df)$ the induced action on $\P TX$.

Let $\nu$ be a probability measure on $\Aut(X)$ that satisfies the moment condition~\eqref{eq:moment}. We~endow $\Omega:=\Aut(X)^\N$ (\resp $\Sigma:= \Aut(X)^\Z$)
with the probability measure $\nu^\N$ (\resp $\nu^\Z$), and set $\cX_+=\Omega \times X$ (\resp $\cX=\Omega \times X$); $\sigma$ will denote the shift (on~$\Omega$ or~$\Sigma$).
For $\omega=(f_i)_{i\geq 0}\in \Omega$, we~keep the notation $f_\omega^{n}$ from
Section \ref{subs:notation}. Then, we~define $F_+\colon \cX_+\to \cX_+$ by
$F_+(\omega,x)=(\sigma(\omega), f^1_\omega(x))$; $F\colon \cX\to \cX$ is defined by the same formula. For further standard notations, we~refer to \cite[\S 7]{stiffness}.

Let $\mu$ be an ergodic $\nu$-stationary measure on $X$. We~introduce the projectivized tangent bundles $\P T \cX_+ = \Omega \times \P TX$ and $\P T\cX = \Sigma\times \P TX$. The bundles $TX$ and $\P TX$ admit measurable trivializations over a set of full measure.
Consider any probability measure $\hat\mu$ on $\P TX$ that is stationary under the random dynamical system induced by $(X,\nu)$ on $\P TX$
and whose projection on $X$ coincides with $\mu$, \ie $\pi_\varstar \hat\mu= \mu $ where $\pi\colon \P TX\to X$ is the natural projection.
Such measures always exist: indeed,
the set of probability measures on $\P TX$ projecting to $\mu$ is compact and convex, and it is non-empty since it contains
the measures $\int \delta_{[v(x)]} d\mu(x)$ for any measurable section
$x\mto [v(x)]$ of $\P TX$. Thus, the operator $\int \P (Df) \,d\nu(f)$ has a fixed point on that set. The stationarity of $\hat\mu$ is equivalent to the invariance of $\nu^\N\times \hat\mu$
under the transformation ${\widehat{F}_+}\colon \Omega\times \P TX\to \Omega\times \P TX$ defined by
\begin{equation*}
{\widehat{F}_+}(\omega, x, [v])= (\sigma(\omega), f^1_\omega(x), \P(D_xf^1_\omega)[v])
\end{equation*}
for any non-zero tangent vector $v\in T_xX$. We~denote by $\hat\mu_x$ the family of probability measures
on the fibers $\P T_x X$ of $\pi$
given by the disintegration of $\hat\mu$ with respect to $\pi$.
The conditional measures of $\nu^\N\times\hat\mu$ with respect to the projection
$\P T\cX_+\to X$ are given by $\hat\mu_{\omega,x} = \nu^\N\times \hat\mu_x$.

\begin{rem}
Even when $\mu$ is $\Gamma_\nu$-invariant, this construction only provides a stationary measure on
$\P TX$. This is exactly what happens for non-elementary subgroups with a parabolic automorphism:
indeed, we~show in Section \ref{subs:proof_hyperbolic} that projectively invariant measures do not exist in this case.
\end{rem}

The tangent action of our random dynamical system induces a stationary product of
matrices in $\GL(2, \C)$. To see this, fix a measurable trivialization $P\colon TX\to X\times \C^2$, given
by linear isomorphisms $P_x\colon T_xX\to \C^2$. It conjugates the action of $DF_+$ to that of a linear
cocycle $A:\cX_+ \times \C^2\to \cX_+ \times \C^2$ over $(\cX_+, F_+, \nu^\N\times \mu)$. In~this context,
Ledrappier establishes in \cite{ledrappier_stationary}
the following ``invariance principle''.

\begin{thm}\label{thm:ledrappier_invariance_principle}
If $\lambda^{-}(\mu) = \lambda^{+}(\mu)$, then for any stationary measure $\hat \mu$ on $\P TX$ projecting to $\mu$, we~have $\P(D_xf)_\varstar \hat\mu_x = \hat\mu_{f(x)}$
for $\mu$-almost every $x$ and $\nu$-almost every~$f$.
\end{thm}

The second ingredient in the proof of Theorem~\ref{thm:hyperbolic} is
a description of such projectively invariant measures; this is where we follow \cite{barrientos-malicet}.
To explain this result a bit of notation is required. Let $V$ and $W$ be hermitian vector spaces of dimension $2$;
we fix two isometric isomorphisms $\iota_V\colon V\to \C^2$ and $\iota_W\colon W\to \C^2$ to the standard
hermitian space $\C^2$, and we endow the projective lines $\P(V)$ and $\P(W)$ with their respective Fubini-Study metrics. If $g\colon V\to W$
is a linear isomorphism, we~set
\begin{equation*}
\llbracket g \rrbracket= \norm{ \P (g) }_{C^1},
\end{equation*}
where $\P(g) \colon \P(V)\to \P(W)$
is the projective linear map induced by $g$ and $\norm{\cdot}_{C^1}$
is the maximum of the norms of $D_z\P(g) \colon T_z\P(V)\to T_{\P (g)(z)}\P(W)$ with respect to the Fubini-Study metrics.
If $\iota_W\circ g\circ \iota_V^{-1} = k_1ak_2$ is the KAK decomposition of
$\iota_W\circ g\circ \iota_V^{-1}$ in $\GL(2, \C)$, we~get
\begin{equation*}
\llbracket g \rrbracket = \frac{\norm{a}^2}{\abs{\det(a)}} = \frac{\norm{\iota_W\circ g\circ \iota_V^{-1}}^2}{\abs{\det(\iota_W\circ g\circ \iota_V^{-1})}},
\end{equation*}
where
$\norm{\cdot}$ is the matrix norm in $\GL_2(\C)$ associated to the Hermitian norm of $\C^2$. In~particular,
\begin{itemize}
\item[(a)] $\llbracket g \rrbracket=1$ if and only if $\P (g)$ is an isometry from $\P(V)$ to $\P(W)$;
\item[(b)] for a sequence $(g_n)$ of linear maps $V\to W$,
$\llbracket g_n\rrbracket$ tends to $+\infty$ with $n$ if and only if $\P(\iota_W\circ g\circ \iota_V^{-1})$
diverges to infinity in $\PGL_2(\C)$.
\end{itemize}
If $f$ is an automorphism of $X$ and $x$ is a point of $X$, then $\kappa_0$ endows $T_xX$ and $T_{f(x)}X$ with
hermitian structures, and we can apply this discussion to $D_xf\colon T_xX\to T_{f(x)}X$. We~are now ready to state the classification of projectively invariant measures.

\begin{thm}\label{thm:classification_proj_invariant}
Let $(X, \nu)$ be a random dynamical system on a complex
surface and let $\mu$ be an ergodic stationary measure. Let
$\hat\mu$ be a stationary measure on $\P TX$ such that $\pi_\varstar {\hat{\mu}}=\mu$ and $\P(D_xf)_\varstar \hat\mu_x = \hat\mu_{f(x)}$
for $\mu$-almost every $x$ and $\nu$-almost every $f$. Then, exactly one of the following two properties is satisfied:
\begin{enumerate}
\item\label{thm:classification_proj_invariant1} For $(\nu^\N\times \mu)$-almost every
$(\omega, x)$, the sequence $\llbracket D_xf^n_\omega\rrbracket$
is unbounded and then:
\begin{enumerate}
\item\label{thm:classification_proj_invariant1a} either there exists a measurable $\Gamma_\nu$-invariant family of lines
$E(x)\subset T_xX$ such that $\hat\mu_x=\delta_{[E(x)]}$ for $\mu$-almost every $x$;
\item\label{thm:classification_proj_invariant1b} or there exists a measurable $\Gamma_\nu$-invariant family of pairs of lines
$E_1(x), E_2(x) \subset T_xX$ and positive numbers $ \lambda_1, \lambda_2$ with
$\lambda_1+\lambda_2 =1$ such that $\hat\mu_x=\lambda_1\delta_{[E_1(x)]} + \lambda_2\delta_{[E_2(x)]}$ for $\mu$-almost every $x$.
\end{enumerate}
\item\label{thm:classification_proj_invariant2} The projectivized tangent action of $\Gamma_\nu$ is reducible to a compact group, that is
there exists a measurable trivialization
of the tangent bundle $(P_x: T_xX \to \C^2)_{x\in X}$, such that for almost
every $f\in \Gamma_\nu$ and every
$x$, $\P\lrpar{P_{f(x)}\circ D_xf\circ P_x\inv}$ belongs to the unitary group $\mathsf{PU}_2(\C)$.
\end{enumerate}
\end{thm}

In Assertion \eqref{thm:classification_proj_invariant1b}, the pair is not ordered: there is no natural distinction of~$E_1$ and~$E_2$,
the elements of $\Gamma_\nu$ may a priori permute these lines.
The proof can be obtained by adapting the arguments of \cite{barrientos-malicet} to the complex case;
full details are given in \cite[\S 7.4]{arxiv-version-of-this-paper}. We~provide a shorter
proof, suggested by
one of the referees, that relies on results of Furstenberg and Zimmer.
Yet another approach, suggested by
another referee, would be to view $\P^1(\C)$ as the boundary of $\mathbb H^3$ and use the notion
of the conformal barycenter of Douady-Earle~\cite{douady-earle}.

\begin{proof}
Consider a probability space $(Y, \mathcal{A}, \m)$ together with an ergodic measure preserving transformation $T\colon Y\to Y$ and a measurable cocycle $\eta\colon (Y, \mathcal A)\to (\GL_2(\C), \mathcal{B}(\GL_2(\C))$. Suppose we are given a measurable map
$\hat\m\colon Y\to \mathrm{Prob}(\P^1(\C))$ from $Y$ to the space of probability measures on $\P^1(\C)$ (equipped with its Borel $\sigma$-algebra)
and that this map is $\eta$-equivariant, \ie
\begin{equation*}
\hat\m_{Ty}=\eta(y)_*\hat\m_y
\end{equation*}
for almost every $y$ in $Y$. The ergodicity of $\m$ and Theorem~3.2.6 of~\cite{Zimmer:book} imply
that, on a subset of full measure in $Y$, $\hat\m$ takes values in a unique $\GL_2(\C)$-orbit in $\mathrm{Prob}(\P^1(\C))$. Let $\lambda\in \mathrm{Prob}(\P^1(\C))$ be a point in this orbit. The equivariance of $\hat\m\colon Y\to \GL_2(\C)_\varstar\lambda$ means that $\hat\m$ is cohomologous to a cocycle taking values in the stabilizer of $\lambda$ in $\GL_2(\C)$, that is,
there exists a measurable map $Y\ni y\mto A_y\in \GL_2(\C)$ such that
$A_{Ty}\inv \eta_y A_y \in \Stab(\lambda)$. Indeed $\hat m_y = (A_y)_\varstar \lambda$ for some $A_y$.
Then, according to Lemma~3.2.1 of~\cite{Zimmer:book} and its Corollary~3.2.2, there are only three possibilities. Either~$\lambda$ is a Dirac mass, or $\lambda$ is an average of two Dirac masses, or the stabilizer of $\lambda$ is compact.

In our situation, we~take $Y$ to be $\cX^+ = \Omega\times X$, $T$ is $F_+$,
$\m$ is $\nu^\N \otimes \mu$, and $\hat \m$ is the family of disintegrations
$(\omega, x)\mto \hat \mu_x$, which is a Borel map since it is the disintegration of a Borel measure on $X\times \P^1$ relative to the projection $X\times \P^1\to X$. Applying the above results, we~get the desired conclusion, except for one point, which is the fact that in the conclusion of the theorem,
the trivialization $P$ in Assertion~\eqref{thm:classification_proj_invariant2} (\resp the lines in Assertion~\eqref{thm:classification_proj_invariant1}) depends only on $x$, and not on $(\omega,x)$, Theorem~\ref{thm:classification_proj_invariant}.
a fact that will be used later.
Since the lines in Assertion~\eqref{thm:classification_proj_invariant1} are directly determined by $\hat \mu_x$ it is obvious that they
depend uniquely on $x$.

Assume that we are in the situation where $\Stab(\lambda)$ is a compact subgroup.
To show that $P$ can be chosen to depend only on $x$ we argue as follows. First, observe that
$\Stab(\lambda) \subset P_0\inv \mathsf{U}_2(\C)P_0 $ for some $P_0\in \GL_2(\C)$, so with notation as above we obtain that $D_xf$ is of the form
$A_{F_+(\omega, x)} P_0 U P_0\inv A_{(\omega, x)}\inv$, for some $U\in \mathsf{U}_2$. We~need to show that $A_{(\omega, x)}$ can be chosen to depend on $x$ only.
Let
\begin{equation*}
E = \bset{(x, A)\in X\times \GL_2(\C)\sep  A_\varstar \lambda = \hat \mu_x}.
\end{equation*} Since
$x\mto \mu_x$ is Borel, $E$ is a Borel subset
of $X\times \GL_2(\C)$ whose fibers in $\GL_2(\C)$ are empty or compact.\footnote{On the null set where the fibers are empty, we~may replace the fiber by some fixed matrix, say $A=\id$.}
From the Borel selection theorem (see~\cite[Th.\,6.9.6]{bogachev_book} or \hbox{\cite[App.\,A]{Zimmer:book}}), there exists a Borel map $x\mto A'_x$ such that for every $x\in X$, $(x, A'_x)$ belongs to $E$; replacing
$A_{(\omega, x)}$ by $A'_x$ concludes the proof.
\end{proof}

\subsection{Proof of Theorem \ref{thm:hyperbolic}}\label{subs:proof_hyperbolic}
By Theorem~\ref{thm:classification_invariant},
$\mu$ is either equivalent to the Lebesgue measure on $X$, or to the $2$-dimensional Lebesgue measure on some components of an invariant totally real surface $Y\subset X$.

\subsubsection{Proof of the hyperbolicity of $\mu$}\label{subsubsec:721}
Let us assume, by~way of contradiction, that $\mu$ is not hyperbolic. Hence its Lyapunov exponents vanish,
and by Theorem~\ref{thm:ledrappier_invariance_principle} and Theorem \ref{thm:classification_proj_invariant}, there
is a measurable set $X'\subset X$
with $\mu(X') = 1$ such that one of the following properties is satisfied along $X'$: \label{alternative_abc}
\begin{enumeratea}
\item\label{case:a} there is a measurable $\Gamma_\nu$-invariant line field $E(x)$;
\item\label{case:b} there exists a measurable $\Gamma_\nu$-invariant splitting $E(x)\oplus E'(x)=T_xX$ of the tangent bundle;
here, the invariance should be taken in the following weak sense: an element $f$ of $\Gamma_\nu$ maps $E(x)$ to
$E(f(x))$ or $E'(f(x))$;
\item\label{case:c} there exists a measurable trivialization $P_x\colon T_xX\to \C^2$ such that in the corresponding coordinates
the projectivized
differential $\P(Df_x)$ takes its values in $\mathsf {PU}_2(\C)$ for all $f\in \Gamma_\nu$ and $\mu$-almost all $x\in X'$.
\end{enumeratea}

Fix a small $\e>0$. By~Lusin's
theorem, there is a compact set $K_\e $ with $\mu(K_\e)>1-\e$
such that the data $x\mto E(x)$, or $x\mto (E(x), E'(x))$ or $x\mto P_x$ in the respective
cases \eqref{case:a}, \eqref{case:b}, and \eqref{case:c}
are continuous on $K_\e$. In~particular, in case \eqref{case:c}, the norms of $P_x$ and $P_x^{-1}$ are bounded by some uniform constant $C(\e)$
on $K_\e$; hence, if $g\in \Gamma_\nu$ and $(x,g(x))$ belongs to $K_\e^2$, $\llbracket Dg_x \rrbracket$ is bounded by $C(\e)^2$.

Fix a pair of parabolic elements $g$ and $h\in \Gamma_\nu$ with distinct invariant fibrations $\pi_g\colon X\to B_g$ and $\pi_h\colon X\to B_h$
respectively (see Remark~\ref{rem:many_parabolic}). These two fibrations are tangent along some curve in $X$, denoted by $\Tang(\pi_g, \pi_h)$.

\begin{itemize}
\item
In a first stage we assume that $X$ is not a torus. According to Section~\ref{subs:halphen}, there is an integer $N>0$ such that $g^N$ and $h^N$
preserve every fiber of their respective invariant fibrations. From now on, we~replace $g$ by $g^N$ and $h$ by $h^N$. Also, $\Gamma_\nu$ is discrete, so we may also assume $\nu(g)\nu(h)>0$ (see footnote~\ref{page:discrete}).

First assume that $\mu$ is absolutely continuous with respect to the Lebesgue measure on $X$, with a positive real analytic density
on the complement of some invariant, proper, Zariski closed subset. We~apply Lemma~\ref{lem:halphen_flat} to $h$ and remark that
$(\pi_h)_\varstar \mu$ can not charge the union of the curves $\mathrm R_j$. Then, we~disintegrate $\mu$ with respect to $\pi_h$ to obtain
conditional measures $\mu_b$, for $b\in B_h$; since $\pi_h$ is holomorphic, the measures~$\mu_b$ are
absolutely continuous with respect to the Haar measure on almost every fiber $\pi_h^{-1}(b)$. By~Lemma~\ref{lem:halphen_flat},
there exists a fiber $\pi_h\inv(b)$ such that (1) the Haar measure of
$K_\varepsilon \cap \pi_h\inv(b)$ is positive,
(2) $b \notin \mathrm{NT}_h$ and (3)
the dynamics of $h$ in $\pi_h\inv(b)$ is uniquely ergodic. These properties hold for $b=\pi_h(z)$, for $\mu$-almost all $z$ in $K_\e$.
Then we can pick $x\in \pi_h\inv(b)$ such that $(h^{k}(x))_{k\geq 0}$ visits $K_\e$ infinitely many times.\footnote{Note that we use the invariance of $\mu$ here, not mere $\nu$-stationarity.} The fifth assertion of Lemma~\ref{lem:halphen_flat}
rules out case \eqref{case:c} because the twisting property implies that the projectivized derivative $\llbracket Dh^n_x \rrbracket$ tends to infinity, while it should be bounded
by $C(\e)^2$ when $h^n(x)\in K_\e$.
Case \eqref{case:b} is also excluded: under the action of $h^n$, tangent vectors projectively
converge to the tangent space of the fibers, so the only possible invariant subspace of dimension $1$ is $\ker(D\pi_h)$.
Thus we are in case \eqref{case:a} and moreover $E(x) = \ker D_x\pi_h$ for $\mu$-almost every $x$. But then,
using $g$ instead of $h$ and the fact that $\mu$ does not charge the curve $\Tang(\pi_g, \pi_h)$, we~get a contradiction.
This shows that the last alternative \eqref{case:a} does not hold either, and this contradiction proves that~$\mu$ is hyperbolic.

If $\mu$ is supported by a $2$-dimensional real analytic subset $Y\subset X$, the same proof applies,
except that we disintegrate
$\mu$ along the singular foliation of $Y$ by circles induced by $\pi_h$ and we use the fact that a generic leaf is a
circle along which
$h$ is uniquely ergodic (see Lemma~\ref{lem:halphen_flat}\eqref{lem:halphen_flat4}).

\item
If $X$ is a torus its tangent bundle is trivial and the differential of an automorphism is
constant. In~an appropriate basis, the differential of a Halphen twist $h$ is of the form
\begin{equation*}\arraycolsep3.5pt
\begin{pmatrix} 1 & \alpha \\ 0 & 1 \end{pmatrix} \quad\text{with } \alpha\neq 0.
\end{equation*}
Thus we are in case \eqref{case:a} with $E(x) = \ker D_x\pi_h$ for $\mu$-almost every $x$. Using another
twist $g$ transverse to $h$ we get a contradiction as before.
\end{itemize}

\subsubsection{Proof of the positivity of the fiber entropy} \label{subs:proof_positivity_entropy}
This follows from classical arguments. Since $\mu$ is invariant the measure $\m:=\nu^\Z\times \mu$ on $\cX$ is
$F$-invariant. In~both cases $\mu\ll \vol_X$ and $\mu\ll \vol_Y$, respectively.
The absolute continuity of the foliation by local Pesin unstable manifolds implies that
the unstable conditionals of $\m$ cannot be atomic, see
\eg \cite[Th.\,B, \S 3]{ledrappier-young3}.
Since the unstable conditionals of
a zero entropy stationary measure are automatically atomic (see \cite[Cor.\,7.14]{stiffness}), we~conclude that
$\mu$ has positive fiber entropy.

This concludes the proof of Theorem~\ref{thm:hyperbolic}\qed

\subsection{A variant of Theorem~\ref{thm:hyperbolic}}
Let us first recall the definition of classical Kummer examples
(see~\cite[\S 4]{finite_orbits} for a thorough treatment)).
Let $A = \C^2/\Lambda$ be a complex torus and let $\eta$ be the involution given by $\eta(z_1, z_2) = (-z_1, -z_2)$; it has $16$
fixed points. Then $A/\langle\eta\rangle$ is a
surface with $16$ singular points, and resolving these singularities (each of them requires a single blow-up)
yields a Kummer surface $X$.
Let $f_A$ be a loxodromic automorphism of $A$ which is induced by a linear transformation of $\C^2$ preserving $\Lambda$;
then $f_A$ commutes to $\eta$ and goes down to an automorphism $f$ of~$X$; such automorphisms will be referred to as
loxodromic, classical, Kummer examples. They preserve the canonical volume $\vol_X$. The Kummer surface $X$
also supports automorphisms which do not come from automorphisms of $A$ (see~\cite{Keum-Kondo} and~\cite{Dolgachev-Keum} for instance).

In the following statement we do not assume that $\Gamma_\nu$ contains a parabolic element.

\begin{thm} \label{thm:hyperbolic_kummer}
Let $(X, \nu)$ be a non-elementary random dynamical system on a Kummer K3 surface satisfying
\eqref{eq:moment} and such that $\Gamma_\nu$ contains a loxodromic classical Kummer example. Then
any ergodic $\Gamma_\nu$-invariant
measure giving no mass to proper Zariski closed subsets of $X$ is hyperbolic.
\end{thm}

\begin{proof}
The proof is similar to that of Theorem \ref{thm:hyperbolic} so we only sketch it. Assume by contradiction that $\mu$ is
not hyperbolic. Since $X$ is a K3 surface, the invariance of the volume
shows that the sum of the Lyapunov exponents of $\mu$ vanishes (see \cite[\S 7.3]{stiffness}), thus both are
equal to $0$, and one of the alternatives of Theorem \ref{thm:classification_proj_invariant} holds, referred to as \eqref{case:a}, \eqref{case:b}, \eqref{case:c} as in in the beginning of Section~\ref{subsubsec:721}.

By assumption, $\Gamma_\nu$ contains a loxodromic, classical Kummer example $f$ associated to a linear automorphism $f_A$ of a torus $A$. This automorphism $f$ is uniformly
hyperbolic in a dense Zariski open subset $U$, which is thus of full $\mu$-measure: its complement is given by the sixteen rational curves
coming from the resolution of the singularities of $A/\eta$. We~denote by $x\mto E^u_f(x)\oplus E^s_f(x)$
the associated splitting of $TX\rest{U}$. The line field $E^u_f$ (\resp $E^s_f$) is everywhere tangent
to an $f$-invariant (singular) holomorphic foliation $\mathcal F^u$ (\resp $\mathcal F^s$) coming from the $f_A$ invariant linear unstable (\resp stable) foliation on~$A$.
Since $f$ is uniformly expanding/contracting on $E^{u/s}_f$, alternative \eqref{case:c} is not possible.

If alternative \eqref{case:a} holds, then $E(x)$ being $f$-invariant on a set of full
measure, it~must coincide with $E^u_f$ or $E^s_f$, say with $E^u_f$. By~continuity any $g\in \Gamma_\nu$ preserves
$E^u_f$ pointwise on $\supp(\mu)$.
Since in addition $\mu$ is Zariski diffuse, $g$
preserves $E^u_f$ everywhere on $X$, so it
preserves the unstable holomorphic foliation $\mathcal F^u$.
From this, we~shall contradict
the fact that $\Gamma_\nu$ is non-elementary. We~use a dynamical argument, based on basic constructions which are surveyed in~\cite{Cantat:Milnor}; one can also derive a contradiction from~\cite{Cantat-Favre}.

Every leaf of $\mathcal F^u$, except a finite number of them, is parametrized by an injective entire holomorphic curve $\varphi\colon \C\to X$, the image of which is Zariski dense.
Fix a Kähler form $\kappa$ on $X$ and consider the positive currents defined by
\begin{equation*}
\alpha\mto \bblrpar{\int_0^R\int_{\disk(0;t)} \varphi^*\kappa \frac{dt}{t}}^{-1} \int_0^R\int_{\disk(0;t)} \varphi^*\alpha \frac{dt}{t}
\end{equation*}
for any smooth $(1,1)$-form $\alpha$. As $R$ goes to $+\infty$, it is known that this sequence of currents converges to a closed positive current $T^+_f$
that does not depend on the parametrization $\varphi$ of the leaf, nor on the leaf itself (provided it is Zariski dense).
This current is uniquely determined by $\mathcal F^u$ and the normalization $\langle T^+_f\vert \kappa \rangle= 1$. Dynamically, it is
the unique closed positive current $T^+_f$ that satisfies $\langle T^+_f\vert \kappa \rangle= 1$ and $f^*T^+_f=\lambda(f)T^+_f$ for some $\lambda(f)>1$.
Its cohomology class $[T^+_f]$ is a non-zero element of $H^{1,1}(X;\R)$ of self-intersection $0$.

Now, pick any element $g$ of $\Gamma_\nu$. Since $g$ preserves $\mathcal F^u$, it permutes its leaves and preserves the ray generated by $[T^+_f]$.
Thus, $\Gamma_\nu$ preserves an isotropic line for the intersection form in $H^{1,1}(X;\R)$, and this contradicts the non-elementarity assumption (see~\cite[\S 2.3]{stiffness}).

Finally, if alternative \eqref{case:b} holds, any $g\in \Gamma_\nu$ preserves
$ \{E^u_f(x), E^s_f(x)\}$ on a set of full measure so,
since $\mu$ is Zariski diffuse,
it must either preserve or
swap these directions. Passing to an index $2$ subgroup both directions are preserved, and we again contradict the non-elementary assumption, as in case \eqref{case:a}.
\end{proof}

\section{Characterization of uniform expansion}\label{sec:characterization}
In this section we build on the previous results, in conjunction with the measure rigidity
results from our previous work~\cite{stiffness},
to find sufficient conditions for as well as obstructions to
uniform expansion for a non-elementary action on a compact complex surface.

\Subsection{Proof of Theorem~\ref{thm:criterion_uniform_expansion_parabolic} and related results}

\Subsubsection{Applying Chung's criterion}

\begin{defi}\label{defi:non_expanding_measure}
Let $\nu$ be a probability measure on $\Aut(X)$. A
$\nu$-stationary measure $\mu$ on $X$ is said to be
\emph{non-expanding} if
every ergodic component $\mu'$ of $\mu$ satisfies:
\begin{enumeratei}
\item\label{defi:non_expanding_measurei} either both Lyapunov exponents of $\mu'$ are non-positive,
\item\label{defi:non_expanding_measureii} or $\mu'$ is hyperbolic and its field of Oseledets stable directions is non-random.
\end{enumeratei}
\end{defi}

Recall that for a hyperbolic stationary measure $\mu$, for $\nu^\N\times \mu$-a.e. $(\omega, x)$, the stable Oseledets subspace $E^s(\omega, x)\subset T_xX$ is defined by
$v\in E^s(\omega, x)$ if and only if $\limsup_{n\to\infty} \unsur{n}\log \norm{(f_\omega^n)_\varstar v}<0$. The field of Oseledets subspaces is said to be non-random if for $\mu$-a.e. $x$,
$\omega\mto E^s(\omega, x)$ is constant mod. $0$.

Theorem~\ref{thm:criterion_chung} asserts that the existence
of non-expanding $\nu$-stationary measures is the obstruction to uniform expansion of $\nu$:

\begin{cor}[of Theorem~\ref{thm:criterion_chung}] \label{cor:criterion_chung}
Let $X$ be a compact complex surface and $\nu$ be a probability measure on the group $\Aut(X)$, satisfying the moment condition \eqref{eq:moment}.
Then~$\nu$ is uniformly expanding if
and only if non-expanding $\nu$-stationary measures do not exist, hence
if and only if every ergodic $\nu$-stationary measure $\mu$ on $X$ satisfies one of the following properties:
\begin{itemize}
\item $\mu$ has a positive Lyapunov exponent and its stable distribution depends non-trivially on the itinerary;
\item the two Lyapunov exponents of $\mu$ are strictly positive.
\end{itemize} \end{cor}

\Subsubsection{Groups with invariant curves}

\begin{pro}\label{pro:invariant_curve}
Let $X$ be a compact complex surface. Let
$\Gamma$ be a subgroup of $\Aut(X)$ that preserves a complex curve $C\subset X$. If $\nu$ is a probability measure on $\Gamma$ satisfying~\eqref{eq:moment}, then $\nu$ is not uniformly expanding.
\end{pro}

\begin{rem}\label{rem:invariant_curve}
We leave the reader check that the proof adapts to the real case in the following sense: if $X$, $\Gamma$ and $C$ are defined over $\R$ and $C(\R)$ is of dimension~$1$ (that~is, neither empty nor a finite set), then $\nu$ is not uniformly expanding in restriction to~$C(\R)$.
\end{rem}

\begin{lem} \label{lem:criterion_chung_curves}
Let $C$ be a compact Riemann surface. Then, $\Aut(C)$ does not support any uniformly expanding probability measure.
\end{lem}

\begin{proof}
Let $\kappa$ be a Kähler form on $C$ that satisfies $\int_C\kappa = 1$. For every $f\in \Aut(C)$,
$\int_C f^*\kappa = 1= \int_C \norm{D_xf}^2 \kappa= 1$, so by the Jensen inequality
$\int_C \log \norm{D_xf } \kappa\leq 0$. Now, if $\nu$ is any probability measure on $\Aut(C)$, then
\[
\int_C \int_{\Aut(C)} \log \norm{D_xf } d\nu(f) \kappa\leq 0,
\]
hence Property~\eqref{eq:UEbisrepetita} cannot be satisfied
by $\nu$ (for any $n_0\geq 1$).
\end{proof}

Note that the same argument applies to conformal diffeomorphisms, in particular for $C^1$
diffeomorphisms of $\mathbb S^1$. Lemma~\ref{lem:criterion_chung_curves} and Remark~\ref{rem:restriction_UE}
imply Proposition~\ref{pro:invariant_curve} when~$C$ is smooth; we now prove Proposition~\ref{pro:invariant_curve} in full generality.

\begin{proof}[Proof of Proposition~\ref{pro:invariant_curve}]
Arguing by contradiction, we~assume that $\nu$ is uniformly expanding. Let $\Gamma_1\leq \Gamma$
be the finite index subgroup fixing each component of $C$,
and each of its branches at each of its singular points; let $\nu_1$ be
the hitting measure on~$\Gamma_1$ associated to
$\nu^{(n_0)}$, where $n_0$ is as in Equation~\eqref{eq:UEbisrepetita}. By~Proposition~\ref{pro:induced_UE}, $\nu_1$~is uniformly
expanding, so by replacing $\nu$ by $\nu_1$ and $C$ by one of its components
we assume now that $C$ is irreducible and all branches at its
singular points are fixed by~$\Gamma$. To get a
contradiction we will construct a stationary measure $\mu$ supported on~$C$
such that the tangential Lyapunov exponent along $TC$ is non-positive.

By Lemma~\ref{lem:criterion_chung_curves} we may assume that
the singular set $\Sing(C)$ is non-empty.
If the genus of $C$ is $\geq 0$, the invariance of $\Sing(C)$ forces $\Gamma\rest{C}$ to be finite, in contradiction with the
uniform expansion of $\nu$. Thus, $C$ is a rational curve;
denote by $\pi:\hat C\to C$ its normalization and by $\hat \Gamma\subset \Aut(\hat{C})\simeq \PGL_2(\C)$
the group induced by $\Gamma$; the measure~$\nu$ induces a
measure $\hat \nu$ on $\hat \Gamma$.
Fix $\hat p\in \hat C$ such that $p:=\pi(\hat p)$ is singular. The germ of curve given by $\hat C$ at $\hat p$ determines one of the branches of $C$ at $p$; our assumptions imply that
$\hat p$ is fixed by $\hat\Gamma$. There are local coordinates $t\in (\C,0)$ for $(\hat C, \hat p)$ and $(z,w)\in (\C^2, (0,0))$ for $(X, p)$ in which
$\pi$ is expressed as a Puiseux expansion
\begin{equation*}
t\mto (\pi_1(t), \pi_2(t)) = (\alpha t^q , \beta t^r) \quad\text{modulo higher order terms},
\end{equation*}
where $1\leq q<r$; if $q=1$ the branch is smooth at $p$. In~these coordinates, the tangent direction to $C$ at $p$ corresponding to the branch determined by $\hat p$ is
given by $(1,0)\in \C^2$.
Let $\lambda_{(\hat C, \hat{p})}$ be the
Lyapunov exponent of $\hat \nu$ at $\hat p$, and $\lambda_{(C, p)}$ be the Lyapunov exponent of $\nu$
in the tangent direction of this branch.

\begin{lem}\label{lem:tangent_singular}
With notation as above
$\lambda_{(C, p)} = q \lambda_{(\hat C, \hat{p})}$. In~particular $\lambda_{(C, p)}$ and $\lambda_{(\hat C, \hat{p})}$
have the same sign.
\end{lem}

\begin{proof}
Pick $f\in \Gamma$, write
$f(z,w) = (f_1(z,w), f_2(z,w))$ in the local coordinates $(z,w)$, and expand $f_1$ in power series: $f_1(z,w) = \sum_{i,j} a_{i,j}z^iw^j$.
Since the branch determined by $\lambda_{(\hat C, \hat{p})}$ is $f$-invariant, we~have $D_{p}f(1,0) = (a_{1,0}, 0)$ with $a_{1,0}\neq 0$. Thus,
\begin{equation*} f_1(\pi(t))
= \sum_{i,j=0}^\infty a_{i,j}\alpha^i\beta^j t^{qi+rj}
= a_{1, 0} \alpha t^q \mod(t^{q+1}).
\end{equation*}
Now, $f$ lifts to an automorphism $\hat f$ of $\hat C$ fixing $\hat p$. Writing
$\hat f(t) = \lambda t \bmod(t^2)$, we~get $\pi_1(\hat f(t)) = \alpha \lambda^q t^q \bmod(t^{q+1})$. Then,
the semi-conjugacy $f_1(\pi(t)) = \pi_1(\hat f(t))$ gives $ \lambda^q = a_{1, 0}$, and we are done.
\end{proof}

We resume the proof of Proposition~\ref{pro:invariant_curve}. We~fix an affine coordinate $s$ on $\hat C\simeq \P^1(\C)$ such that $\hat p = \infty$. Then, every lift $\hat g \in \hat\Gamma$ can be written as an affine map $\hat g(s) = a_g s+b_g$.

\begin{lem}\label{lem:integrability}
The functions $\log \abs{a_g}$ and $\log^+\abs{b_g}$ are $\nu$-integrable and \hbox{$\ee( \log \abs{a_g})\!<\!0$}.
\end{lem}

\begin{proof}
For the spherical metric, the derivative of $\hat g$ at $ \infty$ in $\hat C$ is $1/a_g$. The computations
of Lemma~\ref{lem:tangent_singular} show that the derivative
of $g$ acting on $X$ in the direction of the branch of $C$ at $\pi(\infty)$
is $1/a_g^q$ for some $q\geq 1$. So~\eqref{eq:moment} implies that $\ee( \abs{\log \abs{a_g}})<\infty$.
Since $\nu$ is uniformly expanding, this direction
is repelling on average: by Lemma~\ref{lem:tangent_singular}, we~get
$\ee( \log \abs{a_j\inv})>0$.
To estimate $\abs{b_g}$, we~note that
$\dist_X(\pi(s), p) \asymp \abs{s}^{-q}$ when $s\in \C$ approaches $\infty$. Changing the affine coordinate $s$ if necessary, we~may assume that
$\pi(0)\neq p$. We~get
\begin{equation*}
\unsur{\abs{b_g}^q} \asymp \dist_X (\pi(\hat g(0)), \pi(\infty)) = \dist_X(g(\pi(0)), g(p))\leq \norm{g}_{C^1} \dist_X(\pi(0), p).
\end{equation*}
From this and~\eqref{eq:moment} it follows that $\ee(\log^+\abs{b_g})<\infty$.
\end{proof}

The integrability provided by Lemma~\ref{lem:integrability} now allows us to construct a stationary
measure with full mass in the affine chart $\C\subset C$ with non-positive Lyapunov exponent
(relative to the affine metric).
This is classical, we~briefly recall the argument for completeness (see~\cite{brandt}).
For $\omega = (g_n)_{n\geq 0}$, write $g_n(s) = a_ns+b_n$, and
consider the sequence of right products
$r_n(\omega) = g_0\cdots g_{n-1}$. One easily checks that
\begin{equation}\label{eq:series}
r_n(\omega) (s) = a_{0}\cdots a_{n-1}s + \sum_{j=0}^{n-1} a_{0}\cdots a_{j-1} b_{j}.
\end{equation}
For $\nu^\N$-almost every $\omega$, $\unsur{n} \log \abs{a_{0}\cdots a_{n-1}}$ converges to $\lambda := \ee(\log\abs{a_g})<0$.
Fix $\e < \abs{\lambda}$. Since $\ee(\log^+\abs{b_g})<\infty$, $\sum_{j=0}^\infty \nu\{ \abs{b_g} > e^{\e j} \} <\infty$. By~the Borel-Cantelli Lemma, $\abs{b_{j}}\leq e^{\e j} $ for $\nu^\N$-almost every $\omega$ and for large $j$; hence, the series on
the right hand side of~\eqref{eq:series} converges. It follows that $r_n(\omega) (s) $ converges almost surely to a limit $e_\omega$
that does not depend on $s\in \C$. The distribution of $e_\omega$ is the desired stationary measure $\mu_\C$.
If $\mu$ is any stationary measure with $\mu(\C) = 1$, then $r_n(\omega)_\varstar \mu$ converges to
$\delta_{e_\omega}$ almost surely: this shows that $\mu_\C$ is the unique stationary measure
with $\mu(\C) = 1$; in particular, $\mu_\C$ is ergodic.
Since the affine derivative of $g$ is the constant $a_g$,
the Lyapunov exponent of $\mu_\C$, relative to the affine metric, is equal to~$\lambda$.

To conclude the proof, note that $\mu:=\pi_\varstar (\mu_\C)$ is an ergodic
$\nu$-stationary measure on~$X$ which has a well-defined Lyapunov exponent, thanks to the moment condition~\eqref{eq:moment}.
If $\mu$ gives positive mass to the singular set of $C$, then it must be concentrated on a
single singular point of $C$ (and likewise $\mu_\C$ is a single atom in~$\hat C$). By~Lemma \ref{lem:tangent_singular} the corresponding branch is
attracting on average, which contradicts uniform expansion. Therefore $\mu$ gives no mass to
$\Sing(C)$, and we claim that its Lyapunov exponent $\lambda(\mu)\rest{TC}$
in the direction of $C$ equals $\lambda$ (even if the ratio between the
ambient and affine metrics on $\C\subset C$ is unbounded). Indeed, for
$\mu\times \nu^\N$-almost every $(x,\omega)$ and $v\in T_x^1C$, we~can fix a subsequence $n_j$ such that $f^{n_j}_\omega(x)$ is far
from the singularities of $C$ (hence from $p=\pi(\infty)$). If $j$ is large,
$\unsur{n_j}\log\norm{D_xf^{n_j}_\omega}$ is both close to $\lambda$ and to
$\lambda (\pi_\varstar \mu)\rest{TC}$. We~conclude that $\lambda (\pi_\varstar \mu)\rest{TC}<0$, which again is contradictory. The proof is complete. \end{proof}

\subsubsection{Zariski diffuse measures}
From now on we focus on the case of a minimal Kähler surface
$X$ of Kodaira dimension zero, that is,
a torus, a K3 surface, or an Enriques surface. In~this case
$\Aut(X)$ preserves a canonical volume form $\vol_X$ (see Example~\ref{eg:K3_intro}).\enlargethispage{\baselineskip}%

From Corollary~\ref{cor:criterion_chung}, the obstruction
to uniform expansion is the existence of a non-expanding stationary measure $\mu$. Moreover, in the first case of Definition~\ref{defi:non_expanding_measure}, both
exponents must vanish because we are in a volume preserving setting. In~this situation,
Theorems~\ref{thm:ledrappier_invariance_principle}
and~\ref{thm:classification_proj_invariant} give a precise description of $\mu$.

\begin{thm}\label{thm:rigidity}
Let $X$ be a torus, a K3 surface, or an Enriques surface. Let $\nu$ be a probability measure on
$\Aut(X)$ satisfying~\eqref{eq:moment} such that $\Gamma_\nu$ is non-elementary.
If $\mu$ is a Zariski diffuse $\nu$-stationary measure, the following properties are equivalent
\begin{enumeratea}
\item\label{thm:rigiditya}
$\mu$ is non-expanding;
\item\label{thm:rigidityb}
the fiber entropy $h_\mu(X, \nu)$ vanishes.
\end{enumeratea}
Moreover under these assumptions, $\mu$ is invariant and $h_\mu(f)=0$ for every $f\in \Gamma_\nu$.
\end{thm}

\begin{proof}
As a preliminary step, observe that almost every ergodic component of $\mu$ is
Zariski diffuse: this follows from the fact that there are
only finitely many invariant curves and countably many isolated periodic points. In~addition, by~convexity
of the entropy (see~\cite[Prop.\,4.3.16]{katok-hasselblatt}), if $h_\mu(X, \nu)=0$ then almost every ergodic component of $\mu$ has zero fiber entropy as well.
Thus for both implications we may further assume that $\mu$ is ergodic as a stationary measure.

Let us show that Property \eqref{thm:rigiditya} implies Property \eqref{thm:rigidityb}. Since there is an invariant volume form, either
both Lyapunov exponents of $\mu$ vanish or $\mu$ is hyperbolic. In~the first case,
the invariance principle guarantees that $\mu$ is $\Gamma_\nu$-invariant and
the fibered version of the Ruelle inequality (see \eg \cite[\S 7]{stiffness}) implies that its fiber entropy
vanishes. If $\mu$ is hyperbolic, the invariance of $\mu$ and the vanishing of the entropy follow from
\cite[Th.\,9.1]{stiffness}.
Thus, \eqref{thm:rigiditya} implies \eqref{thm:rigidityb} together with the invariance of $\mu$.

Consider the converse implication. Again, if $\mu$ has zero Lyapunov exponents then it is non-expanding and
invariant. Otherwise it is hyperbolic and
by applying the whole argument of~\cite{br} in the complex case (see Remark~\ref{rmk:conditional} below for a few comments on this generalization), we~infer that if the stable directions of
$\mu$ depend on the itinerary, its conditionals along Pesin unstable manifolds admit a non-trivial
translation invariance; in particular they are non-atomic. It follows that $h_\mu(X, \nu)>0$
(see also~\cite[Rem.\,9.2]{stiffness}). So under assumption \eqref{thm:rigidityb} the stable directions are non-random and, as
already explained, $\mu$ is invariant by~\cite[Th.\,9.1]{stiffness}.

The fact that $h_\mu(f)=0$ for all $f\in \Gamma_\nu$ will be shown in
Theorem~\ref{thm:entropy_every_f}.
\end{proof}

\subsubsection{Refined criterion}

The discussion of the previous paragraphs leads to a version of
Theorem~\ref{thm:criterion_uniform_expansion_parabolic} that does not require $\Gamma_\nu$ to contain parabolic elements:

\begin{thm}\label{thm:criterion_uniform_expansion}
Let $X$ be a compact Kähler surface which is not rational. Let $\nu$ be a probability measure on
$\Aut(X)$ satisfying~\eqref{eq:moment} such that $\Gamma_\nu$ is non-elementary.
Then $\nu$ is uniformly expanding if and only if the three following conditions
hold:
\begin{enumerate}
\item\label{thm:criterion_uniform_expansion1} every finite $\Gamma_\nu$-orbit is uniformly expanding;
\item\label{thm:criterion_uniform_expansion2} there is no $\Gamma_\nu$-invariant algebraic curve;
\item\label{thm:criterion_uniform_expansion3} there is no Zariski diffuse invariant measure $\mu$ with zero fiber entropy.
\end{enumerate}
\end{thm}

\begin{proof} If a compact Kähler surface $X$ is ruled (over a curve of positive genus) or has a positive Kodaira dimension, then $\Aut(X)$ is elementary (in the first case, it preserves the ruling; in the second case, it preserves the Kodaira-Iitaka fibration, acting as a finite group on the base). Thus, the Kodaira dimension of $X$ vanishes. If $X$ is not minimal, the uniqueness of the minimal model shows that there is a $\Aut(X)$-invariant curve, and we know this is incompatible with uniform expansion (Proposition~\ref{pro:invariant_curve}). Now if ${\mathsf{kod}}(X)=0$, $X$ is minimal, and $\Aut(X)$ is non-elementary, then $X$ is a torus, a K3 surface, or an Enriques surface; hence, we~can assume that $X$ is such a surface.

If $\nu$ is uniformly expanding, Property~\eqref{thm:criterion_uniform_expansion1} is obvious, Property~\eqref{thm:criterion_uniform_expansion2} follows from Proposition~\ref{pro:invariant_curve}, and
Property~\eqref{thm:criterion_uniform_expansion3} follows from Corollary~\ref{cor:criterion_chung} and Theorem~\ref{thm:rigidity}.

Conversely, if these properties hold, and if $\mu$ is an ergodic $\nu$-stationary measure then by Property~\eqref{thm:criterion_uniform_expansion2}
$\mu$ is either Zariski diffuse or finitely supported. Then, Theorem~\ref{thm:rigidity} and Property~\eqref{thm:criterion_uniform_expansion1} imply that $\mu$ is not non-expanding, and we conclude with Corollary~\ref{cor:criterion_chung}.
\end{proof}

\begin{proof}[Proof of Theorem~\ref{thm:criterion_uniform_expansion_parabolic}]
This follows directly from Theorem~\ref{thm:hyperbolic} and Theorem~\ref{thm:criterion_uniform_expansion}.
\end{proof}

\begin{rem}\label{rmk:conditional}
The proof of $\eqref{thm:rigidityb}\Rightarrow\eqref{thm:rigiditya}$ in Theorem~\ref{thm:rigidity} relies on the following fact: {\textit{for a hyperbolic stationary measure, if the stable directions of $\mu$ depend on the itinerary, then its unstable conditionals satisfy some non-trivial translation invariance}}. This is the
``easy part'' of the adaptation of~\cite{br} to complex surfaces; the ``difficult part'' would be to obtain
stiffness and some SRB property from this invariance (either on $X$ or on some totally real surface associated to the stationary measure). We~did not provide a proof for this fact because the arguments of~\cite{br} can be applied directly. This program was recently carried out by Roda and will appear shortly.
As a consequence, this fact is also used in the implication ``uniformly expanding implies \eqref{thm:criterion_uniform_expansion3}'' in
Theorem~\ref{thm:criterion_uniform_expansion}. On the other hand, it is not used in Theorem~\ref{thm:criterion_uniform_expansion_parabolic} because in this case the condition~\eqref{thm:criterion_uniform_expansion3} of
Theorem~\ref{thm:criterion_uniform_expansion} is automatically
satisfied, thanks to Theorem~\ref{thm:hyperbolic}; it is not used either for the part of Theorem~\ref{thm:criterion_uniform_expansion} asserting that the assumptions \eqref{thm:criterion_uniform_expansion1}, \eqref{thm:criterion_uniform_expansion2} and \eqref{thm:criterion_uniform_expansion3} imply uniform expansion.
\end{rem}

\begin{rem}
Using Theorem~\ref{thm:hyperbolic_kummer} instead of Theorem~\ref{thm:hyperbolic} gives a version of
Theorem~\ref{thm:criterion_uniform_expansion_parabolic} where the existence of a parabolic element in $\Gamma$ is replaced by the existence of a Kummer element. The details of the adaptation are left to the interested reader.
\end{rem}

\subsection{Uniform expansion along finite orbits}\label{subs:finite_orbits}

Using classical results on random products of matrices, it is easy
to characterize when a fixed point under $\Gamma_\nu$ is uniformly expanding. We~say that a subgroup of $\GL_2(\C)$ is
\emph{strictly triangular} if it is reducible with exactly one invariant direction.

\begin{pro}\label{pro:fixed_UE}
Let $X$ be a torus, a K3 surface, or an Enriques surface. Let
$\nu$ be a probability measure on $\Aut(X)$ satisfying~\eqref{eq:moment}, and let $x_0$ be a fixed point of
$\Gamma_\nu$. Then $\nu$ is uniformly expanding on $T_{x_0}X$ if and only if one of the following holds
\begin{enumeratea}
\item\label{pro:fixed_UEa} the
induced action of $\Gamma_\nu$ on $ T_{x_0}X$ is non-elementary;
\item\label{pro:fixed_UEb} this action is strictly triangular and its invariant direction
is expanding.
\end{enumeratea}
If $\nu$ is symmetric, it is uniformly expanding on $T_{x_0}X$ if and only if \eqref{pro:fixed_UEa} holds.
\end{pro}

In case \eqref{pro:fixed_UEb} there exists $ u\in T_{x_0}X$ such that $f_\varstar u = \lambda_f u$ for every $f\in \Gamma_\nu$, and the expansion means that $\int \log \abs{\lambda_f} d\nu(f) >0$.

\begin{proof}[Proof (see also~\cite{Prohaska-Sert:TAMS})]
By Lemma~\ref{lem:weak_UE2}, to prove uniform expansion it is enough to show that for every $ v\in T_{x_0}X$, $\liminf_{n\to\infty}\unsur{n}\log \norm{(f_\omega^n)_\varstar v} >0$.
The proof is based on the work of Furstenberg and Kifer~\cite{furstenberg-kifer} (see also~\cite[\S 3.7]{bougerol-lacroix}). These references deal with general random products of matrices in $\GL_d(\R)$; in our volume preserving situation the Lyapunov exponents $\lambda_2\leq \lambda_1$ of the random
product in $\GL_2(\C)$
satisfy $\lambda_1+\lambda_2=0$, so they can be read off directly from the action
on
$\P T_{x_0}X$. According to \cite[Th.\,3.5 \& 3.9]{furstenberg-kifer}, there are two possibilities:
\begin{enumeratei}
\item\label{case:i} for every $ v\in T_{x_0}X$ and $\nu^\N$-almost every $\omega$, $\unsur{n}\log \norm{(f_\omega^n)_\varstar v} \to \lambda_1$;
\item\label{case:ii} there
exists a
non-random, $\Gamma_\nu$-invariant filtration $\set{0}=L_2< L_1< L_0 = T_{x_0}X$
and $\beta_1<\beta_0$
such that for $i=0,1$ for any $ v\in L_{i}\setminus L_{i+1}$, for $\nu^\N$-almost every $\omega$,
$\unsur{n}\log \norm{(f_\omega^n)_\varstar v} \to \beta_i$. Furthermore $\beta_0= \lambda_1$.
\end{enumeratei}
We now compare this dichotomy with the classification of subgroups of $\PGL_2(\C)$ (with a slight abuse of notation, we~also denote by $\Gamma_\nu$ the induced subgroup of $\PGL_2(\C)$).

$-$ If $\Gamma_\nu$ is strongly irreducible, we~are in case \eqref{case:i} and there are two possibilities.
If $\Gamma_\nu$ is proximal (hence non-elementary)
then $\lambda_1>0$ and $\nu$ is uniformly expanding. If $\Gamma_\nu$ is not proximal, it is contained in a
compact subgroup and $\nu$ is not uniformly expanding.

$-$ If $\Gamma_\nu$ is irreducible but not strongly irreducible, we~are in case \eqref{case:i} and there are two lines which are permuted by $\Gamma$. In~some affine coordinate $z$ on $\P T_{x_0}X$, $\Gamma_\nu$ is then conjugated to a subgroup of
\begin{equation*}
\set{z\mto \lambda z^{\e}\sep \lambda\in \C^\times, \ \e = \pm 1}
\end{equation*} where $\e=- 1$ with positive probability. In~this case $\lambda_1 = 0$ (see \eg \cite[Prop.\,5.3]{degenerations_SL2}), so $\nu$ is not uniformly expanding.

$-$ If $\Gamma_\nu$ is reducible it preserves one or two directions in $T_{x_0}X$. If $\Gamma$ preserves a direction with a non-positive exponent, then $\nu$ is not uniformly expanding. So, we~can assume that $\Gamma$ preserves a unique direction, and that the corresponding exponent $\beta$ is positive. By~\eqref{case:i} and \eqref{case:ii} we see that $\lim_{n\to\infty}\unsur{n}\log \norm{(f^n_\omega)_\varstar v}\geq \beta$ for any $ v\in T_{x_0}X$ and almost every $\omega$; so
$\nu$ is uniformly expanding.

This covers all possible cases and the proof is complete.
\end{proof}

Let $F$ be a finite set, viewed as a $0$-dimensional manifold, and $V$ be a real or complex vector bundle of dimension $d$ over $F$; identify
$V$ with $ F\times \mathbf{K}^d$, for ${\mathbf{K}} = \R$ or $\C$. Let $\GL(V)$ be the group of bijections of $V$ acting linearly on fibers: it is a semi-direct product $\GL(V) \simeq \mathfrak{S}(F) \ltimes \GL_d({\mathbf{K}})^F$ where $\mathfrak{S}(F)$ acts on
$\GL_d({\mathbf{K}})^F$ by permuting the factors. We~say that a subgroup of $\GL(V)$ is \emph{strongly irreducible} if it acts transitively on $F$
and the stabilizer of any $x\in F$ acts strongly irreducibly on the fiber $\{ x\} \times {\mathbf{K}}^d$ of $V$; or equivalently, if there is no invariant and finite collection of proper positive dimensional subspaces in some fibers of $V$. Similar notations and notions are defined for $\PGL(V)$.

Now, if $F$ is a finite $\Gamma$-orbit on $X$,
consider the induced action of $\Gamma$ on $TX\rest{F}:=\bigcup_{x\in F} T_x X$. We~say that this action is \emph{non-elementary}
if its image in $\PGL\lrpar{TX\rest{F}}$ is strongly irreducible and unbounded. When
$\Gamma$ preserves a volume form on $X$, its image in $\GL\lrpar{TX\rest{F}}$ is unbounded if and only if it is unbounded in $\PGL\lrpar{TX}$. We~say that it is \emph{strictly triangular} if the only proper $\Gamma$-invariant subbundle in $T X\rest{F}$ is given by a $1$-dimensional subbundle $L\subset TX\rest{F}$.

Pick a point $x$ in $F$ and set $\Gamma_x = \Stab_{\Gamma}(\set{x})$. Since $F$ is an orbit, $[\Gamma:\Gamma_x]=\abs{F}$ and the image of $\Gamma$ in $\PGL\lrpar{TX\rest{F}}$ is unbounded if and only if the image of $\Gamma_x$ in $\PGL\lrpar{T_xX}$ is unbounded. Thus, one easily gets the following lemma.

\begin{lem}\label{lem:stab_NE}
If $F$ is a finite $\Gamma$-orbit, the action of
$\Gamma$ on $TX\rest{F}$ is non-elementary (\resp strictly triangular) if and only if for some, hence any, $x\in F$
the action of $\Stab_{\Gamma}(\set{x})$ on $T_{x}X$ is non-elementary (\resp strictly triangular).
\end{lem}

\begin{thm}\label{thm:finite_UE}
Let $X$ be a torus, a K3 surface, or an Enriques surface. Let
$\nu$ be a probability measure on $\Aut(X)$ satisfying~\eqref{eq:moment+}, and $F$ be a finite $\Gamma_\nu$-orbit.
Then $\nu$ is uniformly expanding on $F$ if and only if the
induced action of $\Gamma_\nu$ on $TF$ is
\begin{enumeratea}
\item\label{thm:finite_UEa} either non-elementary (in the sense of the above definition);
\item\label{thm:finite_UEb} or strictly triangular and the field of invariant directions $L\subset TX\rest{F}$ is uniformly expanding.
\end{enumeratea}
If $\nu$ is symmetric, it is uniformly expanding on $F$ if and only if (a) holds.
\end{thm}

\begin{proof}
Let $\Gamma_F$ be the finite index subgroup fixing every point of $F$.
Assume that~$\nu$ is uniformly expanding. Then by Proposition~\ref{pro:induced_UE}, for some
$n_0$, the induced measure $(\nu^{(n_0)})_{\Gamma_F}$ is uniformly expanding. Therefore, by~Proposition~\ref{pro:fixed_UE}, $\Gamma_F$ satisfies Property~\eqref{thm:finite_UEa} or~\eqref{thm:finite_UEb} at every point of $F$, and
we conclude by Lemma~\ref{lem:stab_NE}.
Conversely, assume that~\eqref{thm:finite_UEa} or~\eqref{thm:finite_UEb} holds. Note that by Theorem~\ref{thm:M$_2$},
$\nu_{\Gamma_F}$ satisfies~\eqref{eq:moment+}. By~Lemma~\ref{lem:stab_NE} and Proposition~\ref{pro:fixed_UE},
$\nu_{\Gamma_F}$ is uniformly expanding on $F$,
hence by Proposition~\ref{pro:induced_UE2}, $\nu$ is uniformly expanding on $F$, as desired.
\end{proof}

This theorem shows that when $\nu$ is symmetric
all conditions in Theorem~\ref{thm:criterion_uniform_expansion} depend only on
$\Gamma_\nu$, and not on $\nu$. Thus we obtain:

\begin{cor}\label{cor:independence_nu}
Let $X$ be a torus, a K3 surface, or an Enriques surface. Let $\Gamma$ be a non-elementary subgroup of $\Aut(X)$. Let $\nu$ and $\nu'$ be symmetric probability measures
on $\Aut(X)$ satisfying~\eqref{eq:moment+} such that $\Gamma_\nu = \Gamma_{\nu'}=\Gamma$. Then $\nu$ is uniformly expanding if and only if $\nu'$ is
uniformly expanding.
\end{cor}

\section{Examples of uniformly expanding actions}

\Subsection{A finitary version of Theorem~\ref{thm:criterion_uniform_expansion_parabolic} and application to Wehler surfaces}\label{subs:finitary}

In~\cite[\S\S 7-8]{chung}, Chung uses computer assistance to prove the uniform expansion of
some concrete algebraic actions on real surfaces. In~our situation
Theorem~\ref{thm:criterion_uniform_expansion_parabolic} can be used to
check uniform expansion, but this requires a description of all invariant Zariski-closed subsets.
As already explained, invariant curves can be determined by cohomological computations; for instance, if $X$ is a generic Wehler surface, there is no $\Aut(X)$-invariant curve. Thus the main problem is to study finite orbits.

If the group $\Gamma$ is non-elementary, contains parabolic elements, and has no invariant curve, the main
result of~\cite{finite_orbits} says that $\Gamma$ admits only finitely many finite orbits, except when $(X,
\Gamma)$ is a Kummer example. However, the proof given in~\cite{finite_orbits} does not provide any bound
on the number or the lengths of such orbits; so, there is a~priori no hope of numerically checking
uniform expansion along all of them, nor proving that there are no finite orbits. The next result explains how to overcome this issue.
To~state it, we~denote by $\NS(X;\R)$ the Néron-Severi group of $X$ that is, the subgroup of $H^{1,1}(X;\R)$ obtained by all Chern classes of holomorphic line bundles on $X$; it coincides with the intersection of $H^{1,1}(X;\R)$ with the torsion free part of $H^2(X;\Z)$ (see~\cite[p.\,163]{griffiths-harris}).

\begin{thm}\label{thm:effective}
Let $X$ be a smooth projective surface and
$\Gamma$ be a non-elementary subgroup of $\Aut(X)$ containing parabolic elements, which does not
preserve any algebraic curve. Assume that we are given:
\begin{enumeratei}
\item\label{thm:effectivei}
algebraic equations for $X$, and the formulas defining a generating subset $S$ of~$\Gamma$;
\item\label{thm:effectiveii}
a basis of $\NS(X;\R)$ and the matrices of $s^\varstar\colon \NS(X;\R)\!\to\! \NS(X;\R)$, for~$s$ in~$S$;
\item\label{thm:effectiveiii}
a parabolic element $g\in \Gamma$, given as a word in the generators $s\in S$,
and its invariant fibration $\pi\colon X\to B$.
\end{enumeratei}
Then, there is an analytically computable integer $N(X, \Gamma)$
such that every finite $\Gamma$\nobreakdash-orbit of length greater than $N(X,\Gamma)$ is uniformly expanding (in the sense of Section~\ref{subs:UE_intro}).
\end{thm}

By \emph{analytically computable}, we~mean
computable by a computer able to solve real analytic equations; by \emph{algebraically computable}, we~mean
computable by a computer able to solve algebraic equations. The proof will provide an analytically computable subset containing all possible non-expanding finite orbits.

\begin{example} Let $g$ be a parabolic
element of $\Gamma$, and let $h\in \Gamma$ be a conjugate of~$g$ with a distinct invariant fibration.
Denote by $\mathrm{Tor}_N(g)$ the finite set of
fibers of the $g$-invariant fibration in which $g$ is a periodic translation of period $\leq N$.
Then, the set of finite orbits of $\Gamma$ of length $\leq N$ is algebraically computable since it is contained in
\begin{equation*}
\mathrm{Tor}_N(g)\cap \mathrm{Tor}_N(h) =\{x\in X\sep  g^N(x)=h^N(x)=x \}.
\end{equation*}
\end{example}

A typical application of Theorem~\ref{thm:effective} is to the Wehler family.
Recall from Section~\ref{subs:wehler}
that ${ \cW}_0$ is the family of Wehler surfaces which are smooth and do not contain any
fiber of the three natural projections to
$(\pp^1)^2$.
Under these assumptions the group $\Gamma$ generated by the three basic involutions
$\sigma_1$, $\sigma_2$ and $\sigma_3$ is non-elementary and
has no invariant curve (see~\cite[Prop.\,2.2]{finite_orbits}).
It turns out that in this case $N(X, \Gamma)$ is constant on a Zariski
dense open subset (see Proposition~\ref{pro:N(X)_bounded} below).
This leads to the following theorem, which will be proved in Section~\ref{par:proof_of_9.3}:

\begin{thm}\label{thm:wehler_zariski}
There is a dense Zariski open subset of ${ \cW}_0$
(\resp of the family ${ \cW}_0(\R)$ of real Wehler surfaces),
in which the action of $\Gamma=\langle \sigma_1, \sigma_2, \sigma_3\rangle$ is
uniformly expanding on $X$.
\end{thm}

\subsubsection{Preliminaries on Halphen twists}\label{par:preliminaries2_on_Halphen} Let us resume the discussion from Section~\ref{subs:halphen} and add a few preliminaries on Betti foliations and the non-twisting locus. Let $h$ be a Halphen twist with associated fibration $\pi\colon X\to B$.
Consider a simply connected open subset $U$ of $B^\circ$ together with a section $\sigma\colon U\to X$ of $\pi$
and a continuous frame for the homology of the fibers above $U$.
For $w\in U$, one can identify the fiber $X_w$ to $\C/\Lat(w)$ ($\sigma(w)$ corresponding to the zero of $\C/\Lat(w)$), as in Section \ref{subs:halphen}.
Then, above $U$, there is a unique real-analytic diffeomorphism $\Psi: \pi \inv(U)\to U\times \R^2/\Z^2$
such that
\begin{enumeratea}
\item
$\pi\circ \Psi=\Psi\circ \pi_U$, where $\pi_U$ is the projection onto $U$;
\item
$\Psi$ maps $\sigma$ to the zero section $w\mto (w, (0,0))$ of $\pi_U$, and maps the basis of $H_1(X_w;\Z)$ to the standard basis of $H_1(\R^2/\Z^2; \Z)=\Z^2$;
\item
on each fiber, $\Psi$ is a real analytic isomorphism of real Lie group.
\end{enumeratea}

Above $U$, the \emph{Betti foliation} is the foliation by
submanifolds of the form $\Psi\inv(U\times\nobreak \set{(x,y)})$; these leaves are local holomorphic sections of $\pi$, with $\sigma$ corresponding to $\Psi\inv(U\times \set{(0,0)})$.
Conjugating by
$\Psi$, we~get
\begin{equation}\label{eq:definition_T}
\Psi\circ h\circ \Psi^{-1}\colon (w, (x,y))\mto (w, (x,y)+T(w)),
\end{equation}
where $T\colon U\to \R^2/\Z^2$ is real analytic. By~\cite[Lem.\,3.9]{invariant}, the map $T$
is an (orientation preserving)
branched covering, so it behaves topologically like $w\mto w^k$. In~$U$, the
\emph{non-twisting locus} $\mathrm{NT}_{h}$
is the set $\{w\in U\sep  D_wT = 0\}$; equivalently, $\mathrm{NT}_{h }\cap U = \pi(\set{t_1, \ldots , t_q})$, where
$\set{t_1, \ldots , t_q}$ is the set of tangencies between
the Betti foliation and the section $h \circ \sigma$. These definitions do not depend on the above choices and $\mathrm{NT}_{h}$ can indeed be defined globally on $B^\circ$.
A key fact is that $\mathrm{NT}_{h}$ is a finite subset of $B^\circ$
(see~\cite[Prop.\,3.14]{invariant} or~\cite[Cor.\,7.7.10]{duistermaat}). We~denote by $\abs{\NT_h}$ its cardinality, and by $\mult(\NT_h)$
its cardinality counted with multiplicity, that is, taking into account the
degree of the local branched covering~$T$.

Note that, once $h$ and $\pi$ are given, {\textit{the set $\NT_h\subset B^\circ$ is analytically computable}}: one has to compute the periods of $X_w$
to get $\Lat(w)$, then $\Psi$ is $\R$-linear from $\C/\Lat(w)$ to $\R^2/\Z^2$, and $T$ is then obtain from $h$ by conjugacy.

\subsubsection{Proof of Theorem~\ref{thm:effective}}
As in \cite{invariant}, for $(g,h)\in \Hal(\Gamma)^2$
we set
\begin{equation*}
\STang(\pi_g, \pi_h) =
\Sing(\pi_g)\cup\Sing (\pi_h)\cup \Tang^{\mathrm{tt}}(\pi_g, \pi_h),
\end{equation*}
where $\Sing(\pi_g)$ is the union of all singular and multiple fibers, and
$\Tang^{\mathrm{tt}}(\pi_g, \pi_h)$ is the part of the tangency locus of $\pi_g$ and $\pi_h$ which is not contained in $\Sing(\pi_g)\cup\Sing (\pi_h)$.
Put $\NT_g^X=\pi_g^{-1}(\NT_g)$ (so that $\NT_g^X$ is a curve in $X$) and
likewise $\NT_h^X=\pi_h^{-1}(\NT_h)$.

\begin{lem}\label{lem:NT1}
Let $g,h$ be a pair of Halphen twists in $\Gamma$ with distinct invariant fibrations, and let $x\in X$ be a point with a finite $\Gamma$-orbit.
If this orbit is not uniformly expanding, then
it is contained in
$ \STang(\pi_g, \pi_h)\cup \NT_g^X\cup \NT_h^X$.
\end{lem}

\begin{proof}
We argue by contraposition: replacing $x$ by another point in its orbit if necessary, we~assume that $x\notin \STang(\pi_g, \pi_h)\cup \NT_g^X\cup \NT_h^X$, and we want to show that its orbit is
uniformly expanding.
Since $\Gamma(x)$ is finite, there are positive integers $k$ and $\ell$
such that $g^k$~and~$h^\ell$ are in $\mathrm{Stab}_\Gamma(x)$. By~definition of the non-twisting locus,
$g^k$~and~$h^\ell$
induce parabolic homographies on $\pp(T_xX)$; and since $x\notin \Tang^{tt}(\pi_g,\pi_h)$, the
fixed points of these homographies are distinct; thus, the action of
$\langle g^k, h^\ell\rangle$ on $\pp(T_xX)$ is non-elementary. By~Proposition~\ref{pro:fixed_UE} and Lemma~\ref{lem:stab_NE}, the orbit of $x$ is uniformly expanding. \end{proof}

The intersection number of $\NT_g^X \cup \STang(\pi_g, \pi_h)$ (\resp $\NT_h^X \cup \STang(\pi_g, \pi_h)$) with a smooth fiber $X_w^h$ (\resp $X_w^g$) does not depend on the fiber. Let $n_0(g,h)$ be the maximum of these intersection numbers:
\begin{equation}\label{eq:definition_n0}
n_0=\max\bset{[\NT_g^X \cup \STang(\pi_g, \pi_h)] \cdot [X_w^h],\,  [\NT_h^X \cup \STang(\pi_g, \pi_h]\cdot [X_w^g])}.
\end{equation}
The set $\STang(\pi_g, \pi_h)$ can be computed algebraically, thus
\begin{equation*}
n_0(g,h)\leq A(g,h) \max\bset{\vert \NT_g\vert, \vert \NT_h\vert} + B(g,h),
\end{equation*}
where $A(g,h)$ and $B(g,h)$ can be computed algebraically (by computing the tangency loci and intersection numbers).
Then, we~set
\begin{equation}\label{eq:definition_n(g,h)}
n(g,h)=n_0(g,h) !
\end{equation}

\begin{lem}\label{lem:definition_of_bad(g,h)}
Let $g,h$ be a pair of Halphen twists in $\Gamma$ with distinct invariant fibrations. Let
$x\in X$ be such that $\Gamma(x)$ is finite and not
uniformly expanding. Then $\Gamma(x)$ is contained in
\[
\STang(\pi_g, \pi_h) \cup
(\NT_g^X\cap \NT_h^X)\cup (\NT_g^X\cap \NT_{h^n gh^{-n}}^X) \cup (\NT_h^X\cap \NT_{g^nhg^{-n}}^X),
\]
where $n=n(g,h)$ is defined by~\eqref{eq:definition_n(g,h)}.
\end{lem}

\begin{proof}
The statement of the lemma concerns the orbit $\Gamma(x)$, but we only have to prove it for $x$ itself.
If $x\in \STang(\pi_g, \pi_h) \cup (\NT_g^X \cap \NT_h^X)$ we are done.
Otherwise by Lemma~\ref{lem:NT1}, $x$ belongs to
$\NT_g^X\setminus \NT_h^X$ or $\NT_h^X\setminus \NT_g^X$.
Assume that $x\in \NT_g^X\setminus \NT_h^X$.
The $h$-orbit of $x$ is finite and by Lemma~\ref{lem:NT1} again,
for every $q$, $h^q(x)$ is contained in
$X^h_{x}\cap (\NT_g^X \cup \STang(\pi_g, \pi_h))$ (here we abuse notation and
write $X^h_{x}$ for $X^h_{\pi_h(x)}$). Thus,
$h^{n}(x) = x$, where $n=n(g,h)$. Set $f = h^{n} g h^{-n}$. The fiber $X^f_x$ associated to~$f$ through $x$
is $h^{n}( X^g_x)$, and since $x\notin \NT_h^X$, $X^f_x$ is transverse to
$X^g_x$ at $x$, as well as to~$X^h_x$. Moreover, $x$ belongs to $\NT_f^X$, because $x$ belongs to $\NT_g^X$ and $h^n(x)=x$. Hence $x\in \NT_g^X\cap \NT_{h^ngh^{-n}}^X$. Doing the same
in the case where $x\in \NT_h^X\setminus \NT_g^X$ completes the proof.
\end{proof}

The set $\NT_g$ is analytically computable (by Section \ref{par:preliminaries2_on_Halphen}), and $\Crit(\pi_g)$ is algebraically computable. Similarly, if $h$ is in $\Hal(\Gamma)$, $\Tang^{tt}(\pi_g,\pi_h)\cap \NT_f^X$ is analytically computable.
The previous lemma
shows that all non uniformly expanding finite orbits are contained in
\[
\mathrm{Bad}(g,h):=\STang(\pi_g, \pi_h) \cup
(\NT_g^X\cap \NT_h^X)\cup (\NT_g^X\cap \NT_{h^n gh^{-n}}^X) \cup (\NT_h^X\cap \NT_{g^nhg^{-n}}^X)
\]
for every pair $(g,h)\in \Hal(\Gamma)^2$ with distinct invariant fibrations, where $n = n(g,h)$ as in Equation~\eqref{eq:definition_n(g,h)}.
Intersecting these sets for various choices of $(g,h)$, we~expect to get a finite analytically computable set.
Observe that $\mathrm{Bad}(g,h)$
is the union of $\STang(\pi_g, \pi_h)$ and a finite set, because $\NT_g^X\cap \NT_h^X$ is finite when $\pi_g$ and $\pi_h$ are distinct. So, what remains to do is to
\textit{exhibit an explicit finite set of pairs $(g,h)$ such that the
intersection of the $\STang(\pi_g, \pi_h)$ is finite}. We~first treat the case of Wehler surfaces,
which is sufficient to proceed with Theorem~\ref{thm:wehler_zariski}.

\smallskip

\noindent \textit{Conclusion of the proof of Theorem~\ref{thm:effective} in the Wehler case.}
Fix a Wehler surface $X\in {\cW}_0$ and consider the three pairs $(g_1, g_2)$, $(g_2,g_3)$, $(g_3,g_1)$, where $g_1 = \sigma_2\circ \sigma_3$, $g_2 = \sigma_3\circ\sigma_1$ and $g_3 =\sigma_1\circ\sigma_2$. Note that the $g_i$-invariant fibration is the $i$-th projection~$\pi_i$.

Assume that the intersection of the divisors $\STang(\pi_i, \pi_j)$ contains an irreducible curve $D\subset X$. If $D$ is contained in $\Sing(\pi_i)\cap \Sing(\pi_j)$ with $i\neq j$, then $(\pi_i, \pi_j)$ maps $D$ onto a point and this contradicts the fact that $X\in { \cW}_0$. If $D$ is contained in, say, $\Tang^{tt}(\pi_1,\pi_2)$
and $\Tang^{tt}(\pi_2, \pi_3)$, the three fibrations are pairwise tangent along $D$, and we obtain a contradiction because there is no tangent vector $v\neq 0$ to $(\pp^1)^3$ which is mapped to $0$ by each $D\pi_i$. The last possibility is that $D$ is contained~in, say, $\Tang^{tt}(\pi_1,\pi_2)$ and $\Sing(\pi_3)$. In~this case, there is a point $p$ on $D$ at which $D_p\pi_3\colon T_pX\to T_{\pi_3(p)}\P^1$ is equal to $0$, and
at such a point, the same contradiction applies. This shows that
\begin{equation}\label{eq:finite-bad-set}
F(g_1,g_2,g_3):=\mathrm{Bad}(g_1,g_2)\cap \mathrm{Bad}(g_2,g_3)\cap \mathrm{Bad}(g_3,g_1)
\end{equation}
is finite, with an analytically computable cardinality, and the proof is complete. \qed

\begin{rem}\label{rem:N(X)}
The above proof provides a computation of the integer $N(X, \Gamma)$ involving:
\begin{enumerate}
\item\label{rem:N(X)1} algebraic quantities that are constant on $\cW_0$, like $\STang(\pi_i,\pi_j)\cdot \STang(\pi_j,\pi_k)$,
\item\label{rem:N(X)2} $\vert \mathrm{NT}_{g_i}\vert$ for
$i=1, 2, 3$.
\end{enumerate}
{\textit{Therefore, if $\vert \mathrm{NT}_{g_i}\vert \leq B$, then $N(X,\Gamma)\leq N(B)$ for some $N(B)$ depending only on~$B$.}}

Indeed, the number $n$ in Lemma~\ref{lem:definition_of_bad(g,h)} depends only on $n_0$ (see Equations~\eqref{eq:definition_n(g,h)}) and by Equation~\eqref{eq:definition_n0}) $n_0$ is bounded by a function of $B$.
Then, because the norm of $(g_i^{n_0})^*\colon \NS(X)\to \NS(X)$ is bounded by $Cn_0^2$ for some uniform constant $C$, we~obtain
$\NT_{g_i}^X \cap \NT_{g_j^ng_ig_j^{-n}}^X\leq C' n_0^2 B^2$
for some constant $C'$ and the result follows. \qed
\end{rem}

\subsubsection*{Conclusion of the proof of Theorem~\ref{thm:effective} in the general case}%\label{conclusion_proof_thm_effective}
By assumption, $\Gamma$ is non-elementary and has no invariant curve. Let $\Gamma^*$ be its image in $\GL(\NS(X;\Z))$.

If $g$ is the parabolic element given in assumption~\eqref{thm:effectiveiii} of the theorem, up to sign, there is a unique integral primitive class $c(g)\in \NS(X)$ such that $g^\varstar c(g)=c(g)$ and $c(g)\cdot c(g)=0$. By~the assumptions~\eqref{thm:effectiveii} and~\eqref{thm:effectiveiii}, this class can be computed explicitly. An element $f$ of $\Aut(X)$ preserves the $g$-invariant fibration $\pi$ (permuting its fibers) if and only if it fixes $c(g)$. Since $\Gamma$ is non-elementary, it is not contained in the stabilizer of $c(g)$. Thus, according to Proposition~3.2 of~\cite{eskin-mozes-oh}, there is a computable integer $N$, and a composition $f$ of length $N$ in the generators $s\in S$ that does not preserve $c(g)$. Then, $h:= f\circ g \circ f^{-1}$ is a parabolic element of $\Gamma$ with invariant fibration $\pi\circ f\neq \pi$.

Since $g$ and its invariant fibration $\pi$, as well as $f$, are explicit, we~can compute the degree of the subvariety
$\STang(\pi_g, \pi_h)$ (for the embedding $X\subset \P^m(\C)$ given by assumption (i)). Denote by $(C_i)_{i\in I}$ the irreducible components of $\STang(\pi_g, \pi_h)$; we have $\vert I\vert \leq \deg(\STang(\pi_g, \pi_h))$.

Suppose that for each $C_i$, one can exhibit some $f_i\in \Gamma$ for which $f_i(C_i)\not \subset \STang(\pi_g, \pi_h)$. Then the set of pairs
\begin{equation*}
\set{(g,h)} \cup \set{ (f_igf_i\inv, f_i h f_i\inv)\sep  i\in I}
\end{equation*}
satisfies $\vert \bigcap_i\STang(\pi_{g_i}, \pi_{h_i})\vert < +\infty$, and we are done because the cardinality of this finite set is algebraically computable.
So, we~now fix such an irreducible component~$C_i$, and we construct such an~$f_i$.

First, assume that $C_i$ is an irreducible component of $\Tang^{\mathrm{tt}}(\pi_g, \pi_h)$. Then $C_i$ is generically transverse to $\pi$, hence $\deg(g^n(C_i))$ tends to infinity. We~can thus set \hbox{$f_i=g^{n_i}$} for some large enough $n_i$ (explicitly computable from the action on $\NS(X)$).

The second case is when $C_i$ is an irreducible component of $\Sing(\pi_g)\cup\Sing(\pi_h)$;
in particular, its self-intersection $C_i^2$ is $\leq 0$.
By~\cite[Th.\,D]{finite_orbits}, there exists a loxodromic element $f_0\in \Gamma$
without invariant curve; in particular $f_0^{\abs{I}!}(C_i) \neq C_i$. Since $f_0$ is loxodromic, this inequation is equivalent to
$(f_0^{\abs{I}!})_\varstar [C_i] \neq [C_i]$. Indeed, either $C_i^2=0$ and we readily get a contradiction
since a loxodromic element does not fix any non-zero isotropic class, or $C_i^2<0$
and this follows from $f_0^{\abs{I}!}(C_i) \neq C_i$ since
$[C_i]$ determines~$C_i$ when the self-intersection is negative.
Thus, if we set
\begin{equation*}
W_i:=\bset{f\in \GL(\NS(X, \R))\sep f^{\abs{I}!}_\varstar [C_i] = [C_i]},
\end{equation*}
we see that $\Gamma^*$ is not contained in $W_i$. Proposition 3.2 of \cite{eskin-mozes-oh} then provides a computable element $f\in \Gamma$
such that $f\notin W_i$. Now, if $f^q(C_i)$ were contained in $\STang(\pi_g, \pi_h)$ for $0\leq q\leq \vert I\vert$, we~would find two integers $q_1< q_2\leq \vert I\vert$ such that $f^{q_2-q_1}(C_i)=C_i$; in particular, $f^{\vert I\vert !}(C_i)$ would be equal to $C_i$, a contradiction. Thus, there is an iterate $f_i:=f^{q_i}$, with $q_i \leq \vert I\vert$, such that $f_i(C_i)\not \subset \STang(\pi_g, \pi_h)$, and the proof is complete.
\qed

\subsubsection{Proof of Theorem~\ref{thm:wehler_zariski} }\label{par:proof_of_9.3}
Recall that, for Wehler surfaces, $\Gamma = \bra{\sigma_1, \sigma_2, \sigma_3}$.

\begin{pro}\label{pro:N(X)_bounded}
There exists an analytically computable integer $N$ such that for
any Wehler surface $X\in { \cW}_0$, any
finite $\Gamma$-orbit of length $>N$ is non-elementary (hence uniformly expanding by Theorem~\ref{thm:finite_UE}).
\end{pro}

This uniform bound is the main step towards Theorem~\ref{thm:wehler_zariski}. In~view of Remark~\ref{rem:N(X)}, this proposition
follows from Theorem~\ref{thm:effective} and the following uniformity result.

\begin{pro}\label{pro:NT_bounded}
For any $g\in \set{g_1, g_2, g_3}$, the cardinality of $\NT_g$ is uniformly bounded
in $\cW_0$.
\end{pro}

Let $\cX\subset \cW_0\times (\P^1\times \P^1\times \P^1)$ be
the universal family of Wehler surfaces, as in \cite[\S 2]{finite_orbits}.
As $X$ varies in $ \cW_0$, the automorphisms $g_i$ and their
invariant fibrations $\pi_i$ depend on $X$, but for notational simplicity we drop the dependence in $X$.

From now on, we~fix $g\in \set{g_1, g_2, g_3}$; its invariant fibration $\pi\colon X\to \P^1$ is the restriction of one of the projections $\pi_i$ to $X$; its base does not depend on $X$.

\begin{lem}\label{lem:NT_local}
Let $X_0\in \cW_0$, $w_0\in \NT_g$, and $k$ be the multiplicity of $w_0$ in $\NT_g$.
Let $U\subset \P^1$ be a topological disk such that $\overline U\cap \NT_g = \set{w_0}$ and $\overline U\cap \Crit(\pi)=\emptyset$.
Then, there exists a neighborhood $V$
of $X_0$ in $\cW_0$ such that for any $X$ in $V$, the total multiplicity of
$\NT_g$ in $U$ is equal to $k$.
\end{lem}

\begin{proof}
Fix an open connected neighborhood $V$ of $X_0$ such that for $X$ in $V$,
\begin{itemize}
\item $U$ does not intersect any of the sets $\Crit(\pi)$;
\item there is a section $w\mto \varsigma_X(w)$
of $\cX\to \cW_0\times \P^1$ above $V\times U$, together with a continuous choice of basis for the homology of the fibers of $\pi$ above $U$.
\end{itemize}
Then the sections, the Betti foliations (above $U$), and their lifts to $U\times \C$ all
depend continuously on $X$ in~$V$. In~particular, we~can find a disk $U'\subset \P^1$, with
$w_0\in U'\Subset U$, whose boundary is a
smooth Jordan curve $\gamma$, and such that
for any $X\in V$, the Betti foliation is transverse to $g\circ \varsigma_X$ above $\gamma$. In~particular,
$\NT_g$ is disjoint from $\gamma$.

Now, recall that the map $T$ defined in Equation~\eqref{eq:definition_T}
behaves topologically like $w\mto w^{k+1}$; in such a local coordinate,
$k+1$ is the winding number of the curve $T\circ \gamma$ around~$T(0)$.
Since $\NT_g$ stays disjoint from $\gamma$ for $X$ in $V$,
this winding number is constant in $V$; thus, the number of points of $\NT_g$ enclosed by $\gamma$ (counted with multiplicity) stays constant on $V$. The
lemma follows.
\end{proof}

\begin{lem}\label{lem:effective_analytic}
There exists a proper
semi-algebraic subvariety $\mathcal Z_g \subset \cW_0$ of positive codimension
such that $\mult(\NT_g)$
is locally constant in $ \cW_0\setminus \mathcal Z_g$.
\end{lem}

\skpt
\begin{proof}
\subsubsection*{Step 1:  Keeping away from the singular fibers}
Fix $X_0\in \cW_0$ and $w_0\in \P^1$
a critical value of $\pi$.
It is shown in \cite[Lem.\,3.11]{invariant} that ${\NT}_{g}$ does not accumulate $w_0$.
Here, we~show that outside a semi-algebraic subvariety $\mathcal Z_g\subset \cW_0$ this non-accumulation holds uniformly with respect to~$X$: we shall construct a neighborhood $V\times U$ of $(X_0, w_0)$ such that $U$ is disjoint from $\mathrm{NT}_{g}$ for every $X$ in $V$.
For this, we~review the proof of \cite[Lem.\,3.11]{invariant}
and make it locally uniform in $X$ under appropriate hypotheses on $X_0$.

Define $\cW_1$ to be the dense, Zariski open subset\footnote{To show that $\cW_1$ is dense, we~only have to show that it is non-empty.
This is a consequence of the following fact. {\textit{Let $X$ be in $\cW_0$, let $\pi_1\colon X\to \P^1$ be the first projection, and let $m$ be a critical point of $\pi_1$. Let $F$ be the fiber of $\pi_1$ containing $m$. Then, each of the conditions
\begin{enumerate}
\item\label{cond:1} the singularity of $F$ at $m$ is degenerate (in the sense of Morse, \ie it is not an $A_1$-singularity);
\item\label{cond:2} $F$ contains a second singular point $m'$
\end{enumerate}
defines a proper subset of $\cW_0$.}} In other words, these properties \eqref{cond:1} and \eqref{cond:2} disappear after a generic small perturbation of $X$ in $\cW_0$, which can be checked directly.} of $ \cW_1$
such that for any $X\in \cW_1$ and any
$i\in \set{1,2, 3}$, all singular fibers of $\pi_i$ are of type $I_1$. In~this case
there are $24$ such fibers (the Euler characteristic of a K3 surface is $24$, the contribution to the Euler characteristic of a smooth fiber is $0$, and the contribution of an $I_1$ fiber is $1$). Suppose that $X_0\in \cW_1$.

Fix a small disk $U\subset \P^1$ centered at $w_0$ and containing no other singular value of $\pi\colon X_0\to \P^1$.
Fix a neighborhood $V$ of $X_0$ in $ \cW_1$, and local coordinates on $U$ (depending on $X$), so that (i) this property persists for $X\in V$ and (ii) the unique singular value of $\pi$ in $U$ is $w_0=0$. Let
$X^\#_U $ be the complement in $X_U^g:=\pi\inv(U)$ of the unique singular point of $X_{w_0}^g$. We~fix a reference section $\varsigma_X:U\to X^\#_U$ depending holomorphically on~$X\in V$ and~$w\in U$.

For $X\in V$ and $w\in U\setminus\{w_0\}$ we can write
$X_w^g\simeq \C/\Z\oplus \Z\tau_X(w)$, as in Section~\ref{subs:halphen}. Since the singular fiber $X_{w_0}^g$ is of type $I_1$ and $w_0=0$, the monodromy along a simple loop around~$0$ maps the basis $(1, \tau_X(w))$ to $(1, \tau_X(w)+1)$. Moreover,
$X^\#_U$ is biholomorphic to the quotient of $U\times \C$ by the family of lattices
$\Z\oplus \Z\tau_X(w)$, where $\tau_X(w) = \unsur{2\ii \pi} \log (k_X(w))$ for a function $k_X\colon U\to \C$ which has a single zero at the origin and depends holomorphically on $X\in V$ and $w\in U$.
Since
$g\circ \varsigma_X$ is another section of $\pi$ above $U$, there is a holomorphic function $t_X(w)$
of $X$ and $w$ such that the lift of $g$ to $U\times \C$ is given by $(w, z)\mto (w, z+ t_X(w))$.
The calculations of \cite{invariant}
(see \S 3.3.2 and Lem.\,3.11 there) show that the equation for $\mathrm{NT}_g$ in $U$ is
\begin{equation}\label{eq:NT}
- \ii \log(\abs{k_X(w)}) k_X(w) \, t_X'(w) = k_X'(w) \Ima(t_X(w)).
\end{equation}
We claim that if $\Ima(t_X(0)) \neq 0$, then by reducing $V$ and $U$ if necessary,
$\mathrm{NT}_g\cap X_U^g = \emptyset$. Indeed if $U$ is small enough,
there exist positive constants $\e$, $c$ such that for any $X\in V$,
\begin{equation*}
\abs{k_X(w)}\leq \e, \;\, \abs{k_X'(w)}\geq c, \;\,
\abs{\Ima(t_X(w))}\geq c, \;\, {\text{and}} \;\, \abs{t_X'(w)}\leq c\inv.
\end{equation*}
Reducing $U$ further,
$\e$ can be chosen arbitrary small while
$c$ remains bounded away from~$0$. If $\e\log\e < c^3$, this is not
compatible with the equality~\eqref{eq:NT}, so $\mathrm{NT}_g\cap U = \emptyset$.

\begin{lem}\label{lem:twisting_singular}
The locus
\begin{equation}\label{eq:analytic_equation}
\set{X\in V\sep \Ima(t_X(w_0)) = 0}
\end{equation}
is a semi-algebraic subset of positive codimension.
\end{lem}

\begin{proof}[Proof of Lemma~\ref{lem:twisting_singular}]
Consider the Wehler surfaces $X\subset V\subset \cW_1$, and their equations
\begin{equation*}
A_{222} x^2y^2z^2+ A_{221}x^2y^2z+\cdots + A_{100}x + A_{010}y+A_{001}z + A_{000} = 0.
\end{equation*}
Permuting coordinates if necessary, we~suppose that $\pi\colon X\to \P^1$ is the projection onto the first coordinate.
As $X$ varies near $X_0$, the critical value of $\pi$ near $w_0$ and the corresponding critical point in $X$ can be computed algebraically in terms of the $A_{ijk}$. Using the action of $\PGL(2, \C)^3$ on $\P^1\times \P^1\times \P^1$, we~may
assume that $w_0=0$ (as above) and the unique singular point of the fiber $X_{w_0}^g:=X \cap \set{x= 0}$ is~$(0,0)$.
So, the equation of $X_{w_0}^g$ in $\P^1\times \P^1$ is
\begin{equation*}
ay^2z^2 + b y^2z+cy z^2+ dyz+ey^2+f z^2 = 0,
\end{equation*}
for some coefficients $a, \ldots , f$ given by algebraic expressions in the $A_{ijk}$.
Since $X\in \cW_1$, $X_{w_0}^g$ has two transverse branches at $(0,0)$: their tangent directions are given
by the solutions of $dyz+ey^2+f z^2 = 0$ in~$\P^1$.

One can also write $X_0^g\setminus\{0,0\}$ as the quotient of $\set{0}\times \C$ by the lattice $\Lat(0)=\Z$; in this coordinate,
$g$ acts as multiplication by $\exp(2\ii \pi t_X(0))$. Thus, $\Ima(t_X(0)) = 0$ means that $g$ induces a rotation, instead of a loxodromic homography, on the rational curve $X_0^g$. Writing down $g = \sigma_y\circ \sigma_z$ in coordinates, we~obtain
\begin{equation*}\arraycolsep4.5pt
g(y,z) = \begin{pmatrix}-1 - \sfrac{d^2}{ef}& \sfrac{d}{e} \\ -\sfrac{d}{f} & -1
\end{pmatrix}\begin{pmatrix} y\\z\end{pmatrix} + O\blrpar{\norm{(y,z)}^2}
\end{equation*}
for $(y,z)\in X_0^g$.
Thus, $D_{(0, 0)}g\in \GL(T_{0,0}X)$ has determinant $1$ and trace $ - 2- \sfrac{d^2}{ef}$.
As a consequence, $g$ acts as a rotation on $X_0^g$ if and only if $ 2 + \sfrac{d^2}{ef} \in[-2, 2]$: this is a semi-algebraic condition.
\end{proof}

To conclude, we~let $ \cW_g\subset \cW_0$ be
the intersection of $ \cW_1$ with the complement of the subsets defined by
\begin{equation*}
\set{X\in \cW_1\sep \Ima(t_X(w_i) = 0}
\end{equation*} for each of the 24 singular values $w_i$ of $\pi$. We~finally define $\mathcal Z_g$ to be the complement of $ \cW_g$; by Lemma~\ref{lem:twisting_singular}, it is a proper
semi-algebraic set of positive codimension.

\subsubsection*{Step 2: Conclusion}
Pick $X_0\in \cW_0\setminus \mathcal Z_g$ and cover $\P^1$ by a finite family $F$ of
topological disks, such that for every $U\in F$,
$\overline U$ contains at most one point of $\Crit(\pi)\cup \NT_g$.
If $U\in F$ contains a critical value of $\pi$ (and no point of $\NT_g$),
then, as already explained, this property persists in a neighborhood of $X_0$. By~Step~1, for
$X$ sufficiently
close to $X_0$, $\overline U$ is disjoint from $\NT_g$ as well. For the remaining disks, the local constancy of $\mult(\NT_g)$ follows from Lemma~\ref{lem:NT_local}. The proof is complete.
\end{proof}

\begin{proof}[Proof of Proposition~\ref{pro:NT_bounded}]
We use a semi-continuity argument.
Since the exceptional set $\mathcal Z_g$ defined in Lemma~\ref{lem:effective_analytic} is semi-algebraic,
the open set $\cW_0\setminus \mathcal Z_g$ is also semi-algebraic, so it
admits finitely many connected components (see~\cite[Cor.\,2.7]{bierstone-milman:ihes} for instance).
Thus, by~Lemma~\ref{lem:effective_analytic}, $\mult(\NT_g)$ and therefore $\abs{\NT_g}$
are uniformly bounded
on $\cW_0\setminus \mathcal Z_g$, say $\abs{\NT_g}\leq B$. Now, pick $X_0\in \mathcal Z_g$ (thus $X_0\in \cW_0$) and assume
that for $X_0$ one has $\abs{\NT_g}>B$. We~can then consider a finite number of small topological
disks $U_i$ with disjoint closures in $\P^1$, such that $\abs{\NT_g}\cap \bigcup U_i>B$. By~Lemma~\ref{lem:NT_local}, these non-twisting points persist for
$X$ close enough to $X_0$. Since $\cW_0\setminus \mathcal Z_g$ is dense in~$\cW_0$, this
contradicts the definition of $B$ and the proof is complete.
\end{proof}

\begin{proof}[Proof of Theorem~\ref{thm:wehler_zariski}]
The main point of~\cite[Th.\,A]{finite_orbits} is
that the set of $X\in \cW_0$ possessing a finite orbit of length $\leq B$ is a
proper Zariski closed subset $\mathcal Z_B$ of $ \cW_0$. For~$N$ as in
Proposition~\ref{pro:N(X)_bounded}, for
any $X\in \cW_0\setminus \mathcal Z_N$, all finite orbits of $\Gamma$ are uniformly expanding. We
conclude by applying Theorem~\ref{thm:criterion_uniform_expansion_parabolic} (with $\nu = \unsur{3}(\delta_{\sigma_1}+ \delta_{\sigma_2}+\delta_{\sigma_3})$). The proof of the corresponding statement in $ \cW_0(\R)$ is identical.
\end{proof}

\begin{rem}
We expect that an analogue of Theorem~\ref{thm:wehler_zariski} holds for other families
with large automorphism groups containing parabolic elements, like Enriques surfaces, or the family associated to pentagon folding (see~\cite{Simons}).
\end{rem}

\begin{rem}
The proof of Proposition~\ref{pro:NT_bounded} suggests that
there should exist a notion of multiplicity, {\textit{including singular fibers}},
for which $\mult(\NT_g)$ would be constant on $\cW_0$ and would be an algebraically
computable invariant of the parabolic automorphism $g$. A variant of this question is mentioned in \cite[Rem.\,7.7.4]{duistermaat}.
\end{rem}

\subsection{Thin subgroups}\label{subs:thin}
In this section we consider the total
space $ \cW$ of all Wehler surfaces and the universal family
$\mathcal X\subset \cW\times (\pu\times \pu\times \pu)$. We~change a
little bit the notation: $\Gamma$ will be a subgroup of $\Z/2\Z\ast\Z/2\Z\ast\Z/2\Z$, and
$\Gamma_{X}$ will be the corresponding subgroup of $\Aut(X)$.

Let $E$ be an elliptic curve. Consider the following classical Kummer construction (see~\cite[\S 4]{finite_orbits}): let $\eta$ be
the involution $\eta (x,y)\mto (-x,-y)$ on $A:=E\times E$; the associated Kummer surface is
the desingularization $\widehat{E\times E/\eta}$; the natural $\GL(2, \Z)$ action on $E\times E$ descends
to $E\times E/\eta$ and
induces a non-elementary automorphism group of $\widehat{E\times E/\eta}$.
The surface $E\times E/\eta$ can be realized as a singular Wehler example
(see~\cite[\S 8.2]{Cantat:Panorama-Synthese}); in addition the action of
$\Z/2\Z\ast\Z/2\Z\ast\Z/2\Z$ is induced by a finite index subgroup of $\GL(2, \Z)$.
Let us briefly recall the construction: write
$E$ in Weierstrass form $y^2 = 4x^3 - g_2x - g_3$, with the neutral element of the group law on
$E$ located at infinity.
To $(m_1, m_2)\in E\times E$, $m_i = (x_i, y_i)$, we~associate $m_3 = - (m_1+m_2)$ and
$\phi(m_1, m_2 ) = (x_1, x_2, x_3)$, where $m_3 = (x_3, y_3)$. Then, $\phi$ is $\eta$-invariant and determines
a biregular map
$\phi: E\times E/\eta\to X_E$ onto a singular Wehler surface $X_E$ with $16$ nodal singularities.

Assume that $\Gamma\subset \Z/2\Z\ast\Z/2\Z\ast\Z/2\Z$ is not virtually cyclic. Then for $X\in \cW_0$, $\Gamma_X$ is non-elementary (see~\cite[\S 3]{stiffness}).

\begin{thm}\label{thm:thin}
Let $\Gamma$ be a subgroup of $\Z/2\Z\ast\Z/2\Z\ast\Z/2\Z$ which is not virtually cyclic.
For $X\in \cW_0$ sufficiently close to $X_E$, the subgroup $\Gamma_X$ is uniformly expanding on $X$.
\end{thm}

Thus for every ``abstract'' non-elementary subgroup $\Gamma$ of $\Z/2\Z\ast\Z/2\Z\ast\Z/2\Z$, the open
subset $\cW_{\exp}(\Gamma)$ of those $X\in \cW_0$ for which the action of $\Gamma_X$ is uniformly expanding is non-empty.
The group $\Gamma$ can be arbitrarily thin, in particular it is not assumed to contain parabolic elements. In~view of Theorem~\ref{thm:criterion_uniform_expansion}, it is natural to expect that $\cW_{\exp}(\Gamma)$ is actually dense in the Euclidean topology.

\begin{proof}
The difficulty is that we cannot directly argue that uniform expansion is an open property, because
$X_E$ is singular.

\begin{lem}\label{lem:regular}
Fix $f\in \Z/2\Z\ast\Z/2\Z\ast\Z/2\Z$, and denote also by $f$ the induced fibered map on
the universal family of $(2,2,2)$-surfaces in $(\P^1)^3$.
Then $f$ is regular on a neighborhood of $X_E$.
\end{lem}

\begin{proof}[Proof of Lemma~\ref{lem:regular}] Pick a $(2,2,2)$ surface $X$. If $X$ does not contain any fiber
of the projection $\pi_{12}=(\pi_1,\pi_2)\colon (\P^1)^3\to (\P^1)^2$, then the same property holds in a neighborhood $\mathcal V$ of $X$ in the universal family of $(2,2,2)$-surfaces; furthermore, $\sigma_3$~determines an automorphism\footnote{Indeed,
denote by
${\mathcal{V}}_0$ the projection of $\mathcal V$ in the space of Wehler surfaces $ \cW_0$ (see Section \ref{subs:wehler}),
and let $o\in \cW_0$ be the image of $X$.
For $v\in {\mathcal{V}}_0$, let $X_v \subset {\mathcal V}$ be the corresponding Wehler surface.
Pick a point $(x,y,z)\in X$, say in an affine chart of $(\P^1)^3$ where none of the coordinates is $\infty$. In~other words, $z=[z_0:z_1]\in \P^1$ with $z_1\neq 0$ and $z=z_0/z_1$, and similarly for $x$ and $y$.
For simplicity, we~use affine coordinates $\bfx$, $\bfy$ for the first two coordinates and homogeneous coordinates $[\bfz_0:\bfz_1]$ for the third one. The equation for $X_v\subset \mathcal{V}$ can be written $A_v (\bfx,\bfy)\bfz_0^2+B_v(\bfx,\bfy)\bfz_0\bfz_1+C_v(\bfx,\bfy)\bfz_1^2=0$ and then $\sigma_3(\bfx, \bfy, [\bfz_0:\bfz_1])= (\bfx, \bfy, [-B_v(\bfx,\bfy) \bfz_1-A_v (\bfx, \bfy)\bfz_0 : A_v(\bfx,\bfy)\bfz_1])$. We~have to show that this map is regular near $(x,y,z)$. Since we work near a point at which $z_1\neq 0$, the only problem could be that, at some point $(x',y',[z'_0:z'_1])$ of some $X_v$, we~have both $A_v(x',y')=0$ and $B_v(x',y')=0$, but then $C(x',y')$ should be $0$ too, and $X_v$ should contain a vertical line.} of $\mathcal V$.
Thus, we~only have to prove that $X_E$ does not contain any fiber of the projections $\pi_{ij}$.
Let us show that $X_E$ does not contain any vertical line $\{x=x_0,\, y=y_0\}$. Such a line would provide a family of points $(m_1,m_2)$ on $E\times E$ with fixed first coordinates $x_1=x_0$, $x_2=y_0$, for which the first coordinate of $m_3:=-(m_1+m_2)$ takes arbitrary values. This is impossible. The same argument applies to the lines $\{y=y_0,\, z=z_0\}$ and $\{z=z_0,\, x=x_0\}$ because the relation $m_1+m_2+m_3=0$ is symmetric (equivalently, the equation of $X_E$ given in~\cite[\S 8.2]{Cantat:Panorama-Synthese} is symmetric in $(x,y,z)$).
\end{proof}

There is a finite index subgroup of $\Gamma$ that
fixes each singularity of $X_E$. By~Proposition~\ref{pro:induced_UE2} and the fact that uniform expansion does not depend on the measure, we~can replace $\Gamma$ with this finite index subgroup, endow $\Gamma$ with a finitely supported, symmetric measure $\nu$ with $\Gamma=\Gamma_\nu$, and then we have to prove that $(\Gamma_X, \nu_X)$ is uniformly expanding for $X\in \cW_0$ near $X_E$; here, $\nu_X$ is the measure induced by $\nu$ on $\Gamma_X$.

Endow $\P^1\times \P^1\times \P^1$ with the Fubini study metric, and the Wehler surfaces $X$ with the induced metric. Recall that $T^1X$ denotes the unit tangent bundle.

Assume, by~way of contradiction, that there is a sequence
$X_n\to X_E$ along which~$\nu_{X_n}$ is not uniformly expanding.
For each $n$, let $\varpi$ denote the natural projection $T^1X_n\to X_n$ (\resp $T^1X_E\to X_E$).
Denote by $T^1X_E$ the subset of $T^1(\pu\times \pu\times \pu)$ which coincides with $T^1 \mathrm{Reg}(X_E)$ above the regular part of $X_E$ and coincides with $T^1_x(\pu\times \pu\times \pu)$ above each singularity $x\in \Sing(X_E)$. With this definition it is obvious that
if $x_k$ is any sequence such that $x_k\in \bigcup X_n\cup X_E$ and $x_n\to x\in X_E$
then we have the semicontinuity
$ \limsup T^1_{x_n}X_n\subset T^1X_E$. Theorem~\ref{thm:limit_measure_top_lyapunov}
provides a sequence of stationary measures $\hat \mu_{X_n}$ on $T^1X_n$ (with projections $\mu_{X_n}:=\varpi_*\hat\mu_{X_n}$) such that
\begin{equation*}
\int \log\norm{f_\varstar u} d\nu_{X_n}(f) \, d\hat \mu_{X_n} (u)\leq 0.
\end{equation*}

From Lemma~\ref{lem:regular}, we~can extract a subsequence, still denoted by $(X_n)$,
such that $\lrpar{\hat \mu_{t_n}}$ converges to a stationary measure $\hat \mu_{X_E}$ on $T^1X_E$ satisfying
\begin{equation*}
\int \log\norm{f_\varstar u} d\nu_{X_E}(f) \, d\hat \mu_{X_E} (u)\leq 0.
\end{equation*}
By iterating and using the stationarity of $\hat \mu_{X_E}$, the same inequality holds with $\nu_{X_E}^{(m)}$ instead of $\nu_{X_E}$ for every positive integer $m$. To get the desired contradiction, we~shall show that no such measure exists.

\subsubsection*{Step 1: near the singularities}
Here we show that there exists $n_0\in \N$, $c_0>0$,
and an open neighborhood $U$ of $\Sing (X_E)$ such that if $u\in T^1X_E$ and
$\varpi(u)\in U$, then
\begin{equation}\label{eq:UE_sing}
\int \log\norm{f_\varstar u}\, d\nu_{X_E}^{(n_0)}(f) \geq c_0.
\end{equation}
By Lemma~\ref{lem:regular} and the above mentioned semicontinuity of unit tangent
bundles it is enough to prove this when
$x=\varpi(u)\in \Sing(X_E)$.
Recall that $\Gamma_{X_E}$ fixes $\Sing(X_E)$ pointwise. Around each of its singularities, $X_E$ is locally isomorphic to the quotient
$\C^2/\eta$, $\eta(u,v)=(-u,-v)$, embedded in the standard way in $\C^3$ by
\begin{equation*}
\phi: (u,v)\mto (u^2, uv, v^2) = (\xi,\eta,\zeta),
\end{equation*}
whose image is the quadratic cone $\{\xi\zeta-\eta^2 = 0\}\subset \C^3$. The level-$2$ congruence
subgroup $G$ of $\GL_2(\Z)$
fixes each torsion point of $A:=E\times E$ of order $\leq 2$, and $\Gamma_{X_E}$ is induced by
a non-elementary subgroup $G_0$ of $G$.
The standard linear action of $G$ on $\C^2$ (or more precisely on a neighborhood of any $2$-torsion point of $A$) commutes to $\eta$ and induces a linear action on $\C^3$ via the homomorphism
\begin{equation*}\arraycolsep4.5pt
\phi_\varstar: \begin{pmatrix} a&b\\ c&d\end{pmatrix} \longmapsto
\begin{pmatrix} a^2 & 2ab &b ^2 \\ ac & ad+bc & bc \\ c^2 & 2cd & d^2\end{pmatrix}.
\end{equation*}
Thus, the action of $\Gamma_{X_E}$ on the tangent cone
of $X_E$ at the origin (which is naturally identified with $\{\xi\zeta-\eta^2 = 0\}$)
is, up to a linear conjugacy, given by $\phi_\varstar(G_0)$. Since this cone is Zariski-dense in
$T_xX_E = T_x((\P^1)^3)$ and the action of $\Gamma$ on the universal family of $(2,2,2)$ surfaces is smooth at $(X_E, x)$, we~deduce that the action of $\Gamma$ on $T^1_xX_E$ is also induced
by $\phi_\varstar(G_0)$.

This is a subgroup
of $O(q;\R)\simeq O_{2,1}(\R)$, where $q$ is the quadratic form $q(x,y,z) = xz-y^2$. By~assumption, it is a non-elementary group of isometries of $q$, hence it acts strongly irreducibly and proximally on $\R^3\subset \C^3$ (loxodromic elements of $\GL_2(\Z)$ are mapped to loxodromic elements in $O(q;\R)$). It preserves the real decomposition $ \C^3=\R^3\oplus_\R \ii \R^3$ and the action on $\R^3$ and $\ii \R^3$ are linearly conjugate (by multiplication by $\ii$). Therefore, as in Section \ref{subs:finite_orbits}, the Inequality~\eqref{eq:UE_sing} follows from~\cite{furstenberg-kifer} (see also~\cite[Chap.\,III, Cor.\,3.4(iii)]{bougerol-lacroix}).

\subsubsection*{Step 2: away from the singularities}
We shall show that there exists a neighborhood $U'\subset U$ of $\Sing(X_E)$
and $c>0$ such that
for any fixed $u\in T^1X_E$ such that $\varpi(u)\notin U'$,
\begin{equation}\label{eq:UE_proba}
\pp\Blrpar{\unsur{m}\log \norm{(f_\omega^m)_\varstar u} \geq c} \underset{m\to\infty}{\to} 1.
\end{equation}
By Lemma~\ref{lem:weak_UE2} (see also Remark~\ref{rem:given_v}), this implies that
$\ee\lrpar{\log \norm{(f_\omega^m)_\varstar u}}\geq mc/2$ for large $m$. Then, the first step and
a compactness argument identical to that of
Lemma~\ref{lem:weak_UE} show that uniform expansion holds on $T^1X_E$,
which is the desired contradiction.

Let $U'$ be an open neighborhood of $\Sing(X_E)$ which will be specified later.
There is a constant $\delta = \delta(U')$ such that
\begin{equation}\label{eq:flat1}
\text{if } \varpi(u)\notin U'\text{ and } \varpi(f_\varstar u)\notin U', \text{ then }
\log\norm{f_\varstar u} \geq \log\norm{f_\varstar u}_{\mathrm{flat}} - \delta,
\end{equation}
where $\norm{\cdot}_{\mathrm{flat}}$ is the Riemannian metric on $\mathrm{Reg}(X_E)$ induced by the flat metric of
$E\times E$.

The pull-back of
$\nu$ to $\GL(2, \Z)$ generates $G_0$ and its support is finite. Since $G_0\subset \GL(2, \Z)$ is non-elementary, we~have uniform expansion with respect to the flat metric. By~Lemma~\ref{lem:weak_UE2},
there exists a constant $c_1>0$ and sets of trajectories
$\Omega_m^1\subset \Omega$ such that $\pp(\Omega_m^1)\to 1$ as $m\to \infty$
and
\begin{equation}\label{eq:flat2}
\text{if }\omega\in \Omega_m^1, \quad
\unsur{m}\log \norm{(f_\omega^m)_\varstar u}_{\mathrm{flat}} \geq c_1 .
\end{equation}
Fix $\e>0$ and $c$ such that $0<c<c_1$. We~claim that there is a Margulis function on $X_E$ with poles at
$\Sing(X_E)$.
Indeed, to construct it we view $X_E$ as $E\times E/ \eta$. Let
$\pi:E\times E\to E\times E/\eta \simeq X_E$ be the natural map.
On
$E\times E$, $\pi\inv(\Sing (X_E)) =\mathrm{Fix}(\eta)$ is invariant under the $\GL_2(\Z)$-action, which
is uniformly expanding, so $ u :x\mto - \log d_{\mathrm{flat}}(x, \mathrm{Fix}(\eta))$ defines a
Margulis function. Since $\eta$ is an isometry for the flat metric, $u$ is
$\eta$-invariant so it descends to a Margulis function on $X_E$, as asserted. We~deduce that there is an open set $U' = U'(\e)\subset U$ with the following property:
for large enough $m$,
the set $\Omega_m^2$ of trajectories $\omega\in \Omega$ such that
$(f_\omega^m)(\varpi(u)) \notin U'$ satisfies $\pp(\Omega_m^2)\geq 1-\e/2$. Now,
$U'$ being fixed, for $m$ large enough we have that $\pp(\Omega_m^1\cap \Omega_m^2)\geq 1-\e$ and
by~\eqref{eq:flat1} and~\eqref{eq:flat2}, if $\omega \in \Omega_m^1\cap \Omega_m^2$,
\begin{equation*}
\unsur{m}\log \norm{(f_\omega^m)_\varstar u}\geq c_1 - \frac{\delta(U')}{m}\geq c.
\end{equation*}
Thus, the convergence~\eqref{eq:UE_proba} holds and the proof is complete.
\end{proof}

\section{Applications}\label{sec:applications}

\subsection{Orbit closures}
The following is a version of the orbit closure Theorem~E of~\cite{invariant} in which
periodic orbits are allowed.
Combined with Theorem~\ref{thm:wehler_zariski}, it gives
Theorem~\ref{thm:orbits_wehler}.

\begin{thm}\label{thm:orbit_closures}
Let $X$ be a torus, a K3 surface, or an Enriques surface.
Let $\Gamma\subset \Aut(X)$ be a non-elementary subgroup which contains parabolic elements and does not
preserve any algebraic curve. Assume
that for any finite orbit $O$, the induced action of $\Gamma$ on $TX\rest{O}$ is non-elementary.
Then there exists a finite set $F$ and a real analytic,
totally real, and $\Gamma$-invariant surface $Y \subset X$ with $\Sing(Y)\subset F$
such that for every $x\in X$:
\begin{enumeratea}
\item\label{thm:orbit_closuresa} either $x$ belongs to $F$ (and its orbit is finite);
\item\label{thm:orbit_closuresb} or $x$ belongs to $Y\setminus F$ and $\overline{ \Gamma (x)}$ is a union of components of $Y$;
\item\label{thm:orbit_closuresc} or $\overline{ \Gamma (x)} =X$.
\end{enumeratea}
\end{thm}

\begin{proof}
First observe that under these assumptions, \cite[Th.\,C]{finite_orbits} implies
that there exists a maximal finite invariant subset $F$.
Fix a symmetric measure $\nu$ such that $\Gamma_\nu = \Gamma$ and satisfying the
moment condition~\eqref{eq:moment+}. By~Theorems~\ref{thm:criterion_uniform_expansion_parabolic} and~\ref{thm:finite_UE},
$\nu$ is uniformly expanding. We~now resume the discussion from~\cite[\S 8]{invariant}, in particular Remark~8.6 there: $\STang_\Gamma$ is a finite invariant set and if $x\in X$ is such that $\Gamma(x)$ is infinite
but not dense, then there are two possible situations:
\begin{enumerate}
\item\label{case:1} either $\overline{\Gamma(x)}\setminus \STang_\Gamma$ is discrete outside $\STang_\Gamma$;
\item\label{case:2} or $\overline{\Gamma(x)}\setminus \STang_\Gamma =: Y(x)$ is a totally real analytic surface, whose singular locus $\Sing(Y)$ is discrete outside $\STang_\Gamma$.
\end{enumerate}

In case~\eqref{case:1}, $\Gamma(x)$ is finite.
Indeed $\overline {\Gamma(x)}$ is at most countable,
so if $\mu$ is any cluster value of
$\unsur{n}\sum_{k=0}^{n-1} \nu^k\star\delta_x$, then $\mu$ is a purely atomic
stationary measure. In~this case it follows from Theorem~\ref{thm:margulis_UE} that
the orbit of $x$ must be finite, hence contained in $F$.

If case~\eqref{case:2} holds, we~first claim that $\Sing(Y(x))$ is finite.
Indeed, $\overline {\Sing(Y(x))}$ is a $\Gamma$-invariant countable set, which clusters only at
$\STang_\Gamma$. By~the previous argument, every orbit $\Gamma(y)$ in $\overline {\Sing(Y(x))}$ is finite, so by the
finiteness of the set of finite orbits~\cite[Th.\,C]{finite_orbits} we conclude that
$\Sing(Y(x))$ itself is finite.

Now, let $\mu'$ be a cluster value of
$\unsur{n}\sum_{k=0}^{n-1} \nu^k\star\delta_x$. By~Theorem~\ref{thm:margulis_UE}, $\mu'$
is an atomless stationary measure
supported on $Y(x)$ such that $\mu'(\mathrm{Reg}(Y(x)))=1$. Since $\Gamma$ has no invariant
curve, $\mu'$ is Zariski diffuse. Let $\mu$ be any ergodic component of $\mu'$.
Theorems~\ref{thm:rigidity} and~\ref{thm:hyperbolic} imply that $\mu$ is hyperbolic and its stable directions depend genuinely on the itinerary. Then the
argument of \cite[Th.\,3.1]{br} adapts immediately to show that $\mu$ is
SRB.\footnote{We are \emph{not} claiming that we can extend \cite{br} to non-compact surfaces here. All the necessary
estimates on the Lyapunov norms and Pesin charts
hold by viewing $\mu$ as a hyperbolic stationary measure on the
\emph{compact} complex manifold $X$. The only issue appears when considering the size and intersection
properties of real stable and unstable manifolds in $Y(x)$, starting from \S 9.7 of \cite{br}.
At~this stage Brown and Rodriguez-Hertz
already discard a set of small measure of points with bad properties (see the definition of $\Lambda(\gamma_1)$ on p. 1087); so it is enough to remove from this $\Lambda(\gamma_1)$ the set of
small measure of points too close to $\Sing(Y(x))$, and proceed with their argument.} The
canonical invariant 2-form of $X$ induces a $\Gamma$-invariant measure $\vol_{Y(x)}$
on $Y(x)$ (see Lemma~\ref{lem:volume_Y}). Since $\mathrm{Reg}(Y(x))$ admits a Margulis function, we~conclude from
Proposition~\ref{pro:finite_mass} that the volume $\vol_{Y(x)}$ is finite. Therefore we can copy verbatim the
argument of \cite[Th.\,3.4]{br} to conclude that $\mu$ is $\Gamma$-invariant.
Since \cite[Th.\,C]{invariant} says that there are only finitely many $\Gamma$-invariant measures,
there are only finitely many possible surfaces $Y(x)$. Taking $Y$ to be their union,
the proof is complete.
\end{proof}

\subsection{Ergodicity} \label{subs:ergodicity}
In~\cite{dolgopyat-krikorian}, the original motivation to introduce uniform
expansion was a criterion for ergodicity. The same holds in our setting, with a few caveats which will be explained below.

\begin{thm}[{Dolgopyat-Krikorian~\cite[Cor.\,2]{dolgopyat-krikorian}}, see also~\cite{liu_thesis, chung}]
\label{thm:dolgopyat-krikorian}
Let $X$ be a torus, a K3 surface, or an Enriques surface. Let $\Gamma\subset \Aut(X)$ be a
non-elementary subgroup with a uniformly expanding action on $X$. Then $\vol_X$ is $\Gamma$-ergodic.

Likewise, if $Y\subset X$ is a $\Gamma$-invariant totally real
analytic subset such that $\Gamma$ acts transitively on the
set of irreducible components of $Y$, then $\vol_Y$ is $\Gamma$-ergodic.
\end{thm}

Note that the notion of irreducible component in real analytic geometry is not well-behaved in general
(see~\cite[\S 5.1]{invariant} for a short discussion).
Here we content ourselves with saying that $Y$ is irreducible when $\mathrm{Reg}(Y)$ is connected.
Observe also that the ergodicity of
$\vol_X$ follows directly from Theorem~\ref{thm:classification_invariant} when $\Gamma$
contains a parabolic element.

\begin{proof}[Proof (sketch)]
The proof in~\cite{dolgopyat-krikorian} is a bit sketchy, but it was already
expanded in~\cite{liu_thesis, chung} (see also~\cite{zzhang}).
Here we just make a few comments on (1) the extension to the holomorphic case for the action on $X$,
and (2) how to deal with the possible singularities for the action on totally real surfaces $Y$.

Regarding the action on $X$, let us recall that the proof of~\cite{dolgopyat-krikorian} is a
variation on the Hopf argument in which the asymptotic behavior of the Birkhoff
sums $ \unsur{n}\sum_{k=0}^{n-1} \delta_{f_\omega^n(x)}$ is propagated along chains of
local stable manifolds (associated to different $\omega$'s), to ultimately conclude that there
is a uniform $r$ so that almost every point $x$
is located at distance at least $r$ from the boundary of its ergodic component.
The key technical ingredients are the facts that under the uniform expansion assumption:
\begin{itemize}
\item stable directions at a given point do not concentrate, more precisely there exists
$\alpha>0$ such that for any $x\in X$ and any $[v]\in \pp(T_x X)$, the probability that
$d_{\P^1}\lrpar{[v], [E^s(x,\omega) ]} <\alpha$ is smaller than 1/100: this follows from a compactness argument
(see~\cite[Prop.\,4.4.4]{liu_thesis});
\item the Pesin local stable manifolds have uniformly bounded geometry
(\eg uniformly bounded size in the sense of
\cite[\S 7.4]{stiffness}): this follows from the usual proof of the local stable manifold theorem;
\item the absolute continuity of the local stable foliation in Pesin charts: we can copy the usual proof
or notice that in the holomorphic case this follows from the fact that the holonomy of a holomorphic
motion is quasiconformal.
\end{itemize}
Given these facts, we~can copy the proof of~\cite{dolgopyat-krikorian} by plugging in \S 10.4
the following elementary geometric property, whose proof is left to the reader:
let $w = (w_1, w_2)\in \C^2$ with $\norm{w}<1$ (possibly close to 1)
and $E_w$ be the direction perpendicular to the line $(0w)$; then if
$L$ is a complex line containing $w$, such that the angle in $\P^1$
between the direction of
$L$ and $E_w$ is greater than $\alpha$, then $L\cap B(0, 1)$ contains a disk of radius~$r\geq r(\alpha)$.

For the second statement of the theorem we can directly resort to~\cite{liu_thesis, chung}, except that
we have to take into account the possibility of singular points on $Y$, which affect the size and geometry of local stable manifolds on $Y$. For this, we~may argue exactly as in Theorem~\ref{thm:orbit_closures}:
first, the existence of a Margulis function guarantees that $\vol_Y$ is finite. Next, since
uniform expansion holds on $X$, the size and angle change of local complex stable manifolds is uniformly controlled. Thus, when restricting to $Y$, we~also have a uniform control of this geometry outside any
$\delta$-neighborhood of $\mathrm{Sing}(Y)$. Since the Hopf argument is local, we~get that there is a single
ergodic component outside a $\delta$-neighborhood of $\mathrm{Sing}(Y)$, for every $\delta>0$, and we conclude by letting $\delta$ tend to zero.
\end{proof}

\begin{rem}
The argument of~\cite{dolgopyat-krikorian} works for a random dynamical system on a (real) compact
manifold of dimension $2d$, enjoying a uniform expansion property
along $d$ dimensional tangent subspaces. This assumption
does \emph{not} hold in our setting since along a totally real subspace one
may witness both expansion and
contraction. In~particular the complex uniform expansion condition is not stable under
$C^1$ perturbations by (real) volume preserving diffeomorphisms of~$X$.
Still, the philosophy of the above proof is that
the argument is robust enough so that uniform expansion along complex 1-dimensional
tangent subspaces in a 2-dimensional complex surface guarantees ergodicity.
\end{rem}

\subsection{Equidistribution}

In the following results, given an action of $(\Gamma, \nu)$ on $M$
we say
that random trajectories from $x$ equidistributes towards
$\mu$ if
$\unsur{n}\sum_{k=1}^n \delta_{f^k_\omega} \to \mu$ for $\nu^\N$-almost every
$\omega$, where the convergence is in the
weak$^\varstar$ topology. By~averaging with respect to $\nu^\N$ and applying
the dominated convergence theorem, this implies that
$\unsur{n}\sum_{k=1}^n \nu^k\ast\delta_x\to\mu$ as well.

The following theorem already appears under stronger moment assumptions in \cite[Th.\,D]{chung}.

\begin{thm}
Let $X_\R$ be a smooth
real projective surface and $\nu$ a probability measure on $\aut(X_\R)$ satisfying \eqref{eq:moment+}.
Assume that $\Gamma_\nu$ preserves a smooth volume form $\vol$ on $X(\R)$ and that $\nu$ is uniformly
expanding on $X(\R)$. Then for any $x\in X$ one of the following alternatives holds:
\begin{enumeratea}
\item either $\Gamma_\nu\cdot x$ is finite;
\item or the random trajectories from $x$ equidistribute
towards $\vol_{X'(\R)}$, the normalized induced volume on a
union of components of $X(\R)$.
\end{enumeratea}
\end{thm}

Recall from Theorem~\ref{thm:criterion_uniform_expansion_parabolic} that the uniform expansion assumption
holds when $\Gamma_\nu$
contains parabolic elements, has no invariant curve, and that the induced
action of~$\Gamma_\nu$ on finite orbits is uniformly expanding. In~this case by \cite[Th.\,C]{finite_orbits},
the number of finite orbits is finite. By~Theorem~\ref{thm:wehler_zariski} this applies to generic real Wehler
surfaces and yields Theorem~\ref{thm:equidistribution_wehler}.

\begin{proof}
Breiman's ergodic theorem says that that for $\nu^\N$-almost every $\omega$, any cluster limit $\mu$
of the sequence of empirical averages $\unsur{n} \sum_{k=1}^n \delta_{f_\omega^k(x)}$ is stationary.
Since $\nu$ is uniformly expanding, the existence of a Margulis function (Theorem~\ref{thm:margulis_UE}) shows that if $\Gamma_\nu\cdot x$ is infinite, $\mu$ gives no mass to finite orbits.
Since $\nu$ is uniformly expanding, any ergodic stationary measure $\mu'$ is hyperbolic and its stable directions are non-random, so by \cite[Th.\,3.4]{br}, $\mu'$ is absolutely continuous with respect to $\vol_{X(\R)}$. The ergodicity
Theorem~\ref{thm:dolgopyat-krikorian} shows that for any component $X_0(\R)$ of $X(\R)$, $\vol_{X_0(\R)}$ is ergodic, so
we conclude that, up to scaling, $\mu$ is a finite combination of measures of this type.
\end{proof}

The next result is conditional to the $\nu$-stiffness property of
complex non-elementary uniformly expanding actions. We~expect that
it will be established in the near future.

\begin{thm}\label{thm:equidistribution_complex}
Let $X$ be a K3 or Enriques surface and $\nu$ be a probability
measure on $\aut(X)$ satisfying~\eqref{eq:moment+}.
Assume that
\begin{enumerate}
\item\label{thm:equidistribution_complex1} $\Gamma_\nu$ is non-elementary, contains parabolic elements, has no invariant curve,
and every finite $\Gamma$-orbit is uniformly expanding;
\item $\nu$-stiffness holds, that is, every $\nu$-stationary measure is invariant;
\item\label{thm:equidistribution_complex2} every compact, real analytic, totally invariant surface $Y\subset X$ is smooth.
\end{enumerate}
Then there exists a finite set $F$ and a (possibly singular) totally real analytic surface~$Y$ such that for every
$x\in X$:
\begin{enumeratea}
\item\label{thm:equidistribution_complexa} either $x$ belongs to $F$;
\item\label{thm:equidistribution_complexb} or $x$ belongs to $Y\setminus F$ and its orbit equidistributes towards $\vol_{Y'}$, where $Y'$ is a union of components of $Y$;
\item\label{thm:equidistribution_complexc} or $x\notin F\cup Y$ and its orbit equidistributes towards $\vol_{X}$.
\end{enumeratea}
\end{thm}

The third hypothesis is a weakness of this statement, since we do not know how to study the singularities of invariant real analytic surfaces (except, of course, when we know how to exclude the existence of finite orbits).

\begin{proof}
The sets $F$ and $Y$ were already constructed in Theorem~\ref{thm:orbit_closures}, whose proof
also implies property~\eqref{thm:equidistribution_complexb}.
The classification of invariant measures (Theorem~\ref{thm:classification_invariant}) and the
stiffness property show that the only $\nu$-stationary measure giving no mass to
$Y\cup F$ is $\vol_X$. Therefore the equidistribution property~\eqref{thm:equidistribution_complexc} follows from Breiman's ergodic theorem
and the existence of a Margulis function associated to finite orbits and totally real surfaces (Theorems~\ref{thm:margulis} and~\ref{thm:margulis_totalement_reel}).
\end{proof}

\appendix
\refstepcounter{section}\label{app:zero_entropy}
\section*{Appendix. Rigidity of zero entropy measures}

We complete the proof of Theorem~\ref{thm:rigidity} with the following result
of independent interest.\vspace*{-5pt}\enlargethispage{.5\baselineskip}%

\begin{thm}\label{thm:entropy_every_f}
Let $X$ be a torus or a K3 or an Enriques surface,
and $\nu$ be a probability measure on
$\Aut(X)$ such that $\Gamma_\nu$ is non-elementary. Assume that $\mu$ is a Zariski diffuse
$\nu$-stationary measure such that $h_\mu(X, \nu) = 0$.
Then $\mu$ is $\Gamma_\nu$-invariant and for every
$f\in \Gamma_\nu$, $h_\mu(f)=0$.
\end{thm}

\begin{proof}
As in Theorem~\ref{thm:rigidity}, we~may assume that $\mu$ is ergodic as a stationary measure, and
its $\Gamma_\nu$-invariance was already established there.
Pick $f\in \Gamma_\nu$.
Assume by way of contradiction that $h_\mu(f)>0$,
in particular $f$ must be loxodromic.
If $\mu$ is ergodic for~$f$, then the result follows rather immediately from the
measure rigidity theorem~11.1 in~\cite{stiffness}. Indeed in that theorem
we consider an ergodic measure $\mu$
of positive entropy for $f$ and study the group of automorphisms of $X$ preserving $\mu$,
under the additional assumption that $\mu$ is supported on a real surface. We~reduce the argument to the case of $\Gamma = \langle f, g\rangle$ for some $g$, and divide the proof into 3 cases: (1)~either there is a $\Gamma$-invariant measurable line field, or (2)
there is a $\Gamma$-invariant pair of measurable line fields, or (3) none of the above. In~cases (1) and (2)
we conclude that $\Gamma$ is elementary by adapting the argument of~\cite[Th.\,9.1]{stiffness}: this
does not rely on the additional real structure. In~case (3), since $\mu$ is hyperbolic for $f$, Theorems~\ref{thm:ledrappier_invariance_principle} and~\ref{thm:classification_proj_invariant} imply that
$\mu$ is hyperbolic as a stationary measure and as in the proof of Theorem~\ref{thm:rigidity} we deduce
that $h_\mu(X, \nu)>0$, which
is contradictory. Thus, case (3) does not happen, and we deduce that $\Gamma$ is elementary for every
$g\in \Gamma_\nu$, which is a contradiction. Therefore~$h_\mu(f) = 0$.

What remains to do is to adapt this argument to the case where $\mu$ is not ergodic under $f$.
So consider $f\in \Gamma_\nu$ and assume that $h_\mu(f)>0$ so that $f$ is loxodromic.
As~before there are 3 cases: either (1) there is a $\Gamma_\nu$-invariant measurable line field, or~(2)
there is a $\Gamma_\nu$-invariant pair of measurable line fields, or (3) none of (1) and~(2). We~first observe that as before case (3) does not happen: indeed if there is no invariant line field or pair of invariant line fields, by~Theorems~\ref{thm:ledrappier_invariance_principle} and~\ref{thm:classification_proj_invariant}, either $\mu$ is hyperbolic as a $\nu$-stationary measure, or the
projectivized tangent action of $\Gamma_\nu$ reduces to a compact subgroup. But since
$h_\mu(f)>0$, $f$ admits non-zero Lyapunov exponents on a set of positive measure
so the latter is impossible. Hence $\mu$ is hyperbolic as a $\nu$\nobreakdash-stationary measure, and since
there is no invariant line field, stable directions depend on the itinerary and as before we conclude that
$h_\mu(X, \nu) > 0$, a contradiction. So one of cases (1) or (2) holds.

So assume there exists a measurable $\Gamma_\nu$-invariant line field
$x\mto [E(x)]\in \pp(T_xX)$ and pick $g\in \Gamma_\nu$. Assume further that $g$ is loxodromic. We~will derive a contradiction by
showing that $\langle f, g\rangle$ must be elementary: this is a contradiction because
any non-elementary subgroup of $\Aut(X)$ contains a purely loxodromic non-elementary subgroup.
Let $\mathcal P$ be the measurable partition into ergodic components (under $f$) and
denote by $\mu_P$ the conditional
measure on $P\in \mathcal P$, so
that $\mu = \int \mu_{\mathcal P(x)} d\mu(x)$ is the ergodic decomposition of $\mu$.
Since the entropy function is affine, there exists a $f$-invariant
set $B$ of positive measure such that
for any $x\in B$, $h_{\mu_{\mathcal P(x)}}(f)>0$. In~particular $f$ is non-uniformly hyperbolic along $B$, so
along $B$, $E$ must coincide almost everywhere with one of $E^s_f$ or $E^u_f$. Reducing $B$ to a smaller
invariant subset we may
assume that $E = E^s_f$ almost everywhere along $B$.
For every $n\in \Z$, the automorphism $g^{-n} f g^n$ is loxodromic, preserves $\mu$, is
non-uniformly hyperbolic along $g^{-n}(B)$, and $E$ coincides with $E^s_{g^{-n}fg^n}$ almost everywhere. By~measure preservation
there exists $m\neq n$ such that $\mu(g^{-n}(B)\cap g^{-m}(B))>0$, so
$\mu(B\cap g^{m-n}(B))>0$. Letting $h = g^{m-n}fg^{-(m-n)}$ and
$A = B\cap g^{m-n}(B)$ we are exactly in the situation of Lemma 11.2 of~\cite{stiffness}, and
we conclude that $W^s(f, x) = W^s(h, x)$ for $\mu$-almost every $x\in A$, from which it follows that
$T^+_f= T^+_h$ and finally $(g^{m-n})^*T_f^+ = c T^+_f$. Since $g$ is loxodromic, this implies
that $T^+_f = T^+_g$ or $T^+_f = T^-_g$, and finally that $\langle f, g\rangle$ is elementary, which is
the sought-after contradiction.\enlargethispage{.5\baselineskip}%

Finally, if there is a measurable pair $\set{E_1, E_2}$ of line fields which is $\nu$-a.s.
invariant, we~get a $f$-invariant set $B$ of positive measure along which
$\aset{E_1(x), E_2(x) } = \aset{E^s_f(x) =\nobreak E^u_f(x)}$, and a set
$A = B\cap g^{m-n}(B)$ of positive measure along which
$\aset{E^s_f(x) = E^u_f(x)} = {E^s_h(x) = E^u_h(x)}$, where $h = g^{m-n}fg^{-(m-n)}$, and we conclude as before.
\end{proof}

% End literal retained input author-body.tex

\begin{thebibliography}{99}
\bibitem{BER} M. S. Baouendi, P. Ebenfelt \& L. P. Rothschild - Real submanifolds in complex space and their mappings , Princeton Math. Series , vol. 47 , Princeton University Press, Princeton, NJ, 1999
\bibitem{barrientos-malicet} P. G. Barrientos \& D. Malicet - “Extremal exponents of random products of conservative diffeomorphisms” , Math. Z. 296 (2020) no. 3-4, p. 1185-1207
\bibitem{benard-desaxce} T. Bénard \& N. de Saxcé - “Random walks with bounded first moment on finite-volume spaces” , Geom. Funct. Anal. 32 (2022) no. 4, p. 687-724
\bibitem{benoist-quint3} Y. Benoist \& J.-F. Quint - “Stationary measures and invariant subsets of homogeneous spaces (III)” , Ann. of Math. (2) 178 (2013) no. 3, p. 1017-1059
\bibitem{benoist-quint_book} Y. Benoist \& J.-F. Quint - Random walks on reductive groups , Ergeb. Math. Grenzgeb. (3) , vol. 62 , Springer, Cham, 2016
\bibitem{berger-turaev} P. Berger \& D. Turaev - “On Herman’s positive entropy conjecture” , Adv. Math. 349 (2019), p. 1234-1288
\bibitem{bierstone-milman:ihes} E. Bierstone \& P. D. Milman - “Semianalytic and subanalytic sets” , Publ. Math. Inst. Hautes Études Sci. (1988) no. 67, p. 5-42
\bibitem{blumenthal-xue-young} A. Blumenthal, J. Xue \& L.-S. Young - “Lyapunov exponents for random perturbations of some area-preserving maps including the standard map” , Ann. of Math. (2) 185 (2017) no. 1, p. 285-310
\bibitem{bogachev_book} V. I. Bogachev - Measure theory. Vol. I, II , Springer-Verlag, Berlin, 2007
\bibitem{bougerol-lacroix} P. Bougerol \& J. Lacroix - Products of random matrices with applications to Schrödinger operators , Progress in Probability and Statistics , vol. 8 , Birkhäuser Boston, Inc., Boston, MA, 1985
\bibitem{brandt} A. Brandt - “The stochastic equation $Y_{n+1}=A_nY_n+B_n$ with stationary coefficients” , Adv. in Appl. Probab. 18 (1986) no. 1, p. 211-220
\bibitem{brown-fisher-hurtado} A. Brown, D. Fisher \& S. Hurtado - “Zimmer’s conjecture: subexponential growth, measure rigidity, and strong property (T)” , Ann. of Math. (2) 196 (2022) no. 3, p. 891-940
\bibitem{br} A. Brown \& F. Rodriguez Hertz - “Measure rigidity for random dynamics on surfaces and related skew products” , J. Amer. Math. Soc. 30 (2017) no. 4, p. 1055-1132
\bibitem{cantat_groupes} S. Cantat - “Sur la dynamique du groupe d’automorphismes des surfaces K3” , Transform. Groups 6 (2001) no. 3, p. 201-214
\bibitem{Cantat:BHPS} S. Cantat - “Bers and Hénon, Painlevé and Schrödinger” , Duke Math. J. 149 (2009) no. 3, p. 411-460
\bibitem{Cantat:Panorama-Synthese} S. Cantat - “Quelques aspects des systèmes dynamiques polynomiaux: existence, exemples, rigidité” , in Quelques aspects des systèmes dynamiques polynomiaux , Panoramas \& Synthèses , vol. 30 , Société Mathématique de France, Paris, 2010, p. 13-95
\bibitem{Cantat:Milnor} S. Cantat - “Dynamics of automorphisms of compact complex surfaces” , in Frontiers in complex dynamics , Princeton Math. Series , vol. 51 , Princeton University Press, Princeton, NJ, 2014, p. 463-514
\bibitem{arxiv-version-of-this-paper} S. Cantat \& R. Dujardin - “Preprint version of this paper” , 2023
\bibitem{Simons} S. Cantat \& R. Dujardin - “Dynamics of automorphism groups of compact complex surfaces: examples, classification, and outlook” , 2023, to appear in Proceedings of the Simons Symposia
\bibitem{stiffness} S. Cantat \& R. Dujardin - “Random dynamics on real and complex projective surfaces” , J. reine angew. Math. (2023) no. 802, p. 1-76
\bibitem{finite_orbits} S. Cantat \& R. Dujardin - “Finite orbits for large groups of automorphisms of projective surfaces” , Compositio Math. 160 (2024) no. 1, p. 120-175
\bibitem{invariant} S. Cantat \& R. Dujardin - “Invariant measures for large automorphism groups of projective surfaces” , Transform. Groups 30 (2025) no. 1, p. 75-145
\bibitem{Cantat-Favre} S. Cantat \& C. Favre - “Symétries birationnelles des surfaces feuilletées” , J. reine angew. Math. 561 (2003), p. 199-235
\bibitem{chung} P. N. Chung - “Stationary measures and orbit closures of uniformly expanding random dynamical systems on surfaces” , 2020
\bibitem{Dolgachev-Keum} I. Dolgachev \& J. Keum - “Birational automorphisms of quartic Hessian surfaces” , Trans. Amer. Math. Soc. 354 (2002) no. 8, p. 3031-3057
\bibitem{dolgopyat-krikorian} D. Dolgopyat \& R. Krikorian - “On simultaneous linearization of diffeomorphisms of the sphere” , Duke Math. J. 136 (2007) no. 3, p. 475-505
\bibitem{douady-earle} A. Douady \& C. J. Earle - “Conformally natural extension of homeomorphisms of the circle” , Acta Math. 157 (1986) no. 1-2, p. 23-48
\bibitem{duistermaat} J. J. Duistermaat - Discrete integrable systems. QRT maps and elliptic surfaces , Springer Monographs in Math. , Springer, Berlin, 2010
\bibitem{degenerations_SL2} R. Dujardin \& C. Favre - “Degenerations of ${\rm SL}(2,\mathbb{C})$ representations and Lyapunov exponents” , Ann. H. Lebesgue 2 (2019), p. 515-565
\bibitem{durrett_book} R. Durrett - Probability: theory and examples , Duxbury Press, Belmont, CA, 1996
\bibitem{Eskin-Lindenstrauss:RandomWalk} A. Eskin \& E. Lindenstrauss - “Random walks on locally homogeneous spaces” , preprint
\bibitem{eskin-margulis} A. Eskin \& G. Margulis - “Recurrence properties of random walks on finite volume homogeneous manifolds” , in Random walks and geometry , Walter de Gruyter, Berlin, 2004, p. 431-444
\bibitem{eskin-mozes-oh} A. Eskin, S. Mozes \& H. Oh - “On uniform exponential growth for linear groups” , Invent. Math. 160 (2005) no. 1, p. 1-30
\bibitem{Feller:vol2} W. Feller - An introduction to probability theory and its applications. Vol. II , John Wiley \& Sons, Inc., New York-London-Sydney, 1971
\bibitem{fuchs-litman-silverman-tran} E. Fuchs, M. Litman, J. H. Silverman \& A. Tran - “Orbits on K3 surfaces of Markoff type” , Experiment. Math. 33 (2024) no. 4, p. 663-700
\bibitem{furstenberg-kifer} H. Furstenberg \& Y. Kifer - “Random matrix products and measures on projective spaces” , Israel J. Math. 46 (1983) no. 1-2, p. 12-32
\bibitem{griffiths-harris} P. Griffiths \& J. Harris - Principles of algebraic geometry , Wiley-Interscience, New York, 1978
\bibitem{hall-heyde} P. Hall \& C. C. Heyde - Martingale limit theory and its application , Probability and Math. Statistics , Academic Press, Inc., New York-London, 1980
\bibitem{hurtado:burnside} S. Hurtado, A. Kocsard \& F. Rodríguez-Hertz - “The Burnside problem for ${\rm Diff}_{\omega }(\mathbb{S}^2)$” , Duke Math. J. 169 (2020) no. 17, p. 3261-3290
\bibitem{katok-hasselblatt} A. Katok \& B. Hasselblatt - Introduction to the modern theory of dynamical systems , Encyclopedia of Math. and its Appl. , vol. 54 , Cambridge University Press, Cambridge, 1995
\bibitem{Keum-Kondo} J. Keum \& S. Kondō - “The automorphism groups of Kummer surfaces associated with the product of two elliptic curves” , Trans. Amer. Math. Soc. 353 (2001) no. 4, p. 1469-1487
\bibitem{Kifer} Y. Kifer - Ergodic theory of random transformations , Progress in Probability and Statistics , vol. 10 , Birkhäuser Boston, Inc., Boston, MA, 1986
\bibitem{ledrappier_stationary} F. Ledrappier - “Positivity of the exponent for stationary sequences of matrices” , in Lyapunov exponents (Bremen, 1984) , Lect. Notes in Math. , vol. 1186 , Springer, Berlin, 1986, p. 56-73
\bibitem{ledrappier-young3} F. Ledrappier \& L.-S. Young - “Entropy formula for random transformations” , Probab. Theory Relat. Fields 80 (1988) no. 2, p. 217-240
\bibitem{liu_thesis} X.-C. Liu - Lyapunov exponents approximation, symplectic cocycle deformation and a large deviation theorem , Ph. D. Thesis, IMPA, 2016
\bibitem{meyn-tweedie} S. Meyn \& R. L. Tweedie - Markov chains and stochastic stability , Cambridge University Press, Cambridge, 2009
\bibitem{obata-poletti} D. Obata \& M. Poletti - “Positive exponents for random products of conservative surface diffeomorphisms and some skew products” , J. Dynam. Differential Equations 34 (2022) no. 3, p. 2405-2428
\bibitem{Prohaska-Sert:TAMS} R. Prohaska \& C. Sert - “Markov random walks on homogeneous spaces and Diophantine approximation on fractals” , Trans. Amer. Math. Soc. 373 (2020) no. 11, p. 8163-8196
\bibitem{rebelo-roeder} J. Rebelo \& R. Roeder - “Dynamics of groups of automorphisms of character varieties and Fatou/Julia decomposition for Painlevé 6” , Indiana Univ. Math. J. 73 (2024) no. 6, p. 1967-2038
\bibitem{rosendal:book} C. Rosendal - Coarse geometry of topological groups , Cambridge Tracts in Math. , vol. 223 , Cambridge University Press, Cambridge, 2022
\bibitem{viana} M. Viana - Lectures on Lyapunov exponents , Cambridge Studies in Advanced Math. , vol. 145 , Cambridge University Press, Cambridge, 2014
\bibitem{zzhang} Z. Zhang - “On stable transitivity of finitely generated groups of volume-preserving diffeomorphisms” , Ergodic Theory Dynam. Systems 39 (2019) no. 2, p. 554-576
\bibitem{Zimmer:book} R. J. Zimmer - Ergodic theory and semisimple groups , Monographs in Math. , vol. 81 , Birkhäuser Verlag, Basel, 1984
\end{thebibliography}
\end{document}
